% Propriétés chimiques des alcools — Chimie Tle TI — Cours signé OnBuch+
% Programme officiel MINESEC, Terminale TI. Compiler avec : tectonic L01-proprietes-chimiques-des-alcools.tex
\newcommand{\DOCMATIERE}{Chimie}
\newcommand{\DOCNIVEAU}{Tle TI}
\newcommand{\DOCMODULE}{Module 1 · Chimie organique}
\newcommand{\DOCLECON}{01}
\newcommand{\DOCTITRE}{Propriétés chimiques des alcools}
% OnBuch+ — préambule « cours signés » ENRICHI (compilé avec tectonic / XeLaTeX)
% Système visuel « Pop » d'OnBuch+ : fond crème (jamais blanc), bordures encre
% épaisses, ombres « dures » décalées, titres Archivo Black en capitales.
% Version enrichie pour les cours de chimie : chimie (mhchem, chemfig),
% graphes (pgfplots), boîtes pédagogiques supplémentaires (définition,
% propriété, méthode, attention, le savais-tu, exercice résolu, exercice,
% TP, bilan), page de garde refaite et signature OnBuch+ ancrée en bas.
%
% Macros attendues dans main.tex AVANT \input{preamble} :
%   \DOCMATIERE  \DOCNIVEAU  \DOCMODULE  \DOCLECON (numéro)  \DOCTITRE
\documentclass[11pt]{article}

\usepackage{fontspec}
\usepackage[french]{babel}
\usepackage[a4paper,margin=2cm,top=3.4cm,bottom=2.8cm,headheight=1.4cm,footskip=1.3cm]{geometry}
\usepackage{amsmath,amssymb,amsfonts}
\usepackage[table]{xcolor} % [table] : fournit aussi \rowcolors (alternance de lignes)
\usepackage{pagecolor}
\usepackage{tikz}
\usetikzlibrary{arrows.meta,positioning,calc,shapes.geometric,shapes.misc,shapes.arrows,decorations.pathmorphing,decorations.markings,patterns,shadows,mindmap,trees,fit,backgrounds,matrix,intersections}
\usepackage{pgfplots}
\pgfplotsset{compat=1.18}
\usepackage[version=4]{mhchem}
\usepackage{chemfig}
\usepackage{enumitem}
\usepackage{fancyhdr}
% [most] charge toutes les bibliothèques tcolorbox (listings, minted, raster,
% documentation...) et gonfle inutilement la mémoire TeX sur un document long
% (~300 boîtes) : on ne charge que ce qui est réellement utilisé.
\usepackage[skins,breakable]{tcolorbox}
\usepackage{graphicx}
\usepackage{colortbl}
\usepackage{booktabs}
\usepackage{tabularx}
\usepackage{array}
\usepackage{multicol}
\usepackage{siunitx}
% parse-numbers=false : accepte aussi \SI{2,5 \times 10^{-3}}{...}
\sisetup{locale=FR,output-decimal-marker={,},group-minimum-digits=4,parse-numbers=false}
\DeclareSIUnit\ohm{\ensuremath{\symup{\Omega}}} % U+2126 absent de la police
\usepackage{qrcode}
% \degree : commande fréquemment utilisée par les modèles pour un angle ou une
% température en dehors de \SI{}{} (non fournie par les paquets chargés).
\providecommand{\degree}{^{\circ}}
\usepackage{lastpage}
\usepackage{hyperref}
\hypersetup{hidelinks}

% ============================================================
% Palette « Pop » OnBuch+ — mêmes hex que la classe P (plus_ui.dart)
% ============================================================
\definecolor{poporange}{HTML}{F59321}
\definecolor{popdark}{HTML}{E07A0C}
\definecolor{popcream}{HTML}{FFF6E8}
\definecolor{popink}{HTML}{1C1714}
\definecolor{poppurple}{HTML}{7A5AE0}
\definecolor{popgreen}{HTML}{2E9E6A}
\definecolor{poppink}{HTML}{E0457B}
\definecolor{popblue}{HTML}{2F7DE1}
\definecolor{popgold}{HTML}{C98A12}
\definecolor{popmuted}{HTML}{8A7F76}
\definecolor{popline}{HTML}{EADFD2}
\definecolor{popgray}{HTML}{8A7F76}
\colorlet{popgrey}{popgray}
% Alias tolérants (le modèle nomme parfois les couleurs différemment)
\colorlet{popwhite}{white}
\colorlet{popblack}{popink}
\colorlet{popred}{poppink}
\colorlet{popyellow}{popgold}
% Teintes claires pour fonds de tableaux / zones
\colorlet{popblueL}{popblue!10!white}
\colorlet{poporangeL}{poporange!14!white}
\colorlet{popgreenL}{popgreen!12!white}
\colorlet{poppurpleL}{poppurple!12!white}
\colorlet{poppinkL}{poppink!12!white}
\colorlet{popgoldL}{popgold!14!white}

\pagecolor{popcream}

% ============================================================
% Typographie
% ============================================================
\defaultfontfeatures{Ligatures=TeX}
\setmainfont{PlusJakartaSans-Medium.ttf}[Path=../../../fonts/,BoldFont=PlusJakartaSans-Bold.ttf]
% Maths en sans-serif (Fira Math) avec chiffres et lettres droites de Plus
% Jakarta, pour que les formules se fondent dans le texte.
\usepackage[math-style=ISO,bold-style=ISO]{unicode-math}
\setmathfont{FiraMath-Regular.otf}[Path=../../../fonts/]
\setmathfont{PlusJakartaSans-Medium.ttf}[Path=../../../fonts/,range={up/num,up/{Latin,latin},\mathup}]
\usepackage{icomma} % virgule décimale sans espace en mode math
% Caractères absents de la police de texte (sinon carré vide) : rendus par
% leur équivalent mathématique, en texte comme en formule.
\usepackage{newunicodechar}
\newunicodechar{α}{\ensuremath{\alpha}}
\newunicodechar{Α}{\ensuremath{\alpha}}
\newunicodechar{β}{\ensuremath{\beta}}
\newunicodechar{δ}{\ensuremath{\delta}}
\newunicodechar{✓}{\popcheck{popgreen}}
\newfontfamily\popdisplay{ArchivoBlack-Regular.ttf}[Path=../../../fonts/]
\newfontfamily\poplogo{SpaceGrotesk-Bold.ttf}[Path=../../../fonts/]
\newfontfamily\popsemi{PlusJakartaSans-SemiBold.ttf}[Path=../../../fonts/]
\newfontfamily\popxbold{PlusJakartaSans-ExtraBold.ttf}[Path=../../../fonts/]
\newfontfamily\popxboldls{PlusJakartaSans-ExtraBold.ttf}[Path=../../../fonts/,LetterSpace=3.0]

\setlist[enumerate]{leftmargin=1.7em,itemsep=5pt,topsep=4pt}
\setlist[itemize]{leftmargin=1.7em,itemsep=4pt,topsep=4pt,label={\color{poporange}\textbullet}}
\setlength{\parindent}{0pt}
\setlength{\parskip}{5pt}
\renewcommand{\arraystretch}{1.25}

% chemfig : liaisons un peu plus épaisses, encre Pop
\setchemfig{atom sep=2em,bond style={line width=0.9pt,popink},cram width=3pt}

% pgfplots : style Pop par défaut
\pgfplotsset{
  every axis/.append style={axis line style={popink,line width=1pt},tick style={popink},
    grid style={popline,line width=0.5pt},label style={font=\small},tick label style={font=\footnotesize}},
  popcurve/.style={poporange,line width=1.8pt,no markers,smooth},
}

% ============================================================
% Pill / squiggle
% ============================================================
\newcommand{\poppill}[3]{%
  \tikz[baseline=-0.55ex]\node[draw=popink,line width=1.2pt,rounded corners=2.6pt,%
    fill=#2,inner xsep=6.5pt,inner ysep=3pt]{%
    \popxboldls\fontsize{7}{7}\selectfont\color{#3}\MakeUppercase{#1}};%
}
\newcommand{\popsquiggle}[1]{%
  \tikz[baseline=-0.4ex]{
    \draw[#1,line width=2.6pt,line cap=round]
      plot [smooth] coordinates {(0,0.09) (0.34,0.01) (0.68,0.21) (1.02,0.01) (1.36,0.10)};
  }%
}
% Mot-clé surligné (marqueur orange sous le texte)
\newcommand{\cle}[1]{{\popxbold\color{popdark}#1}}

% ============================================================
% En-tête / pied
% ============================================================
\pagestyle{fancy}
\fancyhf{}
\renewcommand{\headrulewidth}{0pt}
\renewcommand{\footrulewidth}{0pt}
\lhead{\raisebox{-0.15ex}{\poplogo\fontsize{13}{13}\selectfont\color{popink}OnBuch\color{poporange}+}}
\rhead{\poppill{\DOCMATIERE\ \textbullet\ \DOCNIVEAU\ \textbullet\ Leçon \DOCLECON}{white}{popink}}
\lfoot{\popxbold\fontsize{7.6}{7.6}\selectfont\color{popmuted}\MakeUppercase{Cours signé OnBuch+}\ \textbullet\ {\popsemi\DOCTITRE}}
\rfoot{\tikz[baseline=-0.6ex]\node[rounded corners=8.5pt,fill=popink,minimum height=17pt,minimum width=17pt,inner xsep=5pt,inner sep=0]{\popxbold\fontsize{8}{8}\selectfont\color{popcream}\thepage\,/\,\pageref*{LastPage}};}

% ============================================================
% Titres
% ============================================================
\newcommand{\chaptitre}[2]{%
  \begin{tcolorbox}[enhanced,colback=poporange,colframe=popink,boxrule=2.6pt,arc=11pt,
      drop shadow=popink,left=15pt,right=15pt,top=12pt,bottom=11pt,before skip=3pt,after skip=18pt]
    \poppill{#2}{popink}{popcream}\par\vspace{9pt}
    {\raggedright\hyphenpenalty=10000\exhyphenpenalty=10000\popdisplay\fontsize{22}{24}\selectfont\color{popcream}\MakeUppercase{#1}\par}
  \end{tcolorbox}
}
% Section numérotée : pastille orange avec le numéro + titre Archivo
\newcounter{popsec}
\newcommand{\coursec}[1]{%
  \stepcounter{popsec}\setcounter{popsub}{0}%
  \par\vspace{16pt}\noindent
  \begin{minipage}{\linewidth}
  \tikz[baseline=-0.7ex]\node[draw=popink,line width=1.6pt,fill=poporange,rounded corners=4pt,minimum size=22pt,inner sep=0]{\popdisplay\fontsize{12}{12}\selectfont\color{popcream}\thepopsec};%
  \hspace{8pt}{\popdisplay\fontsize{15}{17}\selectfont\color{popink}\MakeUppercase{#1}}\par
  \vspace{4pt}\noindent{\color{poporange}\rule{36pt}{3.6pt}}
  \end{minipage}\par\nopagebreak\vspace{8pt}
}
\newcounter{popsub}
\newcommand{\courssub}[1]{%
  \stepcounter{popsub}%
  \par\vspace{9pt}\noindent{\popxbold\fontsize{11.5}{13}\selectfont\color{popink}{\color{poporange}\thepopsec.\thepopsub}\hspace{6pt}#1}\par\nopagebreak\vspace{4pt}
}

% ============================================================
% Boîtes pédagogiques — même recette : fond blanc, bordure épaisse colorée,
% ombre dure encre, badge plein. Titre optionnel [#1].
% ============================================================
\tcbset{popbase/.style={breakable,enhanced,colback=white,boxrule=2.3pt,arc=9pt,
  shadow={3pt}{-3pt}{0pt}{popink},left=14pt,right=14pt,top=11pt,bottom=11pt,before skip=11pt,after skip=15pt,
  fontupper=\popsemi\fontsize{10.6}{14.5}\selectfont\color{popink},
  fontlower=\popsemi\fontsize{10.6}{14.5}\selectfont\color{popink}}}
% Coche dessinée (le glyphe ✓ n'existe pas dans les polices du document)
\newcommand{\popcheck}[1]{\tikz[baseline=-0.5ex]\draw[#1,line width=1.6pt,line cap=round,line join=round](-0.13,0.0)--(-0.04,-0.09)--(0.13,0.1);}
\newcommand{\popboxhead}[3]{\poppill{#1}{#2}{white}\ifx\relax#3\relax\else\hspace{6pt}{\popxbold\fontsize{10.6}{12}\selectfont\color{popink}#3}\fi\par\vspace{7pt}}

\newtcolorbox{aretenir}[1][]{popbase,colframe=popgreen,before upper={\popboxhead{À retenir}{popgreen}{#1}}}
\newtcolorbox{exemplebox}[1][]{popbase,colframe=poporange,before upper={\popboxhead{Exemple}{poporange}{#1}}}
\newtcolorbox{definition}[1][]{popbase,colframe=popblue,before upper={\popboxhead{Définition}{popblue}{#1}}}
\newtcolorbox{propriete}[1][]{popbase,colframe=poppurple,before upper={\popboxhead{Propriété}{poppurple}{#1}}}
\newtcolorbox{methode}[1][]{popbase,colframe=popgold,before upper={\popboxhead{Méthode}{popgold}{#1}}}
\newtcolorbox{attention}[1][]{popbase,colframe=poppink,before upper={\popboxhead{Attention}{poppink}{#1}}}
\newtcolorbox{experience}[1][]{popbase,colframe=popblue,colback=popblueL,before upper={\popboxhead{Expérience}{popblue}{#1}}}
% « Le savais-tu ? » : carte violette pleine (contexte, culture, Cameroun)
\newtcolorbox{savaistu}[1][]{popbase,colback=poppurple,colframe=popink,
  fontupper=\popsemi\fontsize{10.4}{14.2}\selectfont\color{white},
  before upper={\poppill{Le savais-tu\,?}{white}{poppurple}\ifx\relax#1\relax\else\hspace{6pt}{\popxbold\color{white}#1}\fi\par\vspace{7pt}}}
% Exercice résolu : énoncé en haut, \tcblower puis la solution sur fond crème
\newcounter{popexo}\newcounter{popexr}
\newtcolorbox{exoresolu}[1][]{popbase,colframe=poporange,colbacklower=popcream,
  segmentation style={popink,line width=1.2pt,dashed},
  before upper={\stepcounter{popexr}\popboxhead{Exercice résolu \thepopexr}{poporange}{#1}},
  before lower={\poppill{Solution}{popink}{popcream}\par\vspace{6pt}}}
% Exercice (non résolu en place) — la correction est dans la partie Corrigés
\newtcolorbox{exercice}[1][]{popbase,colframe=popink,
  before upper={\stepcounter{popexo}\popboxhead{Exercice \thepopexo}{popink}{#1}}}
% Corrigé
\newtcolorbox{corrige}[1][]{popbase,colframe=popgreen,colback=popgreenL,
  before upper={\popboxhead{Corrigé}{popgreen}{#1}}}
% Objectifs / prérequis / bilan
\newtcolorbox{objectifs}{popbase,colframe=popink,colback=poporangeL,
  before upper={\popboxhead{Objectifs de la leçon}{popink}{}}}
\newtcolorbox{prerequis}{popbase,colframe=popink,colback=popblueL,
  before upper={\popboxhead{Prérequis}{popblue}{}}}
\newtcolorbox{bilan}{popbase,colframe=popink,colback=white,boxrule=2.8pt,
  before upper={\popboxhead{Fiche bilan}{poporange}{L'essentiel à connaître pour l'examen}}}
% Figure encadrée avec légende
\newtcolorbox{popfigure}[1][]{enhanced,breakable=false,colback=white,colframe=popline,boxrule=1.4pt,arc=8pt,
  left=10pt,right=10pt,top=10pt,bottom=8pt,before skip=10pt,after skip=12pt,halign=center,
  lower separated=false,halign lower=center,
  fontlower=\popsemi\fontsize{9}{11.5}\selectfont\color{popmuted},
  #1}
\newcommand{\legende}[1]{\par\vspace{6pt}{\popsemi\fontsize{9}{11.5}\selectfont\color{popmuted}#1}}

% ============================================================
% Signature OnBuch+
% ============================================================
\newcommand{\poplogoinline}[1]{{\poplogo\fontsize{#1}{#1}\selectfont\color{popink}OnBuch\color{poporange}+}}
\newcommand{\signatureonbuch}{%
  \begin{tcolorbox}[enhanced,colback=popink,colframe=popink,boxrule=2.2pt,arc=10pt,
      drop shadow=poporange,left=14pt,right=14pt,top=10pt,bottom=10pt,before skip=0pt,after skip=0pt]
    \tikz[baseline=-0.7ex]\node[circle,fill=popgreen,draw=popcream,line width=1.3pt,minimum size=22pt,inner sep=0]{\popcheck{white}};%
    \hspace{9pt}\begin{minipage}[c]{0.62\linewidth}
      {\popxbold\fontsize{12}{13}\selectfont\color{popcream}Signé\ \textbullet\ {\poplogo OnBuch\color{poporange}+}}\par\vspace{2pt}
      {\raggedright\popsemi\fontsize{8}{10.5}\selectfont\color{popline}\MakeUppercase{Équipe pédagogique OnBuch+}\\\MakeUppercase{Conforme au programme officiel MINESEC}\par}
    \end{minipage}\hfill
    {\poplogo\fontsize{22}{22}\selectfont\color{popcream}OnBuch\color{poporange}+}
  \end{tcolorbox}%
}

% ============================================================
% Page de garde
% #1 titre · #2 matière · #3 niveau · #4 module · #5 n° leçon · #6 durée ·
% #7 description / objectifs (texte ou itemize)
% ============================================================
\newcommand{\pagegarde}[7]{%
  \thispagestyle{empty}
  \newgeometry{a4paper,margin=1.7cm,top=1.6cm,bottom=1.4cm}%
  \begin{tikzpicture}[remember picture,overlay]
    \fill[poporange!18] ([xshift=-3.2cm,yshift=-3.6cm]current page.north east) circle (4.6cm);
    \fill[poppurple!14] ([xshift=2.4cm,yshift=7.6cm]current page.south west) circle (3.8cm);
  \end{tikzpicture}%
  {\poplogo\fontsize{30}{30}\selectfont\color{popink}OnBuch\color{poporange}+}\hfill
  \raisebox{0.6ex}{\poppill{Cours signé \textbullet\ Premium}{popink}{popcream}}\par
  \vspace{1pt}\popsquiggle{poppurple}\par
  \vspace{8pt}
  \begin{tcolorbox}[enhanced,colback=poporange,colframe=popink,boxrule=2.8pt,arc=13pt,
      drop shadow=popink,left=18pt,right=18pt,top=12pt,bottom=13pt,after skip=10pt]
    \poppill{#2 \textbullet\ #3}{popink}{popcream}\hspace{5pt}\poppill{#4}{white}{popink}\par\vspace{8pt}
    {\popdisplay\fontsize{14}{14}\selectfont\color{popink}LEÇON #5}\par\vspace{4pt}
    {\raggedright\hyphenpenalty=10000\exhyphenpenalty=10000\popdisplay\fontsize{25}{27}\selectfont\color{popcream}\MakeUppercase{#1}\par}
    \vspace{8pt}
    {\popxbold\fontsize{9.5}{9.5}\selectfont\color{popink}\MakeUppercase{Durée conseillée : #6}}
  \end{tcolorbox}
  \begin{tcolorbox}[enhanced,colback=white,colframe=popink,boxrule=2.2pt,arc=10pt,
      drop shadow=popink,left=16pt,right=16pt,top=10pt,bottom=10pt,after skip=8pt]
    \poppill{Dans cette leçon}{poppurple}{white}\par\vspace{5pt}
    \popsemi\fontsize{9.3}{12.6}\selectfont\color{popink}#7
  \end{tcolorbox}
  \noindent\begin{minipage}[t]{0.487\linewidth}
    \begin{tcolorbox}[enhanced,colback=popgreen,colframe=popink,boxrule=2.2pt,arc=10pt,
        drop shadow=popink,left=11pt,right=11pt,top=10pt,bottom=10pt,height=3.1cm,valign=center]
      \tcbox[colback=white,colframe=popink,boxrule=1.3pt,arc=5pt,boxsep=0pt,nobeforeafter,
        left=3pt,right=3pt,top=3pt,bottom=3pt,valign=center]{\qrcode[height=1.9cm]{https://play.google.com/store/apps/details?id=com.luvvix.onbuch&hl=fr}}%
      \hspace{8pt}\begin{minipage}[c]{0.5\linewidth}
        {\popdisplay\fontsize{11}{12.5}\selectfont\color{white}\MakeUppercase{Télécharge OnBuch}}\par\vspace{4pt}
        {\raggedright\popsemi\fontsize{8.4}{11}\selectfont\color{white}Gratuit \textbullet\ Google Play\\Tuteur IA, annales, concours, résultats\par}
      \end{minipage}
    \end{tcolorbox}
  \end{minipage}\hfill
  \begin{minipage}[t]{0.487\linewidth}
    \begin{tcolorbox}[enhanced,colback=poppurple,colframe=popink,boxrule=2.2pt,arc=10pt,
        drop shadow=popink,left=11pt,right=11pt,top=10pt,bottom=10pt,height=3.1cm,valign=center]
      \tikz[baseline=-0.7ex]\node[circle,fill=white,draw=popink,line width=1.3pt,minimum size=1.6cm,inner sep=0]{\popdisplay\fontsize{24}{24}\selectfont\color{poppurple}+};%
      \hspace{8pt}\begin{minipage}[c]{0.6\linewidth}
        {\popdisplay\fontsize{11}{12.5}\selectfont\color{white}\MakeUppercase{Passe à OnBuch+}}\par\vspace{4pt}
        {\raggedright\popsemi\fontsize{8.4}{11}\selectfont\color{white}Tous les cours signés, Tuteur IA illimité, fiches PDF — onglet OnBuch+ de l'app\par}
      \end{minipage}
    \end{tcolorbox}
  \end{minipage}
  \vspace{8pt}
  \popcontenu
  \vfill
  \signatureonbuch
  \restoregeometry
  \newpage
}

% Rangée « Ce cours contient » (page de garde)
\newcommand{\poptile}[3]{% #1 couleur #2 gros texte #3 légende
  \begin{tcolorbox}[enhanced,colback=#1,colframe=popink,boxrule=2pt,arc=9pt,shadow={2.5pt}{-2.5pt}{0pt}{popink},
      left=6pt,right=6pt,top=9pt,bottom=9pt,halign=center,valign=center,height=1.75cm,width=\linewidth,nobeforeafter]
    {\popdisplay\fontsize{12.5}{13}\selectfont\color{white}\MakeUppercase{#2}}\par\vspace{3pt}
    {\centering\popsemi\fontsize{7.8}{9.5}\selectfont\color{white}#3\par}
  \end{tcolorbox}}
\newcommand{\popcontenu}{%
  \par\noindent{\popxboldls\fontsize{7.5}{7.5}\selectfont\color{popmuted}CE COURS SIGNÉ CONTIENT}\par\vspace{7pt}
  \noindent\begin{minipage}[t]{0.235\linewidth}\poptile{popblue}{Cours}{complet, illustré, pas à pas}\end{minipage}\hfill
  \begin{minipage}[t]{0.235\linewidth}\poptile{poporange}{Méthodes}{et exercices résolus}\end{minipage}\hfill
  \begin{minipage}[t]{0.235\linewidth}\poptile{poppink}{Exercices}{type examen + corrigés}\end{minipage}\hfill
  \begin{minipage}[t]{0.235\linewidth}\poptile{popgreen}{Fiche bilan}{l'essentiel à retenir}\end{minipage}\par}

% Sommaire
\newtcolorbox{sommaire}{enhanced,colback=white,colframe=popink,boxrule=2.2pt,arc=10pt,
  drop shadow=popink,left=18pt,right=18pt,top=13pt,bottom=13pt,before skip=6pt,after skip=20pt,
  before upper={\poppill{Sommaire}{poppurple}{white}\par\vspace{9pt}\popsemi\fontsize{10.6}{14.5}\selectfont\color{popink}}}

% Signature de fin de cours — poussée tout en bas de la dernière page
\newcommand{\findecours}{%
  \par\vfill\nopagebreak
  \begin{center}\popsquiggle{poporange}\end{center}\vspace{4pt}
  \signatureonbuch
}
% Glyphes absents de Fira Math : redéfinis à partir de glyphes disponibles
\AtBeginDocument{%
  \def\setminus{\mathbin{\backslash}}\let\smallsetminus\setminus
  \def\vdots{\mathord{\vbox{\baselineskip3.2pt\lineskiplimit0pt\kern4pt\hbox{.}\hbox{.}\hbox{.}}}}%
  \def\ddots{\mathinner{\mkern1mu\raise6.5pt\hbox{.}\mkern2mu\raise3.6pt\hbox{.}\mkern2mu\raise0.7pt\hbox{.}\mkern1mu}}%
  \def\triangle{\mathord{\scalebox{0.9}{$\symup{\Delta}$}}}%
  \def\nabla{\mathord{\raisebox{\depth}{\scalebox{1}[-1]{$\symup{\Delta}$}}}}%
  \def\checkmark{\popcheck{popdark}}%
  \def\star{\mathbin{\ast}}\def\bigstar{\mathord{\ast}}%
  \def\diamond{\mathbin{\scalebox{0.6}{$\Diamond$}}}%
  \def\bigcup{\mathop{\mathchoice{\raisebox{-0.3em}{\scalebox{1.7}{$\cup$}}}{\scalebox{1.25}{$\cup$}}{\cup}{\cup}}\displaylimits}%
  \def\bigcap{\mathop{\mathchoice{\raisebox{-0.3em}{\scalebox{1.7}{$\cap$}}}{\scalebox{1.25}{$\cap$}}{\cap}{\cap}}\displaylimits}%
  \def\lesseqqgtr{\mathrel{\lessgtr}}\def\bigoplus{\mathop{\oplus}\displaylimits}%
}
\providecommand{\ssi}{si et seulement si} % abréviation fréquente

\begin{document}

\pagegarde{Propriétés chimiques des alcools}{Chimie}{Tle TI}{Module 1 · Chimie organique}{01}{4 h (cours + TP)}{Cette leçon explore les grandes réactions des alcools : réaction avec le sodium, déshydratations, oxydations (combustion et oxydation ménagée selon la classe de l'alcool) et estérification. Elle présente aussi les deux voies de préparation des alcools — hydratation des alcènes et fermentation alcoolique — en lien avec la production locale d'éthanol.
\vspace{4pt}
\begin{multicols}{2}\small
\begin{itemize}
\item À la fin de cette leçon, je sais écrire l'équation-bilan de la réaction d'un alcool avec le sodium et interpréter le dégagement de dihydrogène.
\item Je sais mettre en évidence la mobilité de l'hydrogène du groupe hydroxyle des alcools.
\item Je sais distinguer et équilibrer les déshydratations intramoléculaire (alcène) et intermoléculaire (éther-oxyde) selon les conditions opératoires.
\item Je sais réaliser l'expérience de la lampe sans flamme et identifier les produits de l'oxydation catalytique de l'éthanol.
\item Je sais prévoir et écrire les produits de l'oxydation ménagée de chaque classe d'alcool (primaire, secondaire, tertiaire) par le permanganate de potassium.
\item Je sais réaliser l'oxydation ménagée d'un aldéhyde et identifier un aldéhyde par les tests à la liqueur de Fehling et au réactif de Tollens.
\end{itemize}\end{multicols}}

\begin{sommaire}
\begin{multicols}{2}
\begin{enumerate}
\item Rappels et réaction des alcools avec le sodium
\item Déshydratation des alcools
\item Oxydation des alcools : combustion et lampe sans flamme
\item Oxydation ménagée selon la classe de l'alcool
\item Estérification
\item Préparation des alcools
\item Travaux pratiques
\item Méthodes et exercices résolus
\item Exercices
\item Corrigés des exercices
\item Fiche bilan
\end{enumerate}
\end{multicols}
\end{sommaire}

\begin{objectifs}
\begin{itemize}
\item À la fin de cette leçon, je sais écrire l'équation-bilan de la réaction d'un alcool avec le sodium et interpréter le dégagement de dihydrogène.
\item Je sais mettre en évidence la mobilité de l'hydrogène du groupe hydroxyle des alcools.
\item Je sais distinguer et équilibrer les déshydratations intramoléculaire (alcène) et intermoléculaire (éther-oxyde) selon les conditions opératoires.
\item Je sais réaliser l'expérience de la lampe sans flamme et identifier les produits de l'oxydation catalytique de l'éthanol.
\item Je sais prévoir et écrire les produits de l'oxydation ménagée de chaque classe d'alcool (primaire, secondaire, tertiaire) par le permanganate de potassium.
\item Je sais réaliser l'oxydation ménagée d'un aldéhyde et identifier un aldéhyde par les tests à la liqueur de Fehling et au réactif de Tollens.
\item Je sais écrire l'équation-bilan d'une estérification et citer ses caractéristiques (lente, limitée, athermique).
\item Je sais écrire les équations d'hydratation d'un alcène et de fermentation alcoolique du glucose.
\end{itemize}
\end{objectifs}

\begin{prerequis}
\begin{itemize}
\item Nomenclature et formules des alcools ; classes des alcools (primaire, secondaire, tertiaire) vues en début d'année
\item Couples oxydant/réducteur et demi-équations électroniques (Première)
\item Réactions d'oxydoréduction en solution : permanganate de potassium en milieu acide
\item Formules générales des alcènes, acides carboxyliques et esters
\item Notion de catalyseur et d'équilibre chimique (limites d'une transformation)
\end{itemize}
\end{prerequis}

\newpage

\coursec{Rappels et réaction des alcools avec le sodium}

\courssub{Situation d'accroche et problématique}

Imagine un instant que tu te trouves à Bafoussam, au cœur de l'Ouest camerounais. Une coopérative locale transforme le jus de canne à sucre en éthanol par fermentation. Ce précieux liquide alimente les lampes à alcool qui éclairent les cases lors des veillées, sert de désinfectant dans les centres de santé de proximité, mais il est aussi la matière première pour fabriquer l'arôme de banane (acétate d'isoamyle) qui parfume les bonbons « \emph{Koki} » industriels, ou encore pour produire de l'éther, solvant indispensable aux laboratoires de la SONARA. Comment une seule molécule, l'éthanol, peut-elle être le point de départ de transformations si variées : oxydation en vinaigre, estérification en arôme, déshydratation en éther ou en éthylène ?

\medskip
\noindent
\textbf{Problématique :} Quelles sont les propriétés chimiques fondamentales des alcools qui expliquent leur réactivité multiple ? Comment la structure du groupe hydroxyle dicte-t-elle leur comportement face aux métaux, aux acides, aux agents oxydants ou déshydratants ?

\medskip
\noindent
Dans cette première section, nous posons les bases : la structure des alcools et leur réaction caractéristique avec le sodium, qui révèle l'acidité faible mais bien réelle de l'hydrogène porté par l'oxygène.

\courssub{Rappels sur la structure et la classification des alcools}

Un \cle{alcool} est un composé organique caractérisé par la présence d'un \cle{groupe fonctionnel hydroxyle} ($-\text{OH}$) lié à un atome de carbone tétraédrique (hybridation $sp^3$). Sa formule générale s'écrit \textbf{$\text{R}-\text{OH}$}, où $\text{R}$ représente un groupe alkyle (chaîne carbonée saturée) ou aryle.

\begin{definition}[Classes des alcools]
La classe d'un alcool dépend du nombre de groupes alkylés portés par le carbone \textbf{porteur du groupe $\text{OH}$} (carbone $\alpha$) :
\begin{itemize}
    \item \textbf{Alcool primaire ($1^\circ$) :} le carbone $\alpha$ porte \textbf{un seul} groupe alkyle (ex. : éthanol $\ce{CH3-CH2-OH}$, butan-1-ol).
    \item \textbf{Alcool secondaire ($2^\circ$) :} le carbone $\alpha$ porte \textbf{deux} groupes alkylés (ex. : propan-2-ol $\ce{CH3-CH(OH)-CH3}$).
    \item \textbf{Alcool tertiaire ($3^\circ$) :} le carbone $\alpha$ porte \textbf{trois} groupes alkylés (ex. : 2-méthylpropan-2-ol $\ce{(CH3)3C-OH}$).
\end{itemize}
\end{definition}

\begin{popfigure}
\begin{center}
\chemfig{
    % Ethanol (Primaire)
    CH_3-CH_2-[:30]O-[:90]H
}
\qquad
\chemfig{
    % Propan-2-ol (Secondaire)
    CH_3-[:30]C(-[:90]OH)(-[:-30]CH_3)-[:-30]H
}
\qquad
\chemfig{
    % 2-Méthylpropan-2-ol (Tertiaire)
    CH_3-[:30]C(-[:90]OH)(-[:-30]CH_3)(-[:150]CH_3)
}
\end{center}
\legende{Représentation de Cram (chemfig) d'un alcool primaire (éthanol), secondaire (propan-2-ol) et tertiaire (2-méthylpropan-2-ol). Le carbone $\alpha$ est le centre des liaisons.}
\end{popfigure}

\courssub{Polarité de la liaison O--H : comparaison avec l'eau}

Le groupe hydroxyle possède une liaison \ce{O-H} polarisée car l'oxygène est plus électronégatif ($\chi_{\ce{O}} = 3,44$) que l'hydrogène ($\chi_{\ce{H}} = 2,20$). Le doublet non liant sur l'oxygène renforce ce dipôle.

\begin{center}
\begin{tabularx}{\linewidth}{|l|X|X|}
\hline
\rowcolor{popblueL}
\textbf{Propriété} & \textbf{Eau ($\ce{H2O}$)} & \textbf{Alcool ($\ce{R-OH}$)} \\
\hline
Structure & $\ce{H-O-H}$ & $\ce{R-O-H}$ \\
\hline
Polarité $\ce{O-H}$ & Forte & Forte (légèrement moindre si $\text{R}$ est donneur $+I$) \\
\hline
Liaison hydrogène & Donneur \textbf{et} Accepteur (2 donneurs) & Donneur (1 H) \textbf{et} Accepteur (2 doublets sur O) \\
\hline
Acidité ($pK_a$ approx.) & $15,7$ & $16 \sim 18$ (selon classe) \\
\hline
\end{tabularx}
\end{center}

\begin{aretenir}[Acidité relative]
Les alcools sont des acides \textbf{plus faibles que l'eau} ($pK_a$ plus grand). L'effet inductif donneur ($+I$) des groupes alkylés $\text{R}$ augmente la densité électronique sur l'oxygène, déstabilisant l'anion alcoolate $\ce{RO-}$ formé par déprotonation. L'ordre d'acidité est : \textbf{Primaire $>$ Secondaire $>$ Tertiaire} (le tertiaire est le moins acide, le plus basique).
\end{aretenir}

\courssub{Expérience : Action du sodium métallique sur l'éthanol}

C'est une réaction classique de mise en évidence du caractère acide (faible) de l'hydrogène du groupe hydroxyle.

\begin{experience}[Action du sodium sur l'éthanol]
\textbf{Matériel :} Tube à essai, support, bec Bunsen, tuyau de recueil des gaz (par déplacement d'air), allumette, pincette.
\textbf{Produits :} Éthanol \textbf{anhydre} (sans eau !), sodium métallique (morceau de la taille d'un grain de riz, conservé sous huile minérale).
\textbf{Protocole :}
\begin{enumerate}
    \item Verser \SI{2}{mL} d'éthanol anhydre dans le tube à essai.
    \item À l'aide de la pincette, prélever un \textbf{tout petit} morceau de sodium, l'essuyer soigneusement sur papier filtre (retirer l'huile), et le plonger dans l'éthanol.
    \item Boucher immédiatement le tube avec le tuyau de recueil.
    \item Observer : effervescence lente, dégagement d'un gaz incolore. Le tube se tiédit (réaction exothermique).
    \item Recueillir le gaz dans un tube à essai retourné (par déplacement d'air, car $\ce{H2}$ est peu soluble dans l'eau).
    \item \textbf{Test à la flamme :} Approcher une allumette allumée de l'orifice du tube recueillant le gaz.
\end{enumerate}
\textbf{Observations :}
\begin{itemize}
    \item Le sodium \textbf{coule au fond} du tube (densité $\rho_{\ce{Na}} = \SI{0,97}{g.cm^{-3}} > \rho_{\ce{EtOH}} \approx \SI{0,79}{g.cm^{-3}}$), mais l'effervescence (bulles de $\ce{H2}$) le fait remonter et « danser » dans le liquide.
    \item Formation d'un gaz incolore.
    \item \textbf{Test à la flamme :} Le gaz brûle avec une flamme pâle bleue accompagnée d'un \textbf{bruit sec caractéristique (détonation)} : c'est le \textbf{dihydrogène $\ce{H2}$}.
    \item Le mélange réactionnel, une fois évaporé, laisse un solide blanc : l'\textbf{éthanolate de sodium} ($\ce{C2H5ONa}$).
\end{itemize}
\textbf{Interprétation :} Le sodium, métal très réducteur, arrache l'hydrogène de la liaison $\ce{O-H}$ polaire. C'est une réaction d'oxydo-réduction : $\ce{Na -> Na+ + e-}$ (oxydation) et $\ce{2 H+ + 2 e- -> H2}$ (réduction). L'hydrogène de l'alcool se comporte comme $\ce{H+}$.
\end{experience}

\begin{popfigure}
\begin{tikzpicture}[scale=0.9, every node/.style={transform shape}]
    % Support
    \draw[thick] (0,0) -- (0,-0.5) -- (4,-0.5) -- (4,0);
    \draw[thick] (2,-0.5) -- (2,-1.5); % Pied
    % Tube à essai
    \draw[thick, popblue] (1.8,0) -- (1.8,4) -- (2.2,4) -- (2.2,0);
    \draw[thick, popblue] (1.8,4) .. controls (1.8,4.3) and (2.2,4.3) .. (2.2,4); % Fond arrondi
    % Liquide
    \fill[poporange!30] (1.85,0.2) rectangle (2.15,2.5);
    \node[popdark, align=center] at (2,1.3) {Éthanol anhydre};
    % Sodium (au fond)
    \fill[popgray] (1.95,0.4) circle (0.1);
    \node[popdark, right] at (2.2,0.4) {Na (s)};
    % Bulles remontant
    \foreach \y in {0.8, 1.1, 1.4, 1.7, 2.0, 2.3, 2.6, 2.9, 3.2, 3.5} {
        \draw[popblue, thick] (2.0, \y) circle (0.06);
        \draw[popblue, thick] (1.92, \y+0.04) circle (0.04);
    }
    % Tuyau
    \draw[thick, popmuted] (2.0, 4.1) -- (2.0, 5.0) -- (4.0, 5.0) -- (4.0, 3.0);
    % Tube recueil gaz (inversé)
    \draw[thick, popblue] (3.8, 3.0) -- (3.8, 1.0) -- (4.2, 1.0) -- (4.2, 3.0);
    \draw[thick, popblue] (3.8, 1.0) .. controls (3.8, 0.7) and (4.2, 0.7) .. (4.2, 1.0);
    % Eau dans la cuve (schématique)
    \fill[popblue!20] (3.5, 1.0) rectangle (4.5, 0.5);
    % Flamme
    \draw[poporange, thick, decorate, decoration={snake, amplitude=1.5pt, segment length=4pt}] (4.0, 3.0) -- (4.0, 3.8);
    \fill[poporange!80] (3.9, 3.8) .. controls (3.7, 4.2) and (4.3, 4.2) .. (4.1, 3.8);
    \node[popdark, right] at (4.2, 3.5) {Test $\ce{H2}$ : \textbf{« Pouf ! »}};
    % Légendes
    \node[popdark, align=left] at (-0.5, 2) {Dispositif :\\Recueil du gaz\\par déplacement d'air};
\end{tikzpicture}
\legende{Schéma du dispositif : action du sodium sur l'éthanol, recueil du dihydrogène et test de détonation à la flamme.}
\end{popfigure}

\courssub{Équation-bilan et formation de l'alcoolate}

L'équation-bilan de la réaction de l'éthanol avec le sodium s'écrit :

\[
\ce{2 C2H5OH (l) + 2 Na (s) -> 2 C2H5ONa (s) + H2 (g)}
\]

\begin{itemize}
    \item \textbf{Produits :} Le gaz est le \textbf{dihydrogène $\ce{H2}$} (test de la flamme : détonation). Le solide blanc est l'\textbf{éthanolate de sodium} ($\ce{C2H5ONa}$), sel ionique composé du cation $\ce{Na+}$ et de l'anion \textbf{éthanolate} $\ce{C2H5O-}$.
    \item \textbf{État physico-chimique :} La réaction se fait dans l'éthanol \textbf{anhydre}. La présence d'eau ferait réagir le sodium violemment avec l'eau ($\ce{2 Na + 2 H2O -> 2 NaOH + H2}$), masquant la réaction recherchée et étant dangereuse (projection de soude caustique, inflammation de l'hydrogène).
\end{itemize}

\begin{popfigure}
\begin{tikzpicture}[baseline=(current bounding box.center), node distance=1.2cm]
    \node (eth) {\chemfig{CH_3-CH_2-O-H}};
    \node[right=of eth] (plus1) {\textbf{+}};
    \node[right=of plus1] (na) {\chemfig{Na}};
    \node[right=of na] (arrow) {\ce{->[Na][\Delta]}};
    \node[right=of arrow, xshift=0.5cm] (etna) {\chemfig{CH_3-CH_2-O^{-}Na^{+}}};
    \node[right=of etna] (plus2) {\textbf{+}};
    \node[right=of plus2] (h2) {\chemfig{H-H}};
\end{tikzpicture}
\legende{Transformation de l'éthanol en éthanolate de sodium. La liaison $\ce{O-H}$ est rompue, l'hydrogène part sous forme de $\ce{H2}$, l'oxygène garde les deux électrons de la liaison et porte une charge négative, compensée par $\ce{Na+}$.}
\end{popfigure}

\courssub{Généralisation et interprétation : mobilité de l'hydrogène}

Pour tout alcool $\ce{R-OH}$, la réaction avec le sodium s'écrit :

\[
\ce{2 R-OH + 2 Na -> 2 R-O^{-}Na^{+} + H2}
\]

\begin{definition}[Alcoolate]
Un \textbf{alcoolate} (ou alkoxide) est un sel ionique de formule \textbf{$\ce{R-O^{-}M^{+}}$} (souvent $\ce{M = Na, K}$), obtenu par réaction d'un alcool sur un métal alcalin (ou une base forte comme $\ce{NaH}$). L'anion \textbf{alcoolate $\ce{R-O^{-}}$} est la \textbf{base conjuguée} de l'alcool $\ce{R-OH}$.
\end{definition}

\begin{propriete}[Caractère basique de l'alcoolate]
L'ion alcoolate $\ce{RO-}$ est une \textbf{base forte} (bien plus forte que $\ce{OH-}$ dans les milieux aprotiques) et un \textbf{nucléophile} puissant. Il est immédiatement protoné par l'eau :
\[
\ce{RO- + H2O -> ROH + HO-}
\]
C'est pourquoi les alcoolates doivent être manipulés sous atmosphère inerte (azote, argon) et dans des solvants anhydres (éthanol, THF, toluène).
\end{propriete}

\begin{methode}[Équilibrer la réaction Alcool + Sodium]
\begin{enumerate}
    \item Écrire les réactifs : $\ce{ROH + Na}$.
    \item Identifier les produits : \textbf{Alcoolate $\ce{RONa}$} + \textbf{Dihydrogène $\ce{H2}$} (Jamais $\ce{H2O}$ !).
    \item Mettre le coefficient \textbf{2} devant l'alcool et devant le sodium (pour fournir 2 électrons et former $\ce{H2}$).
    \item Mettre le coefficient \textbf{2} devant l'alcoolate.
    \item Vérifier : $\ce{2 ROH + 2 Na -> 2 RONa + H2}$.
\end{enumerate}
\end{methode}

\begin{attention}[Piège classique : $\ce{H2O}$ au lieu de $\ce{H2}$]
\textbf{Erreur fréquente :} Écrire $\ce{ROH + Na -> RONa + H2O}$ ou $\ce{2 ROH + 2 Na -> 2 RONa + H2O}$.
\textbf{Pourquoi c'est faux :} Le sodium est un \textbf{réducteur}. Il donne des électrons aux protons $\ce{H+}$ (issus de la polarisation $\ce{O-H}$) pour former $\ce{H2}$. Il n'y a \textbf{pas d'atome d'oxygène} apporté par le sodium pour former de l'eau. L'oxygène reste dans l'alcoolate.
\textbf{Astuce :} Le test à la flamme (détonation) prouve la présence de $\ce{H2}$, pas de $\ce{H2O}$.
\end{attention}

\courssub{Influence de la classe de l'alcool sur la vitesse de réaction}

La réaction est une réaction acide-base (transfert de proton) dont la vitesse dépend de l'acidité de l'alcool, elle-même liée à la stabilité de l'alcoolate formé.

\begin{propriete}[Ordre de réactivité vis-à-vis du sodium]
\textbf{Alcool primaire $>$ Alcool secondaire $>$ Alcool tertiaire}
\begin{itemize}
    \item \textbf{Primaire (ex. éthanol) :} Réaction \textbf{rapide}, effervescence vive, dégagement de chaleur sensible.
    \item \textbf{Secondaire (ex. propan-2-ol) :} Réaction \textbf{plus lente}, effervescence modérée.
    \item \textbf{Tertiaire (ex. 2-méthylpropan-2-ol) :} Réaction \textbf{très lente}, voire imperceptible à froid. Nécessite souvent un chauffage.
\end{itemize}
\end{propriete}

\begin{aretenir}[Explication par l'effet inductif]
Les groupes alkylés $\text{R}$ ont un effet inductif \textbf{donneur ($+I$)}. Ils « poussent » la densité électronique vers l'oxygène.
\begin{itemize}
    \item Dans un alcool tertiaire, 3 groupes $+I$ enrichissent l'oxygène en densité électronique $\rightarrow$ la liaison $\ce{O-H}$ est \textbf{moins polarisée} $\rightarrow$ l'hydrogène est \textbf{moins $\delta+$} $\rightarrow$ moins attiré par les électrons du sodium.
    \item L'alcoolate tertiaire $\ce{R3C-O-}$ est \textbf{destabilisé} par l'excès de densité électronique sur l'oxygène (effet $+I$ des 3 groupes alkylés) $\rightarrow$ l'équilibre de déprotonation est déplacé vers les réactifs.
\end{itemize}
\emph{Paradoxe apparent :} L'alcool tertiaire est le plus basique (son alcoolate est la base la plus forte), mais c'est celui qui réagit le \textbf{moins vite} avec le sodium car il est le \textbf{moins acide}.
\end{aretenir}

\courssub{Exemple résolu : Dosage de l'alcool par réaction avec le sodium}

\begin{exoresolu}[Détermination de la masse molaire d'un alcool par dégagement d'hydrogène]
\textbf{Énoncé :} On introduit $\SI{0,460}{g}$ d'un alcool primaire saturé $\ce{C_nH_{2n+1}OH}$ dans un tube à essai contenant un excès de sodium métallique. On recueille le dihydrogène dégagé par déplacement d'eau. Au bout de quelques minutes, la réaction est terminée. Le volume de gaz recueilli est $\SI{154}{mL}$ dans les conditions normales de température et de pression (CNTP : $T = \SI{273,15}{K}$, $P = \SI{1,013e5}{Pa}$). On suppose le gaz parfait.
Déterminer la formule brute et le nom de l'alcool.
\textbf{Données :} $V_m = \SI{22,4}{L.mol^{-1}}$ aux CNTP. $M(\ce{C}) = \SI{12,0}{g.mol^{-1}}$, $M(\ce{H}) = \SI{1,0}{g.mol^{-1}}$, $M(\ce{O}) = \SI{16,0}{g.mol^{-1}}$.

\tcblower
\textbf{Solution détaillée :}

\noindent \textbf{1. Quantité de matière de dihydrogène formé}
\[
n_{\ce{H2}} = \frac{V_{\ce{H2}}}{V_m} = \frac{\SI{0,154}{L}}{\SI{22,4}{L.mol^{-1}}} = \SI{6,875e-3}{mol}
\]

\noindent \textbf{2. Relation stœchiométrique (Tableau d'avancement)}
Équation : $\ce{2 ROH + 2 Na -> 2 RONa + H2}$
\begin{center}
\begin{tabular}{|c|c|c|c|c|}
\hline
 & $\ce{2 ROH}$ & $\ce{2 Na}$ & $\ce{2 RONa}$ & $\ce{H2}$ \\
\hline
État initial & $n_0$ & excès & 0 & 0 \\
\hline
État final & 0 & excès & $n_0$ & $n_0/2$ \\
\hline
\end{tabular}
\end{center}
D'après le tableau : $n_{\ce{H2}} = \frac{n_{\text{alcool initial}}}{2}$.
Donc :
\[
n_{\text{alcool}} = 2 \times n_{\ce{H2}} = 2 \times \SI{6,875e-3}{mol} = \SI{1,375e-2}{mol}
\]

\noindent \textbf{3. Masse molaire de l'alcool}
\[
M_{\text{alcool}} = \frac{m_{\text{alcool}}}{n_{\text{alcool}}} = \frac{\SI{0,460}{g}}{\SI{1,375e-2}{mol}} \approx \SI{33,45}{g.mol^{-1}}
\]

\noindent \textbf{4. Détermination de $n$}
Formule brute : $\ce{C_nH_{2n+2}O}$ (alcool primaire saturé $\ce{C_nH_{2n+1}OH}$).
\[
M = n \times 12,0 + (2n+2) \times 1,0 + 16,0 = 14n + 18
\]
\[
14n + 18 = 33,45 \implies 14n = 15,45 \implies n \approx 1,1
\]
$n$ étant un entier, $n = 1$.

\noindent \textbf{5. Conclusion}
L'alcool est le \textbf{méthanol} ($\ce{CH3OH}$, $M = \SI{32,0}{g.mol^{-1}}$). La légère différence sur la masse molaire expérimentale ($\SI{33,45}{g.mol^{-1}}$) s'explique par des imprécisions de mesure de volume (température, pression, lecture ménisque) ou une légère impureté (eau) dans l'échantillon.

\textbf{Réponse :} L'alcool est le \textbf{méthanol} (alcool méthylique).
\end{exoresolu}

\courssub{Applications et contexte : l'éthanolate de sodium au Cameroun}

\begin{savaistu}[L'éthanolate de sodium, base de la chimie fine]
Au laboratoire de chimie organique de l'Université de Yaoundé I ou dans les unités de formulation de l'industrie pharmaceutique (type \emph{Cinpharm} ou laboratoires de galénique), l'éthanolate de sodium ($\ce{EtONa}$) et le \emph{tert}-butylate de potassium ($\ce{tBuOK}$) sont des réactifs standards.
\begin{itemize}
    \item \textbf{Synthèse d'esters (transestérification) :} $EtONa$ catalyse la réaction entre une huile (triglycérides) et du méthanol pour produire du biodiesel (méthyl-esters d'acides gras) et du glycérol. C'est la voie principale pour valoriser l'huile de palme locale en carburant.
    \item \textbf{Condensation de Claisen :} Formation de liaisons $C{-}C$ entre esters.
    \item \textbf{Réaction de Williamson :} $R{-}ONa + R'{-}X \rightarrow R{-}O{-}R' + NaX$ (préparation d'éthers, voir section Déshydratation intermoléculaire).
\end{itemize}
Attention : $\ce{EtONa}$ s'hydrolyse instantanément à l'air humide ($\ce{EtONa + H2O -> EtOH + NaOH}$). Il se stocke sous azote sec.
\end{savaistu}

\begin{exercice}[Entraînement : Rédaction et prévision]
\begin{enumerate}
    \item Écrire l'équation-bilan de la réaction du butan-2-ol (alcool secondaire) avec le sodium. Nommer le sel formé.
    \item Prédire et justifier l'ordre de vitesse de réaction du sodium avec : éthanol, propan-2-ol, 2-méthylpropan-2-ol.
    \item On réalise l'expérience de la « lampe sans flamme » (vue en section 3) avec un fil de cuivre au lieu du sodium. Quelle différence fondamentale existe-t-il entre la réaction $\ce{Cu/EtOH}$ (chauffé) et $\ce{Na/EtOH}$ (à froid) quant à la nature de la transformation chimique (oxydo-réduction vs acide-base) ?
\end{enumerate}
\end{exercice}

\begin{corrige}[Exercice 1]
\begin{enumerate}
    \item \textbf{Équation :} $\ce{2 CH3CH(OH)CH2CH3 + 2 Na -> 2 CH3CH(ONa)CH2CH3 + H2}$.
    \textbf{Nom du sel :} \textbf{Butan-2-olate de sodium} (ou sec-butylate de sodium).
    \item \textbf{Ordre :} Éthanol (primaire) $>$ Propan-2-ol (secondaire) $>$ 2-Méthylpropan-2-ol (tertiaire).
    \textbf{Justification :} L'effet inductif donneur ($+I$) des groupes alkylés augmente la densité électronique sur l'O, diminue la polarité $\ce{O-H}$, diminue l'acidité de l'alcool et déstabilise l'alcoolate formé. L'acidité décroît : $1^\circ > 2^\circ > 3^\circ$, donc la vitesse de réaction avec $\ce{Na}$ suit le même ordre.
    \item $\ce{Na/EtOH}$ (à froid) : \textbf{Réaction acide-base} (transfert $\ce{H+}$). $\ce{Na}$ s'oxyde en $\ce{Na+}$, l'alcool se réduit en $\ce{H2}$. Le carbone de l'alcool ne change pas de nombre d'oxydation.
    $\ce{Cu/EtOH}$ (chauffé, lampe sans flamme) : \textbf{Oxydation catalytique (déshydrogénation oxydante)}. L'éthanol est \textbf{oxydé} en éthanal ($\ce{CH3CHO}$) : $\ce{CH3CH2OH ->[Cu, \Delta] CH3CHO + H2}$. Le cuivre (ou $\ce{CuO}$ réduit en $\ce{Cu}$) joue le rôle de catalyseur/déhydrogénateur. C'est une transformation du carbone (nombre d'oxydation $-1 \to +1$), pas un simple transfert de proton.
\end{enumerate}
\end{corrige}

\coursec{Déshydratation des alcools}

Dans la section précédente, nous avons vu que l'hydrogène du groupe hydroxyle $-\mathrm{OH}$ est \cle{mobile} : le sodium vient le déloger et libère du dihydrogène. Mais le groupe $-\mathrm{OH}$ recèle une seconde fragilité : sous l'action d'un acide concentré et de la chaleur, il peut quitter la molécule \emph{en entraînant avec lui un hydrogène voisin}, pour former une molécule d'eau. C'est la \cle{déshydratation}, une réaction capitale en chimie industrielle — c'est elle qui permet, par exemple, de transformer l'éthanol de fermentation (vin de palme, bili-bili) en éthylène, matière première des plastiques.

\courssub{Notion générale de déshydratation}

\begin{definition}[Déshydratation d'un alcool]
La \cle{déshydratation} d'un alcool est une réaction d'\emph{élimination} au cours de laquelle une molécule d'eau $\ce{H2O}$ est arrachée à partir d'un alcool. Elle se réalise :
\begin{itemize}
\item en présence d'acide sulfurique concentré $\ce{H2SO4}$ ou d'acide phosphorique $\ce{H3PO4}$ (déshydratants), à chaud ;
\item ou par passage des vapeurs d'alcool sur de l'alumine $\ce{Al2O3}$ chauffée (catalyseur hétérogène).
\end{itemize}
Selon les conditions opératoires (température, proportions), on obtient soit un \cle{alcène} (déshydratation \emph{intramoléculaire}), soit un \cle{éther-oxyde} (déshydratation \emph{intermoléculaire}).
\end{definition}

L'intuition est simple : l'acide sulfurique concentré est un « avaleur d'eau ». Chauffé avec un alcool, il arrache $-\mathrm{OH}$ d'un carbone et $-\mathrm{H}$ d'un carbone voisin (voie intra), ou bien il fait se rejoindre deux molécules d'alcool qui perdent ensemble une molécule d'eau (voie inter). \textbf{C'est la température qui décide du destin de l'alcool.}

\courssub{Déshydratation intramoléculaire : obtention d'un alcène}

\begin{definition}[Déshydratation intramoléculaire]
La déshydratation \cle{intramoléculaire} est l'élimination d'une molécule d'eau \emph{au sein d'une même molécule d'alcool} : le groupe $-\mathrm{OH}$ et un atome d'hydrogène porté par le carbone \emph{voisin} (carbone en $\beta$) sont éliminés. Il se forme un \cle{alcène} (apparition d'une double liaison $\ce{C=C}$).
\end{definition}

\textbf{Conditions opératoires :}
\begin{itemize}
\item $\ce{H2SO4}$ concentré, chauffage vers \SI{170}{\celsius} ;
\item ou passage des vapeurs d'alcool sur alumine $\ce{Al2O3}$ vers \SI{400}{\celsius}.
\end{itemize}

\textbf{Exemple fondamental : l'éthanol donne l'éthylène (éthène).}
\[ \ce{CH3-CH2-OH ->[H2SO4\ conc.][170\ \text{\degres C}] CH2=CH2 + H2O} \]
Commentaire : une seule molécule d'éthanol intervient ; on arrache $-\mathrm{OH}$ du carbone 1 et $-\mathrm{H}$ du carbone 2 ; la double liaison s'installe entre les deux carbones. L'équation est équilibrée : $\ce{C2H6O -> C2H4 + H2O}$.

\textbf{Généralisation} (alcool primaire non ramifié) :
\[ \mathrm{R{-}CH_2{-}CH_2{-}OH} \;\longrightarrow\; \mathrm{R{-}CH{=}CH_2} + \ce{H2O} \]

\begin{exemplebox}[Déshydratation du butan-2-ol et règle de Zaïtsev]
Le butan-2-ol \chemfig{CH_3-CH(-[2]OH)-CH_2-CH_3} possède deux carbones $\beta$ différents : deux alcènes peuvent donc se former :
\begin{itemize}
\item le but-1-ène \chemfig{CH_2=CH-CH_2-CH_3} (hydrogène arraché au carbone 1) ;
\item le but-2-ène \chemfig{CH_3-CH=CH-CH_3} (hydrogène arraché au carbone 3).
\end{itemize}
La \cle{règle de Zaïtsev} précise que \textbf{l'alcène majoritaire est le plus substitué}, c'est-à-dire celui dont la double liaison porte le plus grand nombre de groupes alkyles. Ici, le but-2-ène (double liaison disubstituée) est majoritaire devant le but-1-ène (monosubstitué).
\end{exemplebox}

\begin{propriete}[Règle de Zaïtsev (complément Bac)]
Lors de la déshydratation intramoléculaire d'un alcool pouvant conduire à plusieurs alcènes isomères, l'alcène formé \emph{majoritairement} est celui dont la double liaison $\ce{C=C}$ est la plus substituée (le plus entourée de groupes alkyles), car c'est le plus stable.
\end{propriete}

\courssub{Déshydratation intermoléculaire : obtention d'un éther-oxyde}

\begin{definition}[Déshydratation intermoléculaire]
La déshydratation \cle{intermoléculaire} est l'élimination d'une molécule d'eau \emph{entre deux molécules d'alcool} : le groupe $-\mathrm{OH}$ de l'une et l'hydrogène du groupe $-\mathrm{OH}$ de l'autre sont éliminés. Les deux restes se soudent par un atome d'oxygène : il se forme un \cle{éther-oxyde} de formule générale $\mathrm{R{-}O{-}R'}$.
\end{definition}

\textbf{Conditions opératoires :}
\begin{itemize}
\item $\ce{H2SO4}$ concentré, \emph{excès d'alcool}, température modérée vers \SI{140}{\celsius} ;
\item ou vapeurs d'alcool sur alumine $\ce{Al2O3}$ vers \SI{250}{\celsius}.
\end{itemize}

\textbf{Exemple fondamental : l'éthanol donne l'éther éthylique (éthoxyéthane).}
\[ \ce{2 CH3-CH2-OH ->[H2SO4\ conc.][140\ \text{\degres C}] CH3-CH2-O-CH2-CH3 + H2O} \]
Commentaire : \textbf{deux} molécules d'éthanol participent ; une seule molécule d'eau est éliminée pour les deux ; le produit $\ce{C2H5-O-C2H5}$ est l'éther éthylique. Vérification du bilan : $\ce{2 C2H6O -> C4H10O + H2O}$ — les atomes se conservent.

\textbf{Généralisation} (alcools identiques) :
\[ 2\,\mathrm{R{-}OH} \;\longrightarrow\; \mathrm{R{-}O{-}R} + \ce{H2O} \]

\begin{popfigure}
\begin{center}
\begin{tikzpicture}[>=Stealth, node distance=0.4cm]
\node (eth) {\chemfig{CH_3-CH_2-OH}};
\node[right=2.6cm of eth, yshift=1.1cm] (alcene) {\chemfig{CH_2=CH_2}};
\node[right=2.6cm of eth, yshift=-1.1cm] (ether) {\chemfig{CH_3-CH_2-O-CH_2-CH_3}};
\draw[->, very thick, popblue] (eth.east) -- ++(0.7,0) |- (alcene.west)
  node[pos=0.25, above, font=\small\bfseries, text=popblue] {\SI{170}{\celsius}, $\ce{H2SO4}$ conc.}
  node[pos=0.25, below, font=\small, text=popblue] {intramoléculaire : $+\,\ce{H2O}$};
\draw[->, very thick, poporange] (eth.east) -- ++(0.7,0) |- (ether.west)
  node[pos=0.25, above, font=\small\bfseries, text=poporange] {\SI{140}{\celsius}, excès d'alcool}
  node[pos=0.25, below, font=\small, text=poporange] {intermoléculaire : 2 molécules, $+\,\ce{H2O}$};
\node[above=0.15cm of alcene, font=\small\itshape, text=popdark] {éthylène (alcène)};
\node[below=0.15cm of ether, font=\small\itshape, text=popdark] {éther éthylique};
\end{tikzpicture}
\end{center}
\legende{Les deux voies de déshydratation de l'éthanol : la température oriente la réaction vers l'alcène (voie haute) ou vers l'éther-oxyde (voie basse).}
\end{popfigure}

\courssub{Le rôle décisif de la température}

La même matière première (l'éthanol) et le même déshydratant ($\ce{H2SO4}$ concentré) peuvent donner deux produits totalement différents. Tout se joue sur le thermostat :

\begin{center}
\begin{tabularx}{\linewidth}{|l|X|X|X|}
\hline
\rowcolor{popblueL} \textbf{Conditions} & \textbf{Type de déshydratation} & \textbf{Produit obtenu} & \textbf{Équation} \\
\hline
$\ce{H2SO4}$ conc., $\approx$ \SI{140}{\celsius}, excès d'alcool & Intermoléculaire & Éther-oxyde $\mathrm{R{-}O{-}R}$ & $\ce{2 C2H5OH -> C2H5-O-C2H5 + H2O}$ \\
\hline
$\ce{H2SO4}$ conc., $\approx$ \SI{170}{\celsius} & Intramoléculaire & Alcène & $\ce{C2H5OH -> CH2=CH2 + H2O}$ \\
\hline
$\ce{Al2O3}$, $\approx$ \SI{250}{\celsius} & Intermoléculaire & Éther-oxyde & $\ce{2 R-OH -> R-O-R + H2O}$ \\
\hline
$\ce{Al2O3}$, $\approx$ \SI{400}{\celsius} & Intramoléculaire & Alcène & $\mathrm{R{-}CH_2{-}CH_2{-}OH} \to \mathrm{R{-}CH{=}CH_2} + \ce{H2O}$ \\
\hline
\end{tabularx}
\end{center}

\begin{popfigure}
\begin{center}
\begin{tikzpicture}[>=Stealth, scale=1.0]
\draw[->, very thick, popdark] (0,0) -- (10.5,0) node[below left, font=\small] {température};
\foreach \x/\t in {2/100, 4/140, 6/170, 8.5/250}{
  \draw[thick, popdark] (\x,0.12) -- (\x,-0.12) node[below, font=\small] {\SI{\t}{\celsius}};
}
% Zone éther : trait épais orange sous l'axe
\draw[very thick, poporange] (3.1,-0.8) -- (4.9,-0.8);
\node[below, font=\small\bfseries, text=poporange] at (4,-1.1) {Zone éther-oxyde};
% Zone alcène : trait épais bleu sous l'axe
\draw[very thick, popblue] (5.1,-0.8) -- (6.9,-0.8);
\node[below, font=\small\bfseries, text=popblue] at (6,-1.1) {Zone alcène};
% Flèches vers les produits
\node[above, font=\small, text=poporange, align=center] at (4,0.5) {$\ce{C2H5-O-C2H5}$\\ (intermoléculaire)};
\node[above, font=\small, text=popblue, align=center] at (6,0.5) {$\ce{CH2=CH2}$\\ (intramoléculaire)};
\draw[->, thick, poporange] (4,0.15) -- (4,0.35);
\draw[->, thick, popblue] (6,0.15) -- (6,0.35);
% Note Al2O3
\node[font=\small\itshape, text=popmuted, align=center, anchor=west] at (8.7,0.4) {Sur $\ce{Al2O3}$ :\\ éther $\approx$ \SI{250}{\celsius}\\ alcène $\approx$ \SI{400}{\celsius}};
\end{tikzpicture}
\end{center}
\legende{Axe des températures (éthanol $+$ $\ce{H2SO4}$ concentré) : vers \SI{140}{\celsius} l'éther éthylique domine ; vers \SI{170}{\celsius} et au-delà, c'est l'éthylène qui se forme.}
\end{popfigure}

\begin{methode}[Prévoir le produit d'une déshydratation]
\begin{enumerate}
\item \textbf{Regarder la température} : $\approx$ \SI{140}{\celsius} (ou \SI{250}{\celsius} sur $\ce{Al2O3}$) $\Rightarrow$ éther-oxyde ; $\approx$ \SI{170}{\celsius} et plus (ou \SI{400}{\celsius} sur $\ce{Al2O3}$) $\Rightarrow$ alcène.
\item \textbf{Compter les molécules} : éther $\Rightarrow$ deux molécules d'alcool dans l'équation ; alcène $\Rightarrow$ une seule.
\item \textbf{Si plusieurs alcènes sont possibles}, appliquer la règle de Zaïtsev : l'alcène majoritaire est le plus substitué.
\item \textbf{Vérifier l'équilibrage} : une molécule d'eau éliminée dans les deux cas.
\end{enumerate}
\end{methode}

\begin{attention}[Erreurs fréquentes]
\begin{itemize}
\item \textbf{Confondre les deux déshydratations} : avant d'écrire l'équation, demande-toi toujours si \emph{une} ou \emph{deux} molécules d'alcool interviennent. Une molécule $\to$ alcène ; deux molécules $\to$ éther-oxyde.
\item \textbf{Oublier les conditions sur la flèche} : au Baccalauréat, une flèche sans température ni catalyseur coûte des points. Écris toujours $\ce{H2SO4}$ concentré et la température.
\item \textbf{Croire que l'eau vient « de l'acide »} : non, $\ce{H2SO4}$ est un déshydratant ; la molécule d'eau provient entièrement de l'alcool.
\item \textbf{Mal équilibrer la voie intermoléculaire} : il faut un coefficient 2 devant l'alcool.
\end{itemize}
\end{attention}

\courssub{Influence de la classe de l'alcool}

Tous les alcools ne se déshydratent pas avec la même facilité. L'expérience montre que la déshydratation intramoléculaire est d'autant plus facile que l'alcool est \emph{plus substitué} :

\begin{propriete}[Facilité de déshydratation]
La facilité de déshydratation intramoléculaire croît dans l'ordre :
\[ \text{alcool tertiaire} \;>\; \text{alcool secondaire} \;>\; \text{alcool primaire} \]
Un alcool tertiaire (comme le 2-méthylpropan-2-ol) se déshydrate dans des conditions douces, parfois dès la température ambiante en présence d'acide ; un alcool primaire exige des températures élevées (\SI{170}{\celsius} et plus).
\end{propriete}

\begin{exoresolu}[Fabrication industrielle de l'éther éthylique]
\textbf{Énoncé.} Une unité pilote fait réagir \SI{92}{g} d'éthanol en présence d'acide sulfurique concentré maintenu à \SI{140}{\celsius}. On admet que la déshydratation intermoléculaire est totale. Déterminer la quantité et la masse maximales d'éther éthylique $\ce{C2H5-O-C2H5}$ obtenues. On donne : $M(\ce{C}) = \SI{12}{g.mol^{-1}}$, $M(\ce{H}) = \SI{1}{g.mol^{-1}}$, $M(\ce{O}) = \SI{16}{g.mol^{-1}}$.
\tcblower
\textbf{Solution détaillée.}

\emph{Étape 1 : masses molaires.}
\begin{align*}
M(\ce{C2H5OH}) &= 2\times 12 + 6\times 1 + 16 = \SI{46}{g.mol^{-1}}\\
M(\ce{C4H10O}) &= 4\times 12 + 10\times 1 + 16 = \SI{74}{g.mol^{-1}}
\end{align*}

\emph{Étape 2 : quantité d'éthanol.}
\[ n(\ce{C2H5OH}) = \frac{m}{M} = \frac{92}{46} = \SI{2,0}{mol} \]

\emph{Étape 3 : équation et tableau d'avancement.}
\[ \ce{2 C2H5OH ->[H2SO4,\ 140\ \text{\degres C}] C2H5-O-C2H5 + H2O} \]
\begin{center}
\begin{tabular}{lccc}
\hline
 & $\ce{2 C2H5OH}$ & $\ce{C2H5-O-C2H5}$ & $\ce{H2O}$ \\
\hline
État initial (mol) & $2,0$ & $0$ & $0$ \\
État final (mol) & $2,0 - 2x_{max}$ & $x_{max}$ & $x_{max}$ \\
\hline
\end{tabular}
\end{center}
L'éthanol est entièrement consommé : $2,0 - 2x_{max} = 0 \Rightarrow x_{max} = \SI{1,0}{mol}$.

\emph{Étape 4 : masse d'éther.}
\[ m(\text{éther}) = n \times M = 1,0 \times 74 = \SI{74}{g} \]
\textbf{Conclusion :} à partir de \SI{92}{g} d'éthanol (2 mol), on obtient au maximum \SI{74}{g} d'éther éthylique (1 mol) et \SI{18}{g} d'eau. Remarque utile : la somme $74 + 18 = \SI{92}{g}$ vérifie la conservation de la masse.
\end{exoresolu}

\begin{savaistu}[L'éther, premier anesthésique de l'histoire]
L'éther éthylique $\ce{C2H5-O-C2H5}$, produit de la déshydratation intermoléculaire de l'éthanol, a été utilisé pour la première fois comme anesthésique général en chirurgie le 16 octobre 1846, à Boston, par le dentiste William Morton. Ce jour-là, un patient fut opéré d'une tumeur du cou sans ressentir aucune douleur : la chirurgie moderne était née. Aujourd'hui, l'éther a été remplacé par des anesthésiques plus sûrs (il est très inflammable !), mais il reste un solvant de référence au laboratoire. Quant à l'éthylène, produit de la déshydratation intramoléculaire, il est la matière première du polyéthylène, le plastique des sachets et bidons que l'on retrouve partout au Cameroun.
\end{savaistu}

\begin{aretenir}[L'essentiel de la déshydratation]
\begin{itemize}
\item La déshydratation élimine $\ce{H2O}$ d'un alcool, avec $\ce{H2SO4}$ concentré à chaud ou sur $\ce{Al2O3}$.
\item \textbf{Intramoléculaire} (1 molécule, $\approx$ \SI{170}{\celsius} ou \SI{400}{\celsius}/$\ce{Al2O3}$) $\Rightarrow$ \textbf{alcène} : $\ce{C2H5OH -> CH2=CH2 + H2O}$.
\item \textbf{Intermoléculaire} (2 molécules, $\approx$ \SI{140}{\celsius} ou \SI{250}{\celsius}/$\ce{Al2O3}$) $\Rightarrow$ \textbf{éther-oxyde} : $\ce{2 C2H5OH -> C2H5-O-C2H5 + H2O}$.
\item Règle de Zaïtsev : l'alcène majoritaire est le plus substitué.
\item Facilité de déshydratation : tertiaire $>$ secondaire $>$ primaire.
\end{itemize}
\end{aretenir}

\coursec{Oxydation des alcools : combustion et lampe sans flamme}

Dans la section précédente, nous avons vu que l'alcool peut \emph{perdre de l'eau} (déshydratation) pour donner un alcène ou un éther-oxyde : la molécule se transforme, mais son squelette carboné reste globalement intact. Nous allons maintenant étudier ce qui se passe lorsque l'alcool \emph{gagne de l'oxygène} ou \emph{perd de l'hydrogène} : c'est le domaine de l'\cle{oxydation}. Tu connais déjà l'oxydation la plus spectaculaire : quand on allume un peu d'alcool à \SI{90}{\celsius} (ou de kérosène), la flamme bleue qui apparaît traduit une \cle{combustion}. Mais il existe une oxydation beaucoup plus douce, la \cle{lampe sans flamme}, qui transforme l'éthanol en éthanal sans aucune flamme visible. C'est elle qui ouvre la porte de toute la chimie des composés carbonylés.

\courssub{La combustion complète des alcools}

\textbf{Intuition.}\par
Verse quelques gouttes d'alcool à \SI{90}{\celsius} (ou d'éthanol de pharmacie) sur une soucoupe et approche une allumette : une flamme bleue, presque invisible en plein jour, se propage instantanément. C'est la même réaction qui fait tourner les moteurs des véhicules brésiliens roulant à l'éthanol, ou qui chauffe les réchauds à alcool utilisés dans certains ménages. L'alcool réagit avec le \cle{dioxygène} de l'air : c'est une combustion.

\begin{definition}[Combustion complète d'un alcool]
La \cle{combustion complète} d'un alcool est sa réaction avec le dioxygène \ce{O2} qui le transforme intégralement en \cle{dioxyde de carbone} \ce{CO2} et en \cle{eau} \ce{H2O}. C'est une réaction \cle{exothermique} : elle dégage une grande quantité de chaleur. Elle est totale et rapide dès que l'alcool est enflammé.
\end{definition}

Pour un alcool saturé de formule brute générale $C_nH_{2n+1}OH$ (que l'on écrit aussi $C_nH_{2n+2}O$), l'équation-bilan générale s'écrit :

\[ C_nH_{2n+1}OH + \frac{3n}{2}\, O_2 \longrightarrow n\, CO_2 + (n+1)\, H_2O \]

Vérifions le décompte des atomes : à gauche, $n$ atomes de carbone, $(2n+2)$ atomes d'hydrogène et $(1 + 3n)$ atomes d'oxygène ; à droite, $n$ atomes de carbone dans $n\,\ce{CO2}$, $2(n+1) = 2n+2$ atomes d'hydrogène dans $(n+1)\,\ce{H2O}$, et $2n + (n+1) = 3n+1$ atomes d'oxygène. L'équation est bien équilibrée.

\begin{methode}[Équilibrer la combustion d'un alcool $C_nH_{2n+2}O$]
Pour équilibrer rapidement l'équation de combustion d'un alcool :
\begin{enumerate}
\item \textbf{Équilibrer le carbone} : écrire $n\,\ce{CO2}$ (autant de \ce{CO2} que d'atomes de carbone dans l'alcool).
\item \textbf{Équilibrer l'hydrogène} : l'alcool contient $2n+2$ atomes d'hydrogène, donc écrire $(n+1)\,\ce{H2O}$.
\item \textbf{Équilibrer l'oxygène en dernier} : compter les atomes d'oxygène à droite, soit $2n + (n+1) = 3n+1$ ; un oxygène est déjà fourni par l'alcool, il en reste $3n$ à apporter par \ce{O2}, d'où le coefficient $\dfrac{3n}{2}$ devant \ce{O2}.
\end{enumerate}
Si le coefficient $\frac{3n}{2}$ est fractionnaire ($n$ impair), on peut multiplier toute l'équation par 2 pour n'avoir que des coefficients entiers.
\end{methode}

\begin{exemplebox}[Combustion de l'éthanol]
Pour l'éthanol \ce{CH3-CH2-OH} ($n = 2$) : coefficient de \ce{O2} : $\frac{3 \times 2}{2} = 3$ ; coefficient de \ce{CO2} : $2$ ; coefficient de \ce{H2O} : $2+1 = 3$. D'où :
\[ \ce{CH3-CH2-OH + 3 O2 -> 2 CO2 + 3 H2O} \]
Cette réaction dégage environ \SI{1367}{kJ} par mole d'éthanol brûlé (enthalpie standard de combustion). Le pouvoir calorifique massique de l'éthanol est d'environ \SI{30}{kJ.g^{-1}} : c'est moins que l'essence (environ \SI{46}{kJ.g^{-1}}), mais l'éthanol présente un avantage décisif : il peut être produit par \emph{fermentation} de matières végétales (canne à sucre, maïs, manioc), donc renouvelé en permanence. C'est pourquoi on l'utilise comme \cle{biocarburant}, pur ou mélangé à l'essence (carburant E10 : 10 \% d'éthanol).
\end{exemplebox}

\begin{exoresolu}[Chaleur dégagée par la combustion de l'éthanol]
\textbf{Énoncé.} Un réchaud à alcool brûle complètement $m = \SI{23}{g}$ d'éthanol \ce{CH3-CH2-OH}. On donne $M(\ce{CH3-CH2-OH}) = \SI{46}{g.mol^{-1}}$ et l'enthalpie de combustion $\Delta_c H = \SI{-1367}{kJ.mol^{-1}}$.
\begin{enumerate}
\item Écrire l'équation-bilan de la combustion complète.
\item Calculer la quantité de matière d'éthanol brûlé et la chaleur dégagée.
\item Quel volume de dioxygène (CNTP, $V_m = \SI{22,4}{L.mol^{-1}}$) cette combustion a-t-elle consommé ?
\end{enumerate}
\tcblower
\textbf{Solution détaillée.}
\begin{enumerate}
\item Équation-bilan : \[ \ce{CH3-CH2-OH + 3 O2 -> 2 CO2 + 3 H2O} \]
\item Quantité de matière d'éthanol :
\[ n(\ce{CH3-CH2-OH}) = \frac{m}{M} = \frac{23}{46} = \SI{0,50}{mol} \]
La combustion est exothermique : la chaleur \emph{dégagée} vaut, en valeur absolue :
\[ Q = n \times |\Delta_c H| = 0,50 \times 1367 = \SI{683,5}{kJ} \]
\item D'après l'équation, 1 mol d'éthanol consomme 3 mol de dioxygène :
\[ n(\ce{O2}) = 3 \times 0,50 = \SI{1,5}{mol} \]
\[ V(\ce{O2}) = n \times V_m = 1,5 \times 22,4 = \SI{33,6}{L} \]
\end{enumerate}
La combustion de \SI{23}{g} d'éthanol dégage donc environ \SI{684}{kJ} et consomme \SI{33,6}{L} de dioxygène.
\end{exoresolu}

\begin{attention}[Combustion incomplète : un danger mortel]
Si l'apport de dioxygène est insuffisant (pièce mal aérée, réchaud encrassé), la combustion devient \cle{incomplète} : il se forme du \cle{monoxyde de carbone} \ce{CO}, gaz incolore, inodore et très toxique, voire du carbone (suie). Ne jamais utiliser un réchaud à alcool ou à gaz dans une pièce fermée sans aération. Par ailleurs, la combustion \emph{détruit} le squelette carboné : elle ne permet pas de fabriquer de nouvelles molécules organiques. Pour cela, il faut une oxydation plus douce.
\end{attention}

\courssub{La notion d'oxydation ménagée}

\textbf{Intuition.}\par
La combustion, c'est le bulldozer : elle détruit toute la molécule pour n'en tirer que de la chaleur. Mais le chimiste, lui, veut \emph{transformer} l'alcool en une molécule utile (aldéhyde, cétone, acide carboxylique) tout en gardant son architecture carbonée. Il lui faut donc un oxydant « chirurgical », qui n'attaque que le groupe hydroxyle $-\ce{OH}$ et son carbone porteur.

\begin{definition}[Oxydation ménagée]
Une \cle{oxydation ménagée} est une oxydation qui \textbf{ne rompt pas la chaîne carbonée} de la molécule : le nombre d'atomes de carbone est conservé, seul le groupe fonctionnel est modifié. Elle est réalisée par des oxydants en quantité contrôlée (permanganate de potassium \ce{KMnO4}, dichromate de potassium \ce{K2Cr2O7}, oxyde de cuivre \ce{CuO}, dioxygène en présence de catalyseur), à température modérée.
\end{definition}

\begin{center}
\begin{tabularx}{\linewidth}{|l|X|X|}
\hline
\rowcolor{popblueL} \textbf{Critère} & \textbf{Combustion} & \textbf{Oxydation ménagée} \\
\hline
Oxydant & \ce{O2} (excès, flamme) & \ce{KMnO4}, \ce{CuO}, \ce{O2}/catalyseur \\
\hline
Squelette carboné & Détruit (\ce{CO2} + \ce{H2O}) & Conservé \\
\hline
Produits & \ce{CO2} et \ce{H2O} uniquement & Aldéhyde, cétone ou acide carboxylique \\
\hline
Intérêt & Énergie (chaleur) & Synthèse de nouvelles molécules \\
\hline
\end{tabularx}
\end{center}

\courssub{Expérience de la lampe sans flamme}

\textbf{Intuition et principe.}\par
Peut-on oxyder l'éthanol avec le dioxygène de l'air, \emph{sans flamme} ? Oui, à condition de trouver un « passeur » d'oxygène : le \cle{cuivre} joue ce rôle. Chauffé à l'air, il se recouvre d'oxyde de cuivre(II) \ce{CuO} noir ; plongé dans les vapeurs d'éthanol, \ce{CuO} cède son oxygène à l'alcool et redevient du cuivre rouge brillant. La réaction dégage assez de chaleur pour que le fil reste incandescent et entretienne lui-même la réaction : c'est la \cle{lampe sans flamme}.

\begin{experience}[La lampe sans flamme]
\textbf{Dispositif et protocole.}
\begin{enumerate}
\item Verser un peu d'éthanol dans un bécher (ou un tube à essais) et chauffer doucement pour dégager des vapeurs d'alcool, \emph{sans enflammer}.
\item Réaliser une spirale avec un fil de cuivre décapé, fixée à une tige de verre.
\item Chauffer la spirale au rouge dans la flamme d'un bec Bunsen : elle se noircit (formation de \ce{CuO}).
\item Plonger aussitôt la spirale incandescente \emph{au-dessus} du liquide chaud, dans les vapeurs d'éthanol.
\end{enumerate}
\textbf{Observations.}
\begin{itemize}
\item Le fil noir \textbf{redevient rouge brillant} : \ce{CuO} est retransformé en cuivre métallique.
\item Le fil \textbf{reste incandescent} de façon durable : la réaction est exothermique et s'auto-entretient, \textbf{sans flamme}.
\item Une \textbf{odeur piquante de pomme verte} se dégage : c'est l'\cle{éthanal} \ce{CH3-CHO}.
\item De fines gouttelettes se condensent sur les parois froides : de l'\cle{eau} s'est formée.
\end{itemize}
\textbf{Interprétation.} Le cuivre transporte l'oxygène de l'air jusqu'à l'alcool : il s'oxyde en \ce{CuO} à l'air, puis \ce{CuO} oxyde l'éthanol en éthanal en se réduisant en cuivre. Le cycle recommence tant qu'il y a de l'alcool et de l'air.
\end{experience}

\begin{popfigure}
\begin{center}
\begin{tikzpicture}[scale=1.0, every node/.style={font=\small}]
% Bécher
\draw[thick, popdark] (-1.6,0) -- (-1.6,2.6) ;
\draw[thick, popdark] (1.6,0) -- (1.6,2.6);
\draw[thick, popdark] (-1.6,0) -- (1.6,0);
% Liquide éthanol
\fill[popblue!25] (-1.55,0.05) rectangle (1.55,1.0);
\draw[popblue, thick] (-1.55,1.0) -- (1.55,1.0);
\node[popblue] at (0,0.5) {éthanol liquide chaud};
% Vapeurs
\foreach \x in {-1.0,-0.3,0.4,1.0}{
  \draw[popmuted, decorate, decoration={snake, amplitude=0.6mm, segment length=3mm}] (\x,1.1) -- (\x,2.4);}
\node[popmuted, font=\footnotesize] at (0,2.75) {vapeurs d'éthanol};
% Spirale de cuivre
\draw[poporange, very thick] (0,4.6) -- (0,3.4);
\draw[poporange, very thick] (0,3.4) .. controls (0.35,3.35) and (0.35,3.15) .. (0,3.1)
  .. controls (-0.35,3.05) and (-0.35,2.85) .. (0,2.8)
  .. controls (0.35,2.75) and (0.35,2.55) .. (0,2.5)
  .. controls (-0.35,2.45) and (-0.35,2.25) .. (0,2.2);
% Incandescence
\draw[popgold, thick, dashed] (0,2.8) circle (0.75);
\node[poporange, align=center, font=\footnotesize] at (3.6,3.6) {spirale de cuivre\\ incandescente};
\draw[->, poporange] (2.6,3.5) -- (0.5,3.1);
% Chauffage
\draw[popink, thick] (-0.5,-0.55) .. controls (-0.2,-0.1) and (0.2,-0.9) .. (0,-0.35);
\draw[popink, thick] (0.5,-0.55) .. controls (0.2,-0.1) and (-0.2,-0.9) .. (0,-0.35);
\node[popink, font=\footnotesize] at (0,-1.0) {chauffage doux (sans enflammer)};
% Observations
\node[popgreen, align=left, font=\footnotesize, anchor=west] at (-6.4,2.2) {\textbf{Observations :}\\ $\bullet$ fil redevient rouge brillant\\ $\bullet$ incandescence entretenue\\ $\bullet$ odeur de pomme verte (\ce{CH3-CHO})\\ $\bullet$ buée d'eau sur les parois};
\draw[->, popgreen] (-1.8,2.2) -- (-1.65,2.2);
\end{tikzpicture}
\end{center}
\legende{Montage de la lampe sans flamme : la spirale de cuivre chauffée au rouge, plongée dans les vapeurs d'éthanol, oxyde l'alcool en éthanal sans flamme.}
\end{popfigure}

\textbf{Les équations de la réaction.}\par
Le mécanisme global se décompose en deux étapes qui forment un \cle{cycle catalytique}.

\emph{Étape 1 — oxydation du cuivre par le dioxygène de l'air} (le fil rouge devient noir) :
\[ \ce{2 Cu + O2 -> 2 CuO} \]

\emph{Étape 2 — oxydation de l'éthanol par l'oxyde de cuivre(II)} (le fil noir redevient rouge) :
\[ \ce{CH3-CH2-OH + CuO -> CH3-CHO + Cu + H2O} \]

En additionnant membre à membre (le \ce{CuO} formé en étape 1 est consommé en étape 2, et le cuivre est régénéré), on obtient l'équation-bilan de l'\cle{oxydation catalytique} de l'éthanol par le dioxygène :
\[ \ce{2 CH3-CH2-OH + O2 ->[Cu] 2 CH3-CHO + 2 H2O} \]

\begin{popfigure}
\begin{center}
\begin{tikzpicture}[node distance=1.1cm, every node/.style={align=center}]
\node (cu) {\chemfig{Cu}\quad(cuivre rouge)};
\node[right=2.6cm of cu] (cuo) {\chemfig{CuO}\quad(oxyde noir)};
\draw[->, popdark, very thick, bend left=25] (cu) to node[above, font=\footnotesize]{\ce{+ O2} (air, chauffage)} (cuo);
\draw[->, poporange, very thick, bend left=25] (cuo) to node[below, font=\footnotesize]{\ce{+ CH3-CH2-OH} $\longrightarrow$ \ce{CH3-CHO} + \ce{H2O}} (cu);
\node[below=1.5cm of cu, font=\footnotesize\itshape, popmuted] {le cuivre est régénéré : catalyseur};
\end{tikzpicture}
\end{center}
\legende{Cycle catalytique Cu/CuO : le cuivre transporte l'oxygène de l'air vers l'éthanol, qui est oxydé en éthanal \chemfig{CH_3-CH=O} (groupe carbonyle \chemfig{C=O} en bout de chaîne).}
\end{popfigure}

\textbf{Identification des produits.}\par
Pour prouver la nature des produits formés, on réalise deux tests caractéristiques :
\begin{itemize}
\item \textbf{L'éthanal} (aldéhyde) : quelques gouttes de \cle{liqueur de Fehling} (bleue) ajoutées au distillat, chauffées au bain-marie, donnent un \textbf{précipité rouge brique} d'oxyde de cuivre(I) \ce{Cu2O}. On peut aussi utiliser la \cle{2,4-DNPH} : il se forme un \textbf{précipité jaune-orangé}, caractéristique du groupe carbonyle \ce{C=O} (aldéhydes et cétones). Le test à la Fehling, positif, prouve en plus qu'il s'agit d'un aldéhyde.
\item \textbf{L'eau} : elle fait passer le \cle{sulfate de cuivre anhydre} (blanc) au \textbf{bleu} (sulfate de cuivre pentahydraté \ce{CuSO4.5H2O}).
\end{itemize}

\begin{propriete}[Rôle catalytique du cuivre]
Le cuivre est un \cle{catalyseur} : il accélère l'oxydation de l'éthanol par le dioxygène, il participe au mécanisme (cycle \ce{Cu}/\ce{CuO}), mais il est \textbf{intégralement régénéré en fin de réaction}. Il n'apparaît donc pas dans l'équation-bilan ; on l'indique au-dessus de la flèche : \ce{->[Cu]}. La réaction est en outre \textbf{exothermique} : la chaleur dégagée maintient le fil incandescent, d'où le nom de « lampe sans flamme ».
\end{propriete}

\begin{attention}[Erreur fréquente : « le cuivre est consommé »]
Non ! Dire que le cuivre est un réactif consommé est l'erreur classique du baccalauréat. Le cuivre est certes \emph{transformé} temporairement en \ce{CuO}, mais il est \textbf{régénéré} à la fin de chaque cycle : sa masse ne varie pas au cours de l'expérience. C'est un \cle{catalyseur}. Dans l'équation-bilan, seuls l'éthanol et le dioxygène sont les réactifs ; l'éthanal et l'eau sont les produits. Autre piège : ne pas confondre cette oxydation ménagée (qui donne \ce{CH3-CHO}, 2 atomes de carbone conservés) avec la combustion (qui donne \ce{CO2} et détruit la chaîne).
\end{attention}

\begin{exoresolu}[Bilan de matière de la lampe sans flamme]
\textbf{Énoncé.} On réalise l'expérience de la lampe sans flamme avec $n = \SI{0,10}{mol}$ d'éthanol, supposé totalement transformé en éthanal selon : \ce{2 CH3-CH2-OH + O2 ->[Cu] 2 CH3-CHO + 2 H2O}.
\begin{enumerate}
\item Calculer la masse d'éthanal obtenue. On donne $M(\ce{CH3CHO}) = \SI{44}{g.mol^{-1}}$.
\item Un élève pèse la spirale de cuivre avant et après l'expérience : il trouve la même masse. Est-ce surprenant ? Justifier.
\end{enumerate}
\tcblower
\textbf{Solution détaillée.}
\begin{enumerate}
\item D'après l'équation, 2 mol d'éthanol donnent 2 mol d'éthanal, soit une proportion 1 pour 1 :
\[ n(\ce{CH3CHO}) = n(\ce{C2H5OH}) = \SI{0,10}{mol} \]
\[ m(\ce{CH3CHO}) = n \times M = 0,10 \times 44 = \SI{4,4}{g} \]
\item Non, ce n'est pas surprenant : c'est même la preuve expérimentale du \textbf{rôle catalytique} du cuivre. Oxydé en \ce{CuO} puis réduit en \ce{Cu} à chaque cycle, le cuivre est intégralement régénéré : sa masse est conservée. Si le cuivre avait été un réactif, sa masse aurait diminué.
\end{enumerate}
\end{exoresolu}

\begin{savaistu}[Des éthylotests aux pots d'échappement]
Le principe de l'oxydation catalytique de l'alcool a des applications quotidiennes. Les anciens \cle{éthylotests} chimiques contenaient du dichromate de potassium orangé qui, en oxydant l'éthanol de l'air expiré en acide éthanoïque, virait au vert : la couleur trahissait l'alcoolémie. De leur côté, les \cle{pots catalytiques} des automobiles utilisent des métaux catalyseurs (platine, palladium, rhodium) pour achever l'oxydation des vapeurs de carburant imbrûlées et du monoxyde de carbone en \ce{CO2} et \ce{H2O} — exactement la même idée que notre fil de cuivre : un catalyseur qui « passe » l'oxygène sans jamais se consommer. Au Cameroun, la maîtrise de la fermentation puis de la distillation de l'éthanol (à partir de la canne à sucre ou du manioc) intéresse les filières des biocarburants : comprendre l'oxydation de l'éthanol, c'est comprendre à la fois comment il brûle dans un moteur et comment on le transforme en chimie fine.
\end{savaistu}

\begin{aretenir}[L'essentiel de la section]
\begin{itemize}
\item La \cle{combustion complète} d'un alcool $C_nH_{2n+1}OH$ : $C_nH_{2n+1}OH + \frac{3n}{2} O_2 \longrightarrow n CO_2 + (n+1) H_2O$, réaction exothermique qui détruit le squelette carboné (énergie, biocarburants).
\item Une \cle{oxydation ménagée} conserve la chaîne carbonée et ne modifie que le groupe fonctionnel.
\item La \cle{lampe sans flamme} : \ce{2 Cu + O2 -> 2 CuO} puis \ce{CH3-CH2-OH + CuO -> CH3-CHO + Cu + H2O} ; bilan : \ce{2 CH3-CH2-OH + O2 ->[Cu] 2 CH3-CHO + 2 H2O}.
\item Observations clés : fil qui redevient rouge brillant, incandescence entretenue (réaction exothermique), odeur de pomme verte (éthanal), buée (eau).
\item Identification : éthanal par la liqueur de Fehling (précipité rouge brique) ou la DNPH (précipité jaune) ; eau par le sulfate de cuivre anhydre (blanc $\to$ bleu).
\item Le cuivre est un \cle{catalyseur} : régénéré en fin de réaction, sa masse ne change pas.
\end{itemize}
\end{aretenir}

Cette expérience nous a montré qu'un alcool \emph{primaire} comme l'éthanol s'oxyde d'abord en \emph{aldéhyde}. Que se passe-t-il alors avec un alcool secondaire ou tertiaire ? L'oxydation ménagée dépend-elle de la classe de l'alcool ? C'est précisément l'objet de la section suivante.

\coursec{Oxydation ménagée selon la classe de l'alcool}

Dans la section précédente, nous avons vu que la combustion détruit totalement la molécule d'alcool en la réduisant en \ce{CO2} et \ce{H2O}, et que la lampe sans flamme révélait déjà une oxydation plus douce, catalysée par le cuivre incandescent. Mais il existe une troisième voie, bien plus précieuse pour le chimiste : l'\cle{oxydation ménagée}. L'adjectif « ménagée » dit tout : on ménage le squelette carboné, on ne touche qu'au groupe fonctionnel. C'est exactement ce que font, sans le savoir, les vinaigriers de Bafoussam ou de Maroua lorsqu'ils laissent du vin de palme « tourner » à l'air libre : l'alcool se transforme doucement en acide, sans que la chaîne carbonée ne soit brisée. Toute cette section repose sur une idée simple et puissante : \textbf{le résultat de l'oxydation ménagée dépend de la classe de l'alcool}.

\courssub{L'oxydant de référence : le permanganate de potassium en milieu acide}

\begin{definition}[Oxydation ménagée]
Une \cle{oxydation ménagée} est une oxydation qui modifie le groupe fonctionnel d'une molécule organique \textbf{sans rompre aucune liaison carbone--carbone}. Le squelette carboné est conservé ; seul le degré d'oxydation du carbone fonctionnel augmente.
\end{definition}

L'oxydant le plus utilisé au laboratoire est le \cle{permanganate de potassium} \ce{KMnO4}, employé en \textbf{solution acidifiée} (quelques gouttes d'acide sulfurique dilué). Son intérêt est double :

\begin{itemize}
\item c'est un oxydant \textbf{puissant}, capable d'oxyder la plupart des alcools ;
\item c'est un oxydant \textbf{auto-indicateur} : l'ion permanganate \ce{MnO4-} est \textbf{violet} intense, tandis que l'ion manganèse \ce{Mn^2+} auquel il se réduit est \textbf{quasiment incolore}. La décoloration de la solution signe donc visuellement la réaction.
\end{itemize}

Le couple oxydant/réducteur mis en jeu est \ce{MnO4-}/\ce{Mn^2+}, dont la demi-équation électronique en milieu acide s'écrit :

\[ \ce{MnO4- + 8H+ + 5e- -> Mn^2+ + 4H2O} \]

\begin{aretenir}[Le test au permanganate]
Si une solution acidifiée de \ce{KMnO4} versée sur un alcool \textbf{se décolore}, l'alcool a été oxydé. Si elle \textbf{reste violette}, aucune oxydation ménagée n'a eu lieu. Ce simple changement de couleur est un test d'identification de première importance.
\end{aretenir}

\courssub{Oxydation d'un alcool primaire : aldéhyde puis acide carboxylique}

\textbf{Intuition.} Un alcool primaire possède son groupe $-\ce{OH}$ sur un carbone de bout de chaîne, qui porte encore \textbf{deux atomes d'hydrogène}. L'oxydant peut donc lui « arracher » deux hydrogènes successivement : d'abord un premier stade (l'aldéhyde), puis, si l'oxydant est en excès, un second stade (l'acide carboxylique). C'est une véritable \textbf{cascade d'oxydation}.

\begin{definition}[Aldéhyde et cétone]
Un \cle{aldéhyde} possède le groupe carbonyle \ce{C=O} \textbf{en bout de chaîne} : le groupe fonctionnel est $-\ce{CHO}$ (le carbone du carbonyle porte au moins un hydrogène). Une \cle{cétone} possède le groupe carbonyle \ce{C=O} sur un carbone \textbf{non terminal}, lié à deux autres atomes de carbone.
\end{definition}

Prenons l'éthanol \chemfig{CH_3-CH_2-OH}, alcool primaire par excellence.

\textbf{Premier stade} (oxydant en défaut ou distillation immédiate de l'aldéhyde) : formation de l'\textbf{éthanal} :

\[ \ce{CH3-CH2-OH ->[O] CH3-CHO + H2O} \]

\textbf{Second stade} (oxydant en excès, à chaud) : l'aldéhyde est à son tour oxydé en \textbf{acide éthanoïque} :

\[ \ce{CH3-CHO ->[O] CH3-COOH} \]

\begin{popfigure}
\begin{center}
\begin{tikzpicture}[node distance=1.1cm and 1.6cm]
\node (a) {\chemfig{CH_3-CH_2-OH}};
\node[below=0.15cm of a, font=\small\itshape, text=popmuted] {éthanol (alcool primaire)};
\node[right=of a] (b) {\chemfig{CH_3-C(=[1]O)-[7]H}};
\node[below=0.15cm of b, font=\small\itshape, text=popmuted] {éthanal (aldéhyde)};
\node[right=of b] (c) {\chemfig{CH_3-C(=[1]O)-[7]OH}};
\node[below=0.15cm of c, font=\small\itshape, text=popmuted] {acide éthanoïque};
\draw[-{Stealth[length=3mm]}, very thick, poporange] (a) -- (b) node[midway, above, font=\small\bfseries, text=popdark] {[O]} node[midway, below, font=\scriptsize, text=popmuted] {oxydant en défaut};
\draw[-{Stealth[length=3mm]}, very thick, poporange] (b) -- (c) node[midway, above, font=\small\bfseries, text=popdark] {[O]} node[midway, below, font=\scriptsize, text=popmuted] {oxydant en excès};
\end{tikzpicture}
\end{center}
\legende{Cascade d'oxydation ménagée de l'éthanol : alcool primaire $\rightarrow$ aldéhyde $\rightarrow$ acide carboxylique.}
\end{popfigure}

\begin{methode}[Écrire l'équation-bilan d'oxydation par \ce{MnO4-} en milieu acide]
Pour obtenir l'équation-bilan ionique de l'oxydation d'un alcool primaire en aldéhyde :
\begin{enumerate}
\item Écrire la demi-équation de \textbf{réduction} de l'oxydant : \ce{MnO4- + 8H+ + 5e- -> Mn^2+ + 4H2O}.
\item Écrire la demi-équation d'\textbf{oxydation} de l'alcool : \ce{R-CH2OH -> R-CHO + 2H+ + 2e-}.
\item Multiplier chaque demi-équation pour égaliser les électrons : ici $\times 2$ pour la première, $\times 5$ pour la seconde (10 électrons échangés).
\item Additionner membre à membre et \textbf{simplifier} les \ce{H+} et \ce{H2O} présents des deux côtés.
\item Vérifier la conservation des atomes \emph{et} des charges.
\end{enumerate}
\end{methode}

Appliquons la méthode à l'éthanol. Les demi-équations sont :

\begin{align*}
\text{Réduction : } & \ce{MnO4- + 8H+ + 5e- -> Mn^2+ + 4H2O} \quad (\times 2)\\
\text{Oxydation : } & \ce{CH3CH2OH -> CH3CHO + 2H+ + 2e-} \quad (\times 5)
\end{align*}

Après combinaison et simplification :

\[ \ce{2MnO4- + 6H+ + 5CH3CH2OH -> 2Mn^2+ + 5CH3CHO + 8H2O} \]

Si l'oxydant est en excès, l'oxydation va jusqu'à l'acide carboxylique. La demi-équation d'oxydation devient alors \ce{CH3CH2OH + H2O -> CH3COOH + 4H+ + 4e-}, ce qui conduit, après combinaison avec la réduction du permanganate ($\times 4$ et $\times 5$ pour 20 électrons), à :

\[ \ce{4MnO4- + 12H+ + 5CH3CH2OH -> 4Mn^2+ + 5CH3COOH + 11H2O} \]

\begin{exoresolu}[Stœchiométrie de l'oxydation complète de l'éthanol]
\textbf{Énoncé.} On veut oxyder complètement $n = \SI{1,0}{mol}$ d'éthanol en acide éthanoïque à l'aide d'une solution acidifiée de permanganate de potassium. Quelle quantité d'ions permanganate \ce{MnO4-} faut-il utiliser au minimum ?
\tcblower
\textbf{Solution.} D'après l'équation-bilan établie ci-dessus :
\[ \ce{4MnO4- + 12H+ + 5CH3CH2OH -> 4Mn^2+ + 5CH3COOH + 11H2O} \]
le rapport stœchiométrique est $\dfrac{n(\ce{MnO4-})}{n(\ce{CH3CH2OH})} = \dfrac{4}{5}$. Donc :
\[ n(\ce{MnO4-}) = \frac{4}{5} \times n(\text{éthanol}) = \frac{4}{5} \times 1{,}0 = \SI{0,80}{mol} \]
\textbf{Conclusion :} l'oxydation complète de \SI{1}{mol} d'éthanol en acide éthanoïque nécessite $\dfrac{4}{5}$ mol, soit \SI{0,80}{mol}, d'ions permanganate en milieu acide. On retiendra que chaque éthanol cède \textbf{4 électrons} (deux pour passer à l'aldéhyde, deux pour passer à l'acide), tandis que chaque \ce{MnO4-} en capte 5 : d'où le rapport $4/5$.
\end{exoresolu}

\courssub{Oxydation d'un alcool secondaire : la cétone}

Un alcool secondaire porte son groupe $-\ce{OH}$ sur un carbone lié à \textbf{deux autres carbones} et à \textbf{un seul hydrogène}. L'oxydant ne peut donc arracher qu'une seule « paire » d'hydrogènes : l'oxydation s'arrête au stade de la \textbf{cétone}, qui n'est pas oxydable davantage de manière ménagée (il faudrait rompre une liaison \ce{C-C}, ce qui serait une oxydation brutale).

Exemple de référence : le \textbf{propan-2-ol} \chemfig{H_3C-CH(-[2]OH)-CH_3} s'oxyde en \textbf{propanone} (acétone) \chemfig{H_3C-C(=[1]O)-[7]CH_3} :

\[ \ce{CH3-CHOH-CH3 ->[O] CH3-CO-CH3 + H2O} \]

Avec le permanganate en milieu acide, la demi-équation d'oxydation est \ce{CH3CHOHCH3 -> CH3COCH3 + 2H+ + 2e-}, et l'équation-bilan s'écrit :

\[ \ce{2MnO4- + 6H+ + 5CH3CHOHCH3 -> 2Mn^2+ + 5CH3COCH3 + 8H2O} \]

\begin{aretenir}[Bilan des trois classes]
Alcool \textbf{primaire} $\rightarrow$ aldéhyde $\rightarrow$ acide carboxylique (si excès d'oxydant). Alcool \textbf{secondaire} $\rightarrow$ cétone (point final). Alcool \textbf{tertiaire} $\rightarrow$ rien du tout.
\end{aretenir}

\courssub{Alcool tertiaire : l'oxydation impossible}

Le carbone fonctionnel d'un alcool tertiaire, comme le 2-méthylpropan-2-ol \chemfig{H_3C-C(-[2]CH_3)(-[6]OH)-CH_3}, est lié à \textbf{trois atomes de carbone} : il ne porte \textbf{aucun hydrogène}. Or l'oxydation ménagée consiste précisément à retirer l'hydrogène du carbone fonctionnel (avec celui du $-\ce{OH}$) pour créer la double liaison \ce{C=O}. Sans hydrogène disponible, la réaction est impossible : le permanganate \textbf{reste violet}.

\begin{attention}[Le piège classique du baccalauréat]
Erreur fréquente : écrire qu'un alcool tertiaire s'oxyde en « cétone » ou en « acide ». C'est faux ! Sans hydrogène sur le carbone fonctionnel, \textbf{aucune oxydation ménagée n'est possible}. Si l'on chauffe fortement avec un oxydant concentré, on obtient seulement une \emph{dégradation} du squelette (oxydation brutale), hors programme. Au contraire, la persistance de la couleur violette du \ce{KMnO4} est le \textbf{test caractéristique} d'un alcool tertiaire.
\end{attention}

\begin{popfigure}
\begin{center}
\begin{tikzpicture}[font=\small]
\draw[fill=popblueL, rounded corners=2pt] (0,0) rectangle (3.6,3.4);
\draw[fill=popgreenL, rounded corners=2pt] (4.2,0) rectangle (7.8,3.4);
\draw[fill=poppinkL, rounded corners=2pt] (8.4,0) rectangle (12,3.4);
\node[font=\small\bfseries, text=popblue] at (1.8,3.05) {Alcool primaire};
\node[font=\small\bfseries, text=popgreen!60!black] at (6.0,3.05) {Alcool secondaire};
\node[font=\small\bfseries, text=poppink] at (10.2,3.05) {Alcool tertiaire};
\node[align=center] at (1.8,2.35) {$R\text{-}\mathrm{CH_2}\text{-}OH$};
\node[align=center] at (6.0,2.35) {$R\text{-}\mathrm{CHOH}\text{-}R'$};
\node[align=center] at (10.2,2.35) {$R\text{-}\mathrm{C(OH)(R')}\text{-}R''$};
\node[align=center, font=\footnotesize] at (1.8,1.55) {aldéhyde puis\\ \textbf{acide carboxylique}};
\node[align=center, font=\footnotesize] at (6.0,1.55) {\textbf{cétone}};
\node[align=center, font=\footnotesize] at (10.2,1.55) {\textbf{pas d'oxydation}\\ ménagée};
\draw[fill=popcream, draw=popline, rounded corners=1pt] (0.5,0.25) rectangle (3.1,0.95);
\draw[fill=white, draw=popline, rounded corners=1pt] (0.7,0.4) rectangle (1.5,0.8);
\draw[fill=popmuted!30, draw=popline, rounded corners=1pt] (1.7,0.4) rectangle (2.9,0.8);
\node[font=\scriptsize] at (1.1,0.6) {violet};
\node[font=\scriptsize] at (2.3,0.6) {incolore};
\draw[fill=popcream, draw=popline, rounded corners=1pt] (4.7,0.25) rectangle (7.3,0.95);
\draw[fill=white, draw=popline, rounded corners=1pt] (4.9,0.4) rectangle (5.7,0.8);
\draw[fill=popmuted!30, draw=popline, rounded corners=1pt] (5.9,0.4) rectangle (7.1,0.8);
\node[font=\scriptsize] at (5.3,0.6) {violet};
\node[font=\scriptsize] at (6.5,0.6) {incolore};
\draw[fill=popcream, draw=popline, rounded corners=1pt] (8.9,0.25) rectangle (11.5,0.95);
\draw[fill=white, draw=popline, rounded corners=1pt] (9.1,0.4) rectangle (9.9,0.8);
\draw[fill=white, draw=popline, rounded corners=1pt] (10.1,0.4) rectangle (11.3,0.8);
\node[font=\scriptsize] at (9.5,0.6) {violet};
\node[font=\scriptsize] at (10.7,0.6) {reste violet};
\end{tikzpicture}
\end{center}
\legende{Comportement comparatif des trois classes d'alcools face à \ce{KMnO4} acidifié : décoloration pour les alcools primaires et secondaires, persistance du violet pour l'alcool tertiaire.}
\end{popfigure}

\courssub{Oxydation ménagée des aldéhydes et tests d'identification}

Nous avons vu que l'aldéhyde, produit intermédiaire de l'oxydation d'un alcool primaire, est lui-même oxydable en acide carboxylique :

\[ \ce{R-CHO ->[O] R-COOH} \]

Cette grande facilité d'oxydation fait des aldéhydes des \textbf{réducteurs remarquables}, à la différence des cétones qui, elles, ne sont pas oxydables de manière ménagée. C'est précisément cette différence qui fonde les \textbf{tests caractéristiques} permettant de distinguer un aldéhyde d'une cétone :

\begin{itemize}
\item \textbf{Liqueur de Fehling} (solution bleue d'ions cuivre(II) en milieu basique) : à chaud, un aldéhyde la réduit en oxyde cuivreux \ce{Cu2O}, qui forme un \textbf{précipité rouge brique}. Une cétone ne réagit pas.
\item \textbf{Réactif de Tollens} (nitrate d'argent ammoniacal) : un aldéhyde réduit les ions argent en argent métallique qui se dépose sur les parois du tube : c'est le \textbf{miroir d'argent}.
\item \textbf{DNPH} (2,4-dinitrophénylhydrazine) : donne un \textbf{précipité jaune-orangé} avec les aldéhydes \emph{et} les cétones (tous les composés carbonylés). Ce test prouve la présence du groupe \ce{C=O} mais ne distingue pas aldéhyde et cétone : il faut le compléter par Fehling ou Tollens.
\end{itemize}

\begin{attention}[Bien choisir son test]
La DNPH est positive pour les deux familles : ne conclus jamais « c'est un aldéhyde » sur la seule base d'un précipité jaune à la DNPH ! La démarche correcte : DNPH positif $\Rightarrow$ composé carbonylé ; puis Fehling (ou Tollens) positif $\Rightarrow$ aldéhyde ; Fehling négatif $\Rightarrow$ cétone.
\end{attention}

\begin{center}
\rowcolors{1}{popcream}{white}
\begin{tabularx}{\linewidth}{|l|X|X|X|}
\hline
\rowcolor{popblueL} \textbf{Classe de l'alcool} & \textbf{Produit d'oxydation ménagée} & \textbf{Observation avec \ce{KMnO4}/\ce{H+}} & \textbf{Tests d'identification du produit} \\
\hline
Primaire $R-\mathrm{CH_2OH}$ & Aldéhyde $R-\mathrm{CHO}$ puis acide $R-\mathrm{COOH}$ (excès d'oxydant) & Décoloration (violet $\rightarrow$ incolore) & DNPH : précipité jaune ; Fehling : précipité rouge brique ; Tollens : miroir d'argent \\
\hline
Secondaire $R-\mathrm{CHOH}-R'$ & Cétone $R-\mathrm{CO}-R'$ & Décoloration (violet $\rightarrow$ incolore) & DNPH : précipité jaune ; Fehling et Tollens : négatifs \\
\hline
Tertiaire $R-\mathrm{C(OH)(R')-R''}$ & Aucun (oxydation ménagée impossible) & La solution reste violette & Aucun produit à tester \\
\hline
\end{tabularx}
\end{center}

\begin{savaistu}[Du vin de palme au vinaigre : la chimie du quotidien]
Le vinaigre de table n'est rien d'autre qu'une solution d'acide éthanoïque à environ 6 à 8~\% . Il résulte de l'\textbf{oxydation ménagée de l'éthanol} contenu dans le vin, le cidre ou la bière, réalisée par des bactéries du genre \emph{Acetobacter} en présence de l'oxygène de l'air : \ce{CH3CH2OH + O2 -> CH3COOH + H2O}. Au Cameroun, un vin de palme oublié à l'air libre « tourne » en quelques jours et devient piquant : c'est exactement cette réaction, une oxydation ménagée biologique ! C'est aussi pour bloquer ces bactéries que l'on bouche soigneusement les bouteilles de vin et que le vinaigrier artisanal laisse, lui, ses récipients ouverts à l'air.
\end{savaistu}

\begin{exercice}[À toi de jouer]
On dispose de trois flacons étiquetés A, B, C contenant chacun l'un des alcools suivants : butan-1-ol, butan-2-ol, 2-méthylpropan-2-ol. On verse dans chacun quelques gouttes de \ce{KMnO4} acidifié : A et C décolorent la solution, B ne la décolore pas. Le produit obtenu avec A donne un précipité rouge brique avec la liqueur de Fehling ; celui obtenu avec C donne un précipité jaune avec la DNPH mais ne réagit pas avec Fehling. Identifier A, B et C en justifiant chaque étape.
\end{exercice}

\begin{corrige}[Exercice 1]
B ne décolore pas le permanganate : c'est l'alcool \textbf{tertiaire}, le 2-méthylpropan-2-ol. A s'oxyde et son produit réduit la liqueur de Fehling : c'est un \textbf{aldéhyde}, donc A est l'alcool \textbf{primaire}, le butan-1-ol (oxydé en butanal). C s'oxyde en un composé carbonylé (DNPH positif) qui n'est pas un aldéhyde (Fehling négatif) : c'est une \textbf{cétone}, donc C est l'alcool \textbf{secondaire}, le butan-2-ol (oxydé en butanone).
\end{corrige}

\coursec{Estérification}

Dans la section précédente, nous avons vu que l'oxydation ménagée d'un alcool primaire conduit, en passant par l'aldéhyde, à un \cle{acide carboxylique}. Nous disposons donc désormais de deux grandes familles de composés oxygénés : les \cle{alcools} et les \cle{acides carboxyliques}. Que se passe-t-il lorsqu'on les met en présence ? La réponse va te surprendre par son importance pratique : on obtient une molécule à l'odeur agréable, un \cle{ester}, responsable de nombreux parfums de fruits. C'est la réaction d'\textbf{estérification}, une des réactions les plus célèbres de la chimie organique.

\courssub{Intuition et définition}

\textbf{Une situation concrète.}\par

Dans certaines savonneries artisanales et jus de fruits du Cameroun, les arômes de banane, d'ananas ou de pomme ne proviennent pas toujours des fruits eux-mêmes : ce sont des molécules synthétisées en laboratoire, les esters, qui imitent ces parfums. Or un ester se fabrique très simplement, en chauffant un mélange d'acide carboxylique et d'alcool. C'est exactement ce que tu réaliseras en TP : chauffer de l'acide éthanoïque avec de l'éthanol en présence de quelques gouttes d'acide sulfurique concentré dégage une odeur fruitée caractéristique, bien différente de l'odeur piquante du vinaigre de départ.

\begin{definition}[Estérification]
L'\cle{estérification} est la réaction chimique entre un \textbf{acide carboxylique} $\mathrm{R{-}COOH}$ et un \textbf{alcool} $\mathrm{R'{-}OH}$, qui conduit à la formation d'un \textbf{ester} $\mathrm{R{-}COO{-}R'}$ et d'\textbf{eau}.
\[ \mathrm{R{-}COOH + R'{-}OH <=> R{-}COO{-}R' + H_2O} \]
Cette réaction est \textbf{lente}, \textbf{limitée} (c'est un équilibre chimique) et \textbf{athermique} (elle ne dégage ni n'absorbe pratiquement de chaleur).
\end{definition}

Observons de près le mécanisme global : le groupe $-\mathrm{OH}$ de l'acide carboxylique et l'atome d'hydrogène du groupe hydroxyle de l'alcool s'éliminent ensemble pour former une molécule d'eau. Les deux fragments restants se lient alors par une nouvelle liaison $\mathrm{C-O}$. On retiendra la phrase mnémotechnique : \emph{« l'acide perd OH, l'alcool perd H »}.

\begin{attention}[Qui perd quoi ?]
Une erreur classique consiste à croire que c'est l'alcool qui fournit le groupe $-\mathrm{OH}$ de l'eau. Les expériences avec des isotopes marqués ($^{18}\mathrm{O}$) ont montré que c'est bien l'\textbf{acide carboxylique} qui fournit le $-\mathrm{OH}$, et l'\textbf{alcool} qui fournit le $-\mathrm{H}$. Conséquence pratique : dans l'ester $\mathrm{R-COO-R'}$, le groupe alkyle $\mathrm{R'}$ vient de l'alcool, et le reste $\mathrm{R-COO}$ vient de l'acide.
\end{attention}

\courssub{Exemple fondamental : formation de l'éthanoate d'éthyle}

Faisons réagir l'acide éthanoïque (l'acide du vinaigre) avec l'éthanol. L'équation-bilan s'écrit :

\[ \ce{CH3COOH + CH3CH2OH <=> CH3COOCH2CH3 + H2O} \]

L'ester obtenu est l'\textbf{éthanoate d'éthyle}, un liquide à l'odeur fruitée agréable, utilisé comme solvant et arôme.

\begin{popfigure}
\centering
\begin{tikzpicture}[scale=0.9, transform shape]
\node at (0,0) {\chemfig{CH_3-C(=[1]O)(-[7]OH)}};
\node[below=0.15cm of {(0,-0.9)}, popdark] {\footnotesize acide éthanoïque};
\node at (2.6,0) {\Large $+$};
\node at (4.6,0) {\chemfig{CH_3-CH_2-OH}};
\node[popdark] at (4.6,-1.15) {\footnotesize éthanol};
\draw[-{Stealth[length=3mm]}, very thick, popblue] (6.2,0.25) -- (8.4,0.25) node[midway, above, popdark]{\footnotesize $\mathrm{H^+}$, chauffage};
\draw[-{Stealth[length=3mm]}, very thick, poporange] (8.4,-0.35) -- (6.2,-0.35) node[midway, below, popdark]{\footnotesize hydrolyse};
\node at (10.2,0) {\chemfig{CH_3-C(=[1]O)(-[7]O-CH_2-CH_3)}};
\node[popdark] at (10.2,-1.35) {\footnotesize éthanoate d'éthyle};
\node at (13.1,0) {\Large $+$};
\node at (14.3,0) {\chemfig{H_2O}};
\node[popdark] at (14.3,-1.15) {\footnotesize eau};
\end{tikzpicture}
\legende{Équilibre d'estérification entre l'acide éthanoïque et l'éthanol : formation de l'éthanoate d'éthyle et d'eau.}
\end{popfigure}

\begin{attention}[Double flèche obligatoire]
L'erreur la plus fréquente au baccalauréat est d'écrire l'estérification avec une flèche simple \ce{->}. C'est \textbf{faux} : l'estérification est une réaction \textbf{limitée} qui conduit à un état d'équilibre. On doit donc toujours utiliser la double flèche \ce{<=>}. Écrire une flèche simple revient à affirmer que la réaction est totale, ce qui est contraire à l'expérience.
\end{attention}

\courssub{Les trois caractéristiques de l'estérification}

\begin{propriete}[Caractéristiques de l'estérification]
\begin{enumerate}
\item \textbf{Lente} : à température ambiante, l'équilibre n'est atteint qu'au bout de plusieurs jours, voire plusieurs semaines. Même à chaud, il faut plusieurs heures.
\item \textbf{Limitée} : la réaction n'est jamais totale. Elle conduit à un \textbf{équilibre chimique} dans lequel coexistent les réactifs et les produits. La réaction inverse, l'\textbf{hydrolyse} de l'ester, se produit simultanément :
\[ \text{R-COO-R'} + \ce{H2O} \rightleftharpoons \text{R-COOH} + \text{R'-OH} \]
\item \textbf{Athermique} : la réaction ne dégage pratiquement pas de chaleur ($\Delta H \approx 0$). Conséquence capitale : \textbf{la température ne modifie pas la composition de l'équilibre} ; elle ne fait qu'accélérer son établissement.
\end{enumerate}
\end{propriete}

\textbf{Pourquoi un équilibre ?}\par

Dès que des molécules d'ester et d'eau apparaissent dans le milieu, elles peuvent réagir entre elles pour régénérer l'acide et l'alcool : c'est l'\cle{hydrolyse}. Au début, l'estérification domine car l'ester et l'eau sont rares. Puis la vitesse d'hydrolyse augmente à mesure que les produits s'accumulent, jusqu'au moment où les deux réactions, directe et inverse, se déroulent à la même vitesse : les quantités n'évoluent plus macroscopiquement. C'est l'\textbf{état d'équilibre}, un équilibre \emph{dynamique}.

\begin{exemplebox}[Le rendement limite de 67 \%]
Partons d'un mélange équimolaire : $n = \SI{1,0}{mol}$ d'acide éthanoïque et $n = \SI{1,0}{mol}$ d'éthanol. À l'équilibre, on constate expérimentalement qu'il s'est formé environ $\SI{0,67}{mol}$ d'ester. Le rendement de l'estérification, pour un alcool primaire en proportions stœchiométriques, est donc d'environ :
\[ \eta = \frac{n_{\text{ester formé}}}{n_{\text{ester théorique}}} = \frac{0{,}67}{1{,}0} \approx 67\ \% \]
Il reste à l'équilibre environ $\SI{0,33}{mol}$ d'acide et $\SI{0,33}{mol}$ d'alcool. Avec un alcool secondaire, le rendement limite tombe à environ 60 \%, et à peine 5 \% avec un alcool tertiaire (encombrement stérique).
\end{exemplebox}

\courssub{Comment améliorer le rendement et la vitesse ?}

Puisque l'estérification est lente et limitée, le chimiste dispose de deux types de leviers : ceux qui \textbf{accélèrent} l'atteinte de l'équilibre (sans le déplacer) et ceux qui \textbf{déplacent} l'équilibre dans le sens de la formation de l'ester.

\begin{propriete}[Moyens d'action sur l'estérification]
\textbf{Pour accélérer la réaction (sans modifier le rendement à l'équilibre) :}
\begin{itemize}
\item \textbf{Élever la température} (chauffage à reflux) : la réaction est plus rapide, mais l'équilibre final reste le même car la réaction est athermique.
\item \textbf{Ajouter un catalyseur} : quelques gouttes d'\textbf{acide sulfurique concentré} (ions \ce{H+}) accélèrent à la fois l'estérification et l'hydrolyse. Le catalyseur ne déplace pas l'équilibre, il permet seulement de l'atteindre plus vite.
\end{itemize}
\textbf{Pour améliorer le rendement (déplacer l'équilibre vers l'ester) :}
\begin{itemize}
\item Introduire l'un des réactifs \textbf{en excès} (par exemple l'alcool, souvent le moins cher) : d'après la loi de Le Chatelier, l'équilibre se déplace dans le sens de la consommation de l'excès, donc vers la formation de l'ester.
\item \textbf{Éliminer l'eau} (ou l'ester) \textbf{au fur et à mesure} de sa formation, par distillation : l'équilibre se déplace continuellement pour régénérer le produit retiré, et le rendement peut approcher 100 \%.
\end{itemize}
\end{propriete}

\begin{attention}[Catalyseur et température ne déplacent pas l'équilibre]
Piège récurrent : croire que chauffer ou ajouter de l'acide sulfurique augmente le rendement. C'est faux ! Ces deux moyens ne font qu'\textbf{accélérer} l'établissement de l'équilibre. Seuls l'excès d'un réactif ou l'élimination d'un produit augmentent réellement le rendement. Retiens bien la distinction : \emph{cinétique} (vitesse) d'un côté, \emph{thermodynamique} (position de l'équilibre) de l'autre.
\end{attention}

\courssub{Le chauffage à reflux : le montage du TP}

Pour réaliser l'estérification au laboratoire, on chauffe le mélange acide + alcool + quelques gouttes d'acide sulfurique concentré. Mais ces espèces sont volatiles : chauffer dans un récipient ouvert ferait s'échapper les réactifs ! On utilise donc un \textbf{chauffage à reflux} : les vapeurs se condensent dans un réfrigérant vertical et retombent dans le ballon. On accélère ainsi la réaction à chaud, sans aucune perte de matière.

\begin{popfigure}
\centering
\begin{tikzpicture}[scale=1.0, transform shape]
% Ballon
\draw[very thick, popdark] (0,0) circle (1.1);
\draw[very thick, popdark] (-0.35,1.02) -- (-0.35,1.9);
\draw[very thick, popdark] (0.35,1.02) -- (0.35,1.9);
% Liquide dans le ballon
\fill[poporangeL] (-1.02,-0.35) arc (200:340:1.06) -- cycle;
\draw[poporange, thick] (-1.02,-0.35) arc (200:340:1.06);
\node[popdark] at (0,-0.55) {\footnotesize mélange réactionnel};
% Réfrigérant
\draw[very thick, popdark] (-0.45,1.9) -- (-0.45,5.2);
\draw[very thick, popdark] (0.45,1.9) -- (0.45,5.2);
\draw[very thick, popdark] (-0.45,5.2) arc (180:0:0.45);
% Serpentin
\draw[thick, popblue, decorate, decoration={coil, segment length=6pt, amplitude=4pt}] (0,2.1) -- (0,5.1);
% Entrée / sortie d'eau
\draw[very thick, popblue, -{Stealth}] (-1.7,2.3) -- (-0.45,2.3) node[midway, above, popblue]{\footnotesize entrée d'eau};
\draw[very thick, popblue, -{Stealth}] (0.45,4.9) -- (1.7,4.9) node[midway, above, popblue]{\footnotesize sortie d'eau};
% Chauffe-ballon
\draw[very thick, popdark] (-1.35,-1.35) -- (-1.35,-1.7) -- (1.35,-1.7) -- (1.35,-1.35);
\draw[popgold, thick, decorate, decoration={coil, segment length=8pt, amplitude=3pt}] (-1.1,-1.52) -- (1.1,-1.52);
\node[popdark] at (0,-2.05) {\footnotesize chauffe-ballon};
% Vapeurs
\draw[-{Stealth}, popmuted, dashed] (0.15,0.6) -- (0.15,1.7);
\draw[-{Stealth}, popmuted, dashed] (-0.15,2.0) -- (-0.15,0.9);
\node[popmuted, right] at (0.55,1.35) {\footnotesize vapeurs};
\end{tikzpicture}
\legende{Montage de chauffage à reflux : les vapeurs se condensent dans le réfrigérant à eau et retombent dans le ballon ; on chauffe sans perdre de matière.}
\end{popfigure}

\begin{experience}[Synthèse de l'éthanoate d'éthyle au laboratoire]
\textbf{Protocole} : dans un ballon, introduire environ $\SI{0,2}{mol}$ d'acide éthanoïque, $\SI{0,2}{mol}$ d'éthanol, quelques gouttes d'acide sulfurique concentré et quelques grains de pierre ponce. Adapter le réfrigérant à eau et chauffer à reflux pendant environ 30 minutes.\par
\textbf{Observations} : après refroidissement et versement du contenu du ballon dans une solution de carbonate de sodium, on observe un liquide huileux surnageant, dégageant une \textbf{odeur fruitée} caractéristique : c'est l'ester formé.\par
\textbf{Interprétation} : l'odeur fruitée prouve la formation d'une espèce nouvelle, l'éthanoate d'éthyle. L'acide sulfurique a joué le rôle de catalyseur ; le reflux a permis de chauffer sans perte de matière.
\end{experience}

\courssub{Nomenclature des esters}

\begin{methode}[Nommer un ester $\mathrm{R-COO-R'}$]
Le nom d'un ester comporte deux mots, construits ainsi :
\begin{enumerate}
\item \textbf{Premier mot} : nom de l'acide carboxylique dont dérive l'ester, dans lequel la terminaison « -oïque » est remplacée par « \textbf{-oate} ». Exemple : acide éthanoïque $\rightarrow$ éthano\textbf{ate}.
\item \textbf{Deuxième mot} : « \textbf{de} » suivi du nom du groupe \textbf{alkyle} $\mathrm{R'}$ provenant de l'alcool (méthyle, éthyle, propyle\ldots). Exemple : éthanol $\rightarrow$ d'\textbf{éthyle}.
\end{enumerate}
Ainsi, $\ce{CH3COOCH2CH3}$ se nomme \textbf{éthanoate d'éthyle}.
\end{methode}

\begin{exemplebox}[Quelques esters et leurs arômes]
\begin{center}
\begin{tabularx}{\linewidth}{|l|X|X|}
\hline
\rowcolor{popblueL} \textbf{Formule semi-développée} & \textbf{Nom} & \textbf{Arôme} \\
\hline
$\ce{CH3COO-CH2-CH2-CH(CH3)2}$ & éthanoate d'isoamyle (ou de 3-méthylbutyle) & banane \\
\hline
$\ce{CH3COO-CH2-CH3}$ & éthanoate d'éthyle & pomme, solvant \\
\hline
$\ce{HCOO-CH2-CH3}$ & méthanoate d'éthyle & framboise \\
\hline
$\ce{CH3-(CH2)2-COO-CH2-CH3}$ & butanoate d'éthyle & ananas \\
\hline
\end{tabularx}
\end{center}
\end{exemplebox}

\begin{savaistu}[Les esters, parfumeurs de la nature]
De nombreux arômes naturels de fruits sont des esters : l'éthanoate d'isoamyle donne son parfum à la banane, le butanoate d'éthyle à l'ananas. L'industrie agroalimentaire synthétise ces esters pour aromatiser jus de fruits, bonbons et yaourts — y compris certaines boissons très populaires au Cameroun. Les esters jouent aussi un rôle capital ailleurs : les \textbf{graisses et huiles} (comme l'huile de palme) sont des triesters du glycérol, et leur hydrolyse en présence de soude produit le \textbf{savon} : c'est la saponification, à la base des savonneries artisanales.
\end{savaistu}

\courssub{La réaction inverse : l'hydrolyse des esters}

L'estérification étant un équilibre, sa réaction inverse existe : c'est l'\textbf{hydrolyse}.

\begin{definition}[Hydrolyse d'un ester]
L'\cle{hydrolyse} d'un ester est la réaction de l'ester avec l'eau, qui régénère l'acide carboxylique et l'alcool :
\[ \mathrm{R{-}COO{-}R'} + \ce{H2O} \rightleftharpoons \mathrm{R{-}COOH} + \mathrm{R'{-}OH} \]
Comme l'estérification, elle est \textbf{lente}, \textbf{limitée} et \textbf{athermique} : c'est exactement la même réaction, lue dans l'autre sens. Un excès d'eau déplace l'équilibre vers l'hydrolyse.
\end{definition}

\begin{exoresolu}[Rendement d'une estérification]
\textbf{Énoncé.} On chauffe à reflux, en présence d'acide sulfurique, un mélange contenant $n_0 = \SI{0,50}{mol}$ d'acide éthanoïque et $n_0 = \SI{0,50}{mol}$ d'éthanol. À l'équilibre, le dosage de l'acide restant montre qu'il en reste $\SI{0,165}{mol}$.
\begin{enumerate}
\item Écrire l'équation-bilan de la réaction.
\item Dresser le tableau d'avancement et déterminer l'avancement final $x_f$.
\item Calculer le rendement de l'estérification.
\item Proposer deux moyens d'augmenter ce rendement.
\end{enumerate}
\tcblower
\textbf{Solution détaillée.}
\begin{enumerate}
\item Équation-bilan (double flèche, car réaction limitée) :
\[ \ce{CH3COOH + C2H5OH <=> CH3COOC2H5 + H2O} \]
\item Tableau d'avancement (quantités en mol) :
\begin{center}
\begin{tabularx}{\linewidth}{|l|X|X|X|X|}
\hline
\rowcolor{popblueL} \textbf{État} & $\ce{CH3COOH}$ & $\ce{C2H5OH}$ & $\ce{CH3COOC2H5}$ & $\ce{H2O}$ \\
\hline
Initial & $0{,}50$ & $0{,}50$ & $0$ & $0$ \\
\hline
En cours & $0{,}50 - x$ & $0{,}50 - x$ & $x$ & $x$ \\
\hline
Final & $0{,}50 - x_f$ & $0{,}50 - x_f$ & $x_f$ & $x_f$ \\
\hline
\end{tabularx}
\end{center}
L'acide restant vaut $0{,}50 - x_f = 0{,}165$, d'où :
\[ x_f = 0{,}50 - 0{,}165 = \SI{0,335}{mol} \]
Il s'est donc formé $\SI{0,335}{mol}$ d'éthanoate d'éthyle.
\item Si la réaction était totale, on obtiendrait $x_{\max} = \SI{0,50}{mol}$ d'ester. Le rendement est :
\[ \eta = \frac{x_f}{x_{\max}} = \frac{0{,}335}{0{,}50} = 0{,}67 \quad \text{soit} \quad \eta \approx 67\ \% \]
On retrouve la valeur caractéristique d'un mélange équimolaire avec un alcool primaire.
\item Pour améliorer le rendement, on peut : (a) introduire l'éthanol \textbf{en excès} ; (b) \textbf{éliminer l'eau} (ou l'ester) par distillation au fur et à mesure de sa formation. En revanche, ajouter davantage d'acide sulfurique ou chauffer plus fort ne changerait pas le rendement : cela ne ferait qu'atteindre plus vite le même équilibre.
\end{enumerate}
\end{exoresolu}

\begin{aretenir}[L'essentiel sur l'estérification]
\begin{itemize}
\item L'estérification est la réaction d'un \textbf{acide carboxylique} avec un \textbf{alcool}, qui donne un \textbf{ester} et de l'\textbf{eau} : $R\text{-}COOH + R'\text{-}OH \rightleftharpoons R\text{-}COO\text{-}R' + \ce{H2O}$.
\item Elle est \textbf{lente}, \textbf{limitée} (équilibre, double flèche $\ce{<=>}$) et \textbf{athermique}.
\item Le rendement limite est d'environ \textbf{67 \%} pour un mélange équimolaire avec un alcool primaire.
\item Le \textbf{chauffage} et les ions \ce{H+} (acide sulfurique) \textbf{accélèrent} la réaction sans déplacer l'équilibre ; l'\textbf{excès d'un réactif} ou l'\textbf{élimination d'un produit} améliorent le rendement.
\item Nomenclature : \textbf{alkanoate d'alkyle} (ex. éthanoate d'éthyle).
\item La réaction inverse est l'\textbf{hydrolyse} de l'ester.
\end{itemize}
\end{aretenir}

\begin{exercice}[À toi de jouer]
On fait réagir de l'acide méthanoïque \ce{HCOOH} avec du propan-1-ol \ce{CH3-CH2-CH2-OH}.
\begin{enumerate}
\item Écrire l'équation-bilan de la réaction d'estérification.
\item Nommer l'ester obtenu.
\item Citer les trois caractéristiques de cette réaction.
\item On part de $\SI{0,30}{mol}$ de chaque réactif et on obtient à l'équilibre $\SI{0,20}{mol}$ d'ester. Calculer le rendement.
\end{enumerate}
\end{exercice}

\begin{corrige}[Exercice : estérification de l'acide méthanoïque]
\begin{enumerate}
\item $\ce{HCOOH + CH3-CH2-CH2-OH <=> HCOO-CH2-CH2-CH3 + H2O}$
\item L'ester formé est le \textbf{méthanoate de propyle} (l'acide méthanoïque donne « méthanoate », le propan-1-ol donne « propyle »).
\item La réaction est \textbf{lente}, \textbf{limitée} (équilibre chimique) et \textbf{athermique}.
\item Rendement : $\eta = \dfrac{0{,}20}{0{,}30} \approx 0{,}67$, soit environ $67\ \%$, valeur typique d'un mélange équimolaire avec un alcool primaire.
\end{enumerate}
\end{corrige}

\coursec{Préparation des alcools}

Après avoir étudié les réactions caractéristiques des alcools — acido-basicité avec le sodium, déshydratation, oxydation et estérification — nous nous penchons maintenant sur \emph{comment on les obtient} à l'échelle industrielle et artisanale. Deux voies majeures s'offrent à nous : l'hydratation d'un alcène (voie pétrochimique) et la fermentation alcoolique (voie biologique), cette dernière étant au cœur de nos traditions et de notre économie locale (bili-bili, vin de palme, éthanol carburant).

\courssub{Hydratation d'un alcène : addition d'eau sur une double liaison}

\textbf{Principe général.}\par
L'hydratation d'un alcène consiste en l'addition d'une molécule d'eau (\ce{H2O}) sur la double liaison carbone-carbone (\ce{C=C}) en milieu acide. C'est une réaction d'\textbf{addition électrophile} catalysée par un acide fort (acide sulfurique \ce{H2SO4} concentré, acide phosphorique \ce{H3PO4} sur support solide) ou par un catalyseur hétérogène (zéolites, phosphore de tungstène) en phase gazeuse à haute température.

\begin{definition}[Hydratation d'un alcène]
L'hydratation d'un alcène est une réaction d'addition d'eau sur une double liaison \ce{C=C} en présence d'un catalyseur acide, conduisant à un alcool. C'est l'opération inverse de la déshydratation intramoléculaire étudiée précédemment.
\end{definition}

\textbf{Équation générale et conditions.}\par
Pour l'éthylène (alcène le plus simple), la réaction s'écrit :
\[ \ce{CH2=CH2 + H2O ->[H3PO4, 300 ^{\circ}C, pression] CH3-CH2OH} \]
\emph{Conditions industrielles :} phase gazeuse, catalyseur \ce{H3PO4} sur silice/alumine, \SI{300}{\celsius}, \SI{70}{bar}. Rendement par passage \approx 5 \%, mais recyclage de l'éthylène non réagi.

Sous forme générale pour un alcène linéaire $C_nH_{2n}$ :
$$ C_nH_{2n} + H_2O \xrightarrow{\text{catalyseur acide}} C_nH_{2n+1}OH $$

\begin{popfigure}
\centering
\[ \ce{CH2=CH-CH3 + H2O ->[H2SO4][\Delta] CH3-CH(OH)-CH3} \]
\legende{Hydratation du propène : l'alcool majeur est le propan-2-ol (règle de Markovnikov), le propan-1-ol (\ce{CH3CH2CH2OH}) est minoritaire.}
\end{popfigure}

\textbf{Cas des alcènes asymétriques : la règle de Markovnikov.}\par
Lorsque l'alcène n'est pas symétrique (ex. : propène \ce{CH2=CH-CH3}, but-1-ène, but-2-ène), deux alcools isomères \emph{peuvent} se former. L'expérience montre qu'un seul prédomine largement : c'est l'alcool le plus substitué (le groupe \ce{-OH} se fixe sur le carbone de l'ancienne double liaison qui portait le \emph{moins} d'hydrogènes).

\begin{methode}[Prévoir le produit majoritaire de l'hydratation d'un alcène asymétrique (Règle de Markovnikov)]
\begin{enumerate}
\item Identifier les deux atomes de carbone de la double liaison \ce{C=C}.
\item Compter le nombre d'atomes d'hydrogène portés par \textbf{chacun} de ces deux carbones.
\item L'atome d'hydrogène (\ce{H}) de la molécule d'eau s'ajoute sur le carbone \textbf{le plus hydrogéné} (celui qui a déjà le plus d'H).
\item Le groupe hydroxyle (\ce{-OH}) s'ajoute donc sur l'autre carbone (le \textbf{moins hydrogéné}, le plus substitué).
\item Écrire la formule semi-développée de l'alcool formé.
\end{enumerate}
\end{methode}

\begin{exemplebox}[Hydratation du propène]
\emph{Étape 1 :} Propène = \ce{CH2=CH-CH3}. Carbone 1 (\ce{CH2}) porte 2 H. Carbone 2 (\ce{CH}) porte 1 H.
\emph{Étape 2 :} Le carbone 1 est le plus hydrogéné. Le \ce{H} de l'eau s'y fixe.
\emph{Étape 3 :} Le \ce{-OH} se fixe sur le carbone 2.
\emph{Produit majoritaire :} \ce{CH3-CH(OH)-CH3} = \textbf{propan-2-ol} (alcool secondaire).
\emph{Produit minoritaire :} \ce{CH3-CH2-CH2OH} = propan-1-ol (alcool primaire).
\end{exemplebox}

\begin{attention}[Piège classique : inversion de Markovnikov]
Ne confondez pas « carbone le plus hydrogéné » et « carbone le plus substitué ».
\begin{itemize}
\item \ce{H} (de \ce{H2O}) $\to$ Carbone \textbf{plus hydrogéné} (moins substitué).
\item \ce{OH} (de \ce{H2O}) $\to$ Carbone \textbf{moins hydrogéné} (plus substitué).
\end{itemize}
L'alcool formé majoritairement est toujours l'alcool \textbf{le plus substitué} (tertiaire $>$ secondaire $>$ primaire). C'est la conséquence de la stabilité des carbocations intermédiaires (\ce{C+} tertiaire plus stable que secondaire, etc.).
\end{attention}

\begin{exoresolu}[Rendement de l'hydratation du propène]
\emph{Énoncé.} On hydrate \SI{42,0}{g} de propène (\ce{C3H6}, $M = \SI{42,0}{g.mol^{-1}}$) en présence d'acide sulfurique concentré. On récupère \SI{36,0}{g} de propan-2-ol (\ce{C3H8O}, $M = \SI{60,0}{g.mol^{-1}}$). Calculer le rendement de la réaction.
\tcblower
\emph{Solution.}
\begin{enumerate}
\item \textbf{Équation bilan :} $\ce{CH2=CH-CH3 + H2O -> CH3-CH(OH)-CH3}$
\item \textbf{Avancement maximal :} $n_{\text{propène, initial}} = \frac{42,0}{42,0} = \SI{1,00}{mol}$. L'eau est en excès. $x_{\max} = \SI{1,00}{mol}$.
\item \textbf{Masse théorique de propan-2-ol :} $m_{\text{th}} = x_{\max} \times M_{\text{propan-2-ol}} = 1,00 \times 60,0 = \SI{60,0}{g}$.
\item \textbf{Rendement :} $\eta = \frac{m_{\text{obtenu}}}{m_{\text{th}}} \times 100 = \frac{36,0}{60,0} \times 100 = \mathbf{60,0\,\%}$.
\end{enumerate}
\emph{Commentaire :} Le rendement n'est pas quantitatif car l'équilibre est limité et des réactions secondaires (polymérisation, formation d'éthers) se produisent.
\textbf{Réponse :} Le rendement de l'hydratation est de \textbf{60,0 \%}.
\end{exoresolu}

\courssub{Fermentation alcoolique : la voie biologique}

\textbf{Contexte et définition.}\par
Depuis des millénaires, l'homme utilise la capacité de micro-organismes (levures) à transformer les sucres en éthanol. Au Cameroun, cette tradition est vivante : \emph{bili-bili} (mil, maïs, sorgho), \emph{vin de palme} (sève de palmier à huile ou raphia), \emph{ngoñ} (banane plantain), jus de bissap fermenté. C'est la \textbf{fermentation alcoolique}.

\begin{definition}[Fermentation alcoolique]
La fermentation alcoolique est une dégradation \textbf{anaérobie} (sans oxygène) du glucose (\ce{C6H12O6}) en éthanol (\ce{C2H5OH}) et dioxyde de carbone (\ce{CO2}), catalysée par des enzymes (zymase) produites par des levures (principalement \emph{Saccharomyces cerevisiae}).
\end{definition}

\textbf{Équation-bilan et conditions opératoires.}\par
L'équation chimique globale, établie par Gay-Lussac puis précisée par Pasteur, est :
\[ \ce{C6H12O6 ->[Levures, 30\textdegree C, anaerobiose] 2 C2H5OH + 2 CO2 ^} \]

\begin{aretenir}[Conditions optimales de la fermentation alcoolique]
\begin{itemize}
\item \textbf{Anaérobiose stricte :} absence d'\ce{O2} (présence d'\ce{O2} $\to$ respiration cellulaire $\to$ \ce{CO2} + \ce{H2O}, pas d'alcool).
\item \textbf{Température :} \SIrange{25}{35}{\celsius} (optimum \approx \SI{30}{\celsius}). Au-delà de \SI{40}{\celsius}, les enzymes se dénaturent ; en dessous de \SI{15}{\celsius}, la réaction est trop lente.
\item \textbf{pH :} légèrement acide (\SIrange{4}{5}{}), ce qui inhibe les bactéries concurrentes (acide lactique, acétique).
\item \textbf{Substrat :} sucres fermentescibles (glucose, fructose, saccharose, maltose). L'amidon doit être hydrolysé au préalable (maltage ou amylases).
\item \textbf{Levures :} \emph{Saccharomyces cerevisiae} (levure de boulanger/brasseur) ou levures sauvages présentes sur les fruits/écorces.
\end{itemize}
\end{aretenir}

\begin{attention}[Erreur fréquente : l'équation mal équilibrée]
Vérifiez systématiquement le bilan d'atomes :
\begin{center}
\begin{tabular}{c|c|c|c}
 & \ce{C} & \ce{H} & \ce{O} \\ \hline
Glucose (\ce{C6H12O6}) & 6 & 12 & 6 \\
2 Éthanol (2 \ce{C2H5OH}) & 4 & 12 & 2 \\
2 \ce{CO2} & 2 & 0 & 4 \\ \hline
\textbf{Total Produits} & \textbf{6} & \textbf{12} & \textbf{6} \\
\end{tabular}
\end{center}
\textbf{N'oubliez jamais le \ce{CO2} gazeux !} C'est lui qui fait lever le pain, pétiller la bière et le bili-bili, et qui permet de suivre la réaction (perte de masse, bulles).
\end{attention}

\textbf{Limite intrinsèque : la toxicité de l'éthanol.}\par
Les levures produisent de l'éthanol, mais celui-ci devient toxique pour elles au-delà d'une certaine concentration. La fermentation s'arrête spontanément vers \SIrange{12}{15}{\%} vol. (degré alcoolique). Pour obtenir des spiritueux plus forts (whisky, vodka, alcool médical \SI{96}{\%}), il faut \textbf{distiller} le moût fermenté.

\begin{exemplebox}[Calcul du rendement massique théorique : du glucose à l'éthanol]
\emph{Données :} $M_{\text{glucose}} = \SI{180}{g.mol^{-1}}$ ; $M_{\text{éthanol}} = \SI{46}{g.mol^{-1}}$ ; $\rho_{\text{éthanol}} = \SI{0,79}{g.mL^{-1}}$.
\emph{Question :} Quelle masse et quel volume d'éthanol pur peut-on obtenir théoriquement à partir de \SI{180}{g} de glucose (1 mol) ?
\emph{Résolution :}
\begin{align*}
n_{\text{glucose}} &= \frac{180}{180} = \SI{1,00}{mol} \\
n_{\text{éthanol, th}} &= 2 \times n_{\text{glucose}} = \SI{2,00}{mol} \quad (\text{selon équation bilan}) \\
m_{\text{éthanol, th}} &= 2,00 \times 46 = \mathbf{\SI{92}{g}} \\
V_{\text{éthanol, th}} &= \frac{m}{\rho} = \frac{92}{0,79} \approx \mathbf{\SI{116,5}{mL}} \approx \SI{117}{mL}
\end{align*}
En pratique, le rendement est inférieur (biomasse levures, glycérol, acide acétique, pertes \ce{CO2}, distillation incomplète). Un bon rendement industriel tourne autour de \SI{90}{--}\SI{95}{\%} du théorique.
\end{exemplebox}

\courssub{Applications locales et filière éthanol au Cameroun}

Le Cameroun possède des atouts agronomiques majeurs (canne à sucre, maïs, manioc, patate douce) pour développer une filière bioéthanol. La \textbf{SOSUCAM} (Société Sucrière du Cameroun) produit de l'éthanol à partir de la mélasse (résidu de la fabrication du sucre de canne) dans ses unités de Nkoteng et Mbandjock.

\begin{popfigure}
\centering
\begin{tikzpicture}[
    node distance=1.5cm and 2.5cm,
    box/.style={draw=popblue, thick, rounded corners, fill=popblueL, minimum width=2.8cm, minimum height=1.2cm, align=center, font=\small},
    arrow/.style={-Stealth, thick, popdark},
    labelbox/.style={font=\footnotesize, text=popmuted, align=center}
]
% Nœuds principaux
\node[box] (canne) {Canne à sucre\\ (Culture)};
\node[box, right=of canne] (broyage) {Broyage / Diffusion\\ (Extraction du jus)};
\node[box, right=of broyage] (sucre) {Cristallisation\\ (Sucre + Mélasse)};
\node[box, below=of sucre] (melasse) {Mélasse\\ (50\% sucres)};
\node[box, right=of melasse] (ferment) {Fermentation\\ (Levures, 30°C, anaérobiose)};
\node[box, right=of ferment] (vin) {Vin de mélasse\\ (8--10\% vol. éthanol)};
\node[box, right=of vin] (distil) {Distillation\\ (Colonne à plateaux)};
\node[box, right=of distil] (ethanol) {Éthanol\\ (96\% vol. puis 99,5\%)};
\node[box, below=of ethanol, xshift=-1.5cm] (usages) {Usages :\\ Solvants, Désinfectants,\\ Biocarburant (E10),\\ Boissons spiritueuses};

% Flèches principales
\draw[arrow] (canne) -- (broyage) node[midway, above, labelbox] {Récolte};
\draw[arrow] (broyage) -- (sucre) node[midway, above, labelbox] {Purification};
\draw[arrow] (sucre) -- node[left, labelbox] {Sous-produit} (melasse);
\draw[arrow] (melasse) -- (ferment) node[midway, above, labelbox] {Dilution + Nutriments};
\draw[arrow] (ferment) -- (vin) node[midway, above, labelbox] {\ce{CO2} dégagé};
\draw[arrow] (vin) -- (distil) node[midway, above, labelbox] {Chauffage};
\draw[arrow] (distil) -- (ethanol) node[midway, above, labelbox] {Concentration};
\draw[arrow] (ethanol) -- (usages);

% Boucle retour vinasse (optionnel, simplifié)
\node[labelbox, below=0.5cm of ferment] (vinasse) {Vinasses $\to$ Engrais / Biogaz};
\draw[arrow, popgreen] (ferment) |- (vinasse);
\end{tikzpicture}
\legende{Filière locale de production d'éthanol à partir de la canne à sucre (SOSUCAM) : de la culture à la distillation.}
\end{popfigure}

\begin{savaistu}[L'éthanol camerounais : au-delà de la boisson]
\begin{itemize}
\item \textbf{Carburant :} Le Cameroun expérimente l'incorporation d'éthanol dans l'essence (objectif E10 : 10\% éthanol) pour réduire l'importation de pétrole et les émissions de \ce{CO2}.
\item \textbf{Secteur médical :} L'éthanol \SI{96}{\%} (alcool à 96°) est l'antiseptique de référence dans nos formations sanitaires (désinfection cutanée, matériel). La production locale (SOSUCAM, UCAM) sécurise l'approvisionnement.
\item \textbf{Chimie verte :} L'éthanol est une plateforme chimique : déshydratation $\to$ éthylène (plastiques), oxydation $\to$ acide acétique, estérification $\to$ solvants (acétate d'éthyle).
\item \textbf{Valorisation des déchets :} Les vinasses (résidus de distillation) sont riches en potassium et matière organique $\to$ engrais organiques (retour au champ) ou méthanisation (biogaz pour l'énergie de l'usine).
\end{itemize}
\end{savaistu}

\courssub{Bilan de la leçon : Carte mentale des transformations des alcools}

Pour réviser efficacement, construis cette carte mentale (à reproduire sur ta fiche de révision). Elle relie la structure de l'alcool (primaire, secondaire, tertiaire) à ses réactions types.

\begin{popfigure}
\centering
\begin{tikzpicture}[
    mindmap,
    concept color=popblue,
    level 1 concept/.append style={level distance=4.5cm, sibling angle=90, font=\small\bfseries},
    level 2 concept/.append style={level distance=3.2cm, sibling angle=45, font=\footnotesize},
    level 3 concept/.append style={level distance=2.8cm, font=\scriptsize},
    every node/.style={scale=0.85},
    text=white
]
\node[concept] {ALCOOLS\\ R-\ce{OH}}
    child[concept color=poporange] { node[concept] {Préparation}
        child { node[concept] {Hydratation alcènes\\ (Markovnikov)} }
        child { node[concept] {Fermentation\\ Glucose $\to$ Éthanol} }
    }
    child[concept color=popgreen] { node[concept] {Réactions du \ce{-OH}}
        child { node[concept] {Acidité\\ + Na $\to$ Alcoolate} }
        child { node[concept] {Estérification\\ + Acide $\to$ Ester} }
        child { node[concept] {Déshydratation\\ $\to$ Alcène / Éther} }
    }
    child[concept color=poppurple] { node[concept] {Oxydation (Classe)}
        child[concept color=poppink] { node[concept] {Primaire\\ $\to$ Aldéhyde $\to$ Acide} }
        child[concept color=poppink] { node[concept] {Secondaire\\ $\to$ Cétone} }
        child[concept color=poppink] { node[concept] {Tertiaire\\ Pas d'oxydation ménagée} }
    }
    child[concept color=popgold] { node[concept] {Combustion\\ $\ce{CO2} + \ce{H2O}$} };
\end{tikzpicture}
\legende{Carte mentale récapitulative : préparation, réactions du groupe hydroxyle et oxydation selon la classe de l'alcool.}
\end{popfigure}

\begin{exercice}[Entraînement Bac : Préparation et réactions croisées]
On part de \SI{56,0}{g} de but-1-ène (\ce{C4H8}, $M = \SI{56,0}{g.mol^{-1}}$).
\begin{enumerate}
\item Écrire l'équation de l'hydratation du but-1-ène en précisant le produit \textbf{majoritaire} (règle de Markovnikov). Nommer cet alcool.
\item Cet alcool majoritaire est oxydé par le permanganate de potassium en milieu acide chauffé. Écrire l'équation de cette oxydation ménagée. Nommer le produit organique formé.
\item Calculer la masse théorique de ce produit final obtenue à partir des \SI{56,0}{g} de but-1-ène initial (rendement 100\% à chaque étape).
\end{enumerate}
\begin{corrige}[Exercice 1]
\begin{enumerate}
\item \textbf{Hydratation :} $\ce{CH2=CH-CH2-CH3 + H2O ->[H^+] CH3-CH(OH)-CH2-CH3}$.
Produit majoritaire : \textbf{butan-2-ol} (alcool secondaire). Le \ce{OH} se porte sur le C2 (moins hydrogéné de la double liaison).
\item \textbf{Oxydation ménagée (alcool secondaire) :} $\ce{CH3-CH(OH)-CH2-CH3 + [O] -> CH3-CO-CH2-CH3 + H2O}$.
Avec \ce{KMnO4}/\ce{H+} : $\ce{5 CH3CH(OH)CH2CH3 + 2 MnO4^- + 6 H+ -> 5 CH3COCH2CH3 + 2 Mn^{2+} + 8 H2O}$.
Produit : \textbf{butan-2-one} (méthyléthylcétone).
\item \textbf{Calcul :}
$n_{\text{but-1-ène}} = 56,0 / 56,0 = \SI{1,00}{mol}$.
1 mol but-1-ène $\to$ 1 mol butan-2-ol $\to$ 1 mol butan-2-one.
$M_{\text{butan-2-one}} = \SI{72,0}{g.mol^{-1}}$.
$m_{\text{th}} = 1,00 \times 72,0 = \mathbf{\SI{72,0}{g}}$.
\end{enumerate}
\end{corrige}
\end{exercice}

\paragraph{Mots-clés à retenir pour le Bac :} \cle{Hydratation}, \cle{Markovnikov}, \cle{Fermentation alcoolique}, \cle{Anaérobiose}, \cle{Levures (Saccharomyces)}, \cle{Distillation}, \cle{Rendement}, \cle{SOSUCAM}, \cle{Bilan C/H/O}.

\newpage

\coursec{Travaux pratiques}

\courssub{Objectifs du TP}

À l'issue de cette séance de travaux pratiques, tu dois être capable de :

\begin{itemize}
\item réaliser l'\cle{oxydation catalytique} des vapeurs d'éthanol sur cuivre incandescent (expérience de la \cle{lampe sans flamme}) et identifier le produit formé ;
\item distinguer expérimentalement les trois \cle{classes d'alcools} (primaire, secondaire, tertiaire) par leur comportement vis-à-vis du \cle{permanganate de potassium} acidifié ;
\item mettre en œuvre les tests d'identification des \cle{aldéhydes} et des \cle{cétones} (liqueur de Fehling, DNPH) ;
\item écrire les équations-bilans des réactions d'oxydation ménagée observées ;
\item consigner des observations de manière rigoureuse dans un tableau d'exploitation.
\end{itemize}

\courssub{Matériel et produits}

\begin{center}
\rowcolors{2}{popblueL}{white}
\begin{tabularx}{\linewidth}{|l|X|}
\hline
\rowcolor{popblue!60} \textbf{Matériel} & \textbf{Produits chimiques} \\
\hline
Tubes à essai (au moins 5) avec portoir & Éthanol à \SI{95}{\celsius} (ou alcool médical) \\
Bécher de \SI{50}{mL} & Propan-2-ol (alcool isopropylique) \\
Fil de cuivre (environ \SI{15}{cm}) enroulé en spirale & 2-méthylpropan-2-ol (butanol tertiaire) \\
Bec Bunsen ou lampe à alcool & Solution de \ce{KMnO4} à \SI{0,01}{mol.L^{-1}} environ \\
Pince en bois ou en fer & Acide sulfurique dilué (\SI{1}{mol.L^{-1}}) \\
Pipettes compte-gouttes & Liqueur de Fehling (fraîchement préparée) \\
Papier filtre, baguettes de verre & Solution de DNPH (2,4-dinitrophénylhydrazine) \\
Spatule, entonnoir à solide & Eau distillée \\
\hline
\end{tabularx}
\end{center}

\begin{savaistu}[Alternatives avec le matériel local]
Si le propan-2-ol ou le 2-méthylpropan-2-ol manque au laboratoire, on peut les remplacer respectivement par de l'alcool isopropylique vendu en pharmacie (désinfectant) et démontrer simplement la différence primaire/tertiaire avec l'éthanol et un alcool tertiaire fourni par le professeur. À défaut de DNPH, la liqueur de Fehling seule permet déjà de distinguer aldéhyde (précipité rouge brique) et cétone (pas de réaction). La liqueur de Fehling peut être préparée sur place en mélangeant une solution de sulfate de cuivre(II) et une solution de tartrate de sodium et potassium en milieu basique.
\end{savaistu}

\courssub{Sécurité}

\begin{attention}[Consignes de sécurité impératives]
\begin{itemize}
\item \textbf{Éthanol et alcools} : liquides \emph{inflammables} (pictogramme flamme). Tenir les flacons éloignés de toute flamme ; ne jamais chauffer directement un flacon d'alcool.
\item \textbf{Acide sulfurique} : \emph{corrosif} (pictogramme corrosion). Porter blouse, lunettes et gants ; verser toujours l'acide dans l'eau, jamais l'inverse.
\item \textbf{Permanganate de potassium} : \emph{oxydant} et tachant (pictogramme flamme sur cercle). Il tache durablement la peau et les vêtements ; manipuler avec soin.
\item \textbf{Liqueur de Fehling} : basique et contenant du cuivre ; \emph{nocive} (pictogramme point d'exclamation). Éviter tout contact avec la peau.
\item \textbf{Éthanal dégagé} : vapeurs irritantes ; travailler près d'une aération et ne pas inhaler directement au-dessus du bécher.
\item Le fil de cuivre rougi reste brûlant plusieurs minutes : le saisir uniquement avec la pince.
\item En cas de projection, rincer abondamment à l'eau et prévenir immédiatement l'enseignant.
\end{itemize}
\end{attention}

\courssub{Schéma du dispositif}

\begin{popfigure}
\begin{center}
\begin{tikzpicture}[scale=0.95, every node/.style={font=\small}]
% Bécher (paroi unique)
\draw[thick, popblue] (0,0) -- (0,3) -- (2.4,3) -- (2.4,0);
\draw[thick, popblue] (0,0) -- (2.4,0); % fond
% Liquide éthanol
\fill[poporange!40] (0.06,0.06) rectangle (2.34,1.5);
\draw[poporange] (0.06,1.5) -- (2.34,1.5);
\node[popdark] at (1.2,0.7) {éthanol tiède};
% Spirale de cuivre
\draw[very thick, poppink] (1.2,4.6) -- (1.2,3.6);
\draw[very thick, poppink, decorate, decoration={coil, aspect=0.4, segment length=4pt, amplitude=5pt}] (1.2,3.6) -- (1.2,2.1);
\node[right, poppink, align=left] at (1.35,4.2) {spirale de cuivre\\ rougie à la flamme};
% Vapeurs
\draw[->, popmuted, dashed] (0.5,1.6) .. controls (0.3,2.3) .. (0.6,2.9);
\draw[->, popmuted, dashed] (1.9,1.6) .. controls (2.1,2.3) .. (1.8,2.9);
\node[left, popmuted, align=right] at (-0.3,2.4) {vapeurs\\ d'éthanol};
% Chauffage doux
\draw[poporange, thick] (1.0,-0.55) .. controls (1.1,-0.15) and (1.2,-0.35) .. (1.2,-0.05);
\draw[poporange, thick] (1.4,-0.55) .. controls (1.3,-0.15) and (1.2,-0.35) .. (1.2,-0.05);
\node[below, poporange] at (1.2,-0.6) {chauffage doux};
% Papier DNPH
\draw[thick, popgreen] (2.9,2.6) rectangle (4.1,3.0);
\node[right, popgreen, align=left] at (4.2,2.8) {papier imbibé de DNPH\\ (test des vapeurs)};
\end{tikzpicture}
\end{center}
\legende{Dispositif de la lampe sans flamme : la spirale de cuivre incandescente, plongée dans les vapeurs d'éthanol, continue de rougeoyer sans flamme grâce à la chaleur dégagée par l'oxydation catalytique.}
\end{popfigure}

\courssub{Protocole}

\textbf{Partie 1 : la lampe sans flamme (oxydation catalytique de l'éthanol)}

\begin{enumerate}
\item Verser environ \SI{5}{mL} d'éthanol dans un bécher propre et sec.
\item Chauffer \emph{doucement} le bécher quelques secondes (sans porter à ébullition) afin de dégager des vapeurs d'éthanol, puis éloigner la flamme.
\item À l'aide de la pince, rougir la spirale de fil de cuivre dans la flamme du bec Bunsen jusqu'à ce qu'elle soit incandescente.
\item Plonger immédiatement la spirale rougie au-dessus du liquide, dans les vapeurs du bécher. Observer : le fil \emph{continue de rougeoyer} alors qu'il n'y a aucune flamme.
\item Renouveler l'opération 2 à 3 fois. Sentir prudemment (en éventant avec la main) : une odeur piquante de pomme verte se dégage.
\item Approcher à l'orifice du bécher un papier filtre imbibé de DNPH : observer la formation d'un précipité jaune-orangé.
\end{enumerate}

\textbf{Partie 2 : oxydation ménagée par le permanganate selon la classe de l'alcool}

\begin{enumerate}
\item Numéroter trois tubes à essai T1, T2, T3.
\item Introduire dans chaque tube \SI{2}{mL} de solution violette de \ce{KMnO4}, puis 2 gouttes d'acide sulfurique dilué (agiter).
\item Ajouter goutte à goutte, en agitant après chaque goutte :
\begin{itemize}
\item dans T1 : 5 gouttes d'\textbf{éthanol} (alcool primaire) ;
\item dans T2 : 5 gouttes de \textbf{propan-2-ol} (alcool secondaire) ;
\item dans T3 : 5 gouttes de \textbf{2-méthylpropan-2-ol} (alcool tertiaire).
\end{itemize}
\item Observer l'évolution de la coloration dans chaque tube pendant 5 minutes.
\item Tester le contenu de T1 avec la liqueur de Fehling (chauffage au bain-marie) et celui de T2 avec la DNPH.
\item Consigner toutes les observations dans le tableau d'exploitation.
\end{enumerate}

\courssub{Observations et résultats attendus}

\begin{center}
\rowcolors{2}{popgreenL}{white}
\begin{tabularx}{\linewidth}{|l|X|X|X|}
\hline
\rowcolor{popgreen!60} \textbf{Tube} & \textbf{Alcool (classe)} & \textbf{Observation avec \ce{KMnO4} acidifié} & \textbf{Test complémentaire} \\
\hline
T1 & Éthanol \chemfig{CH_3-CH_2-OH} (primaire) & Décoloration progressive de la solution violette en 1 à 2 min & Fehling : précipité rouge brique de \ce{Cu2O} \\
T2 & Propan-2-ol \chemfig{H_3C-CH(-[2]OH)-CH_3} (secondaire) & Décoloration de la solution violette en 2 à 3 min & DNPH : précipité jaune-orangé ; Fehling : négatif \\
T3 & 2-méthylpropan-2-ol \chemfig{CH_3-C(-[2]CH_3)(-[6]CH_3)-OH} (tertiaire) & Aucune décoloration : la solution reste violette & Aucun test positif \\
\hline
\end{tabularx}
\end{center}

Pour la partie 1, on observe que la spirale de cuivre reste incandescente dans les vapeurs d'éthanol \emph{sans flamme}, qu'une odeur piquante caractéristique (éthanal) se dégage, et que le papier à la DNPH se couvre d'un précipité jaune-orangé.

\courssub{Exploitation}

\begin{exoresolu}[Questions guidées d'exploitation]
\textbf{Énoncé.}
\begin{enumerate}
\item Pourquoi la spirale de cuivre continue-t-elle de rougeoyer dans les vapeurs d'éthanol alors qu'il n'y a pas de flamme ? Quel est le rôle du cuivre ?
\item Écrire l'équation-bilan de l'oxydation catalytique de l'éthanol en éthanal.
\item À partir des demi-équations électroniques, établir l'équation-bilan de l'oxydation de l'éthanol par le permanganate en milieu acide (couple \ce{MnO4-}/\ce{Mn^2+}).
\item Écrire l'équation-bilan de l'oxydation ménagée du propan-2-ol et nommer le produit obtenu.
\item Pourquoi le 2-méthylpropan-2-ol ne réagit-il pas avec le permanganate de potassium acidifié ?
\item Comment les tests à la Fehling et à la DNPH permettent-ils de distinguer les produits formés dans T1 et T2 ?
\end{enumerate}
\tcblower
\textbf{Solution détaillée.}
\begin{enumerate}
\item L'oxydation de l'éthanol en éthanal est une réaction \emph{exothermique} : la chaleur dégagée entretient l'incandescence du fil de cuivre, d'où le nom de \og lampe sans flamme \fg. Le cuivre joue le rôle de \cle{catalyseur} : il accélère la réaction sans être consommé (il est régénéré en fin de réaction).
\item \[ 2\,\ce{CH3-CH2-OH} + \ce{O2} \xrightarrow{\ce{Cu}} 2\,\ce{CH3-CHO} + 2\,\ce{H2O} \]
\item Demi-équations :
\begin{align*}
\ce{MnO4- + 8H+ + 5e- &-> Mn^2+ + 4H2O} \quad (\times 2)\\
\ce{CH3-CH2-OH &-> CH3-CHO + 2H+ + 2e-} \quad (\times 5)
\end{align*}
D'où l'équation-bilan :
\[ \ce{2MnO4- + 5CH3-CH2-OH + 6H+ -> 2Mn^2+ + 5CH3-CHO + 8H2O} \]
La décoloration observée correspond à la transformation des ions permanganate violets en ions \ce{Mn^2+} quasiment incolores.
\item Le propan-2-ol, alcool secondaire, s'oxyde en \textbf{propanone} (acétone), une cétone. L'équation-bilan de l'oxydation ménagée s'écrit :
\[ \ce{CH3-CHOH-CH3 + [O] -> CH3-CO-CH3 + H2O} \]
(Si l'oxydant est le permanganate en milieu acide : $\ce{3CH3-CHOH-CH3 + 2MnO4- + 4H+ -> 3CH3-CO-CH3 + 2Mn^2+ + 5H2O}$).
\item Le 2-méthylpropan-2-ol est un alcool \emph{tertiaire} : le carbone fonctionnel ne porte \textbf{aucun atome d'hydrogène}. L'oxydation ménagée, qui consiste à retirer l'hydrogène du groupe hydroxyle \emph{et} celui du carbone fonctionnel, est donc impossible : il n'y a pas de réaction, la solution reste violette.
\item Dans T1, le produit est l'\textbf{éthanal}, un aldéhyde : il réduit la liqueur de Fehling (précipité rouge brique de \ce{Cu2O}) et donne un précipité jaune avec la DNPH. Dans T2, le produit est la \textbf{propanone}, une cétone : elle réagit avec la DNPH (précipité jaune, test commun aux composés carbonylés) mais \emph{pas} avec la liqueur de Fehling. Le test de Fehling, négatif pour les cétones, permet donc de les distinguer des aldéhydes.
\end{enumerate}
\end{exoresolu}

\courssub{Conclusion}

Ce TP a permis de mettre en évidence expérimentalement deux propriétés fondamentales des alcools. D'une part, l'expérience de la \cle{lampe sans flamme} montre que les vapeurs d'éthanol s'oxydent sur cuivre incandescent en \cle{éthanal}, avec dégagement de chaleur qui entretient la réaction : le cuivre agit comme catalyseur. D'autre part, l'action du permanganate de potassium acidifié permet de distinguer les trois classes d'alcools : un alcool \cle{primaire} s'oxyde en aldéhyde (puis en acide carboxylique si l'oxydant est en excès), un alcool \cle{secondaire} s'oxyde en cétone, tandis qu'un alcool \cle{tertiaire} ne subit pas d'oxydation ménagée. Les tests à la liqueur de Fehling et à la DNPH confirment la nature des produits formés. Ces réactions d'oxydation ménagée constituent ainsi une méthode d'identification des classes d'alcools, très utilisée au laboratoire comme dans l'industrie (contrôle de la qualité de l'éthanol produit par fermentation du jus de canne ou de palme).

\newpage

\coursec{Méthodes et exercices résolus}

\courssub{Fiche méthode 1 : Écrire l'équation-bilan de la réaction d'un alcool avec le sodium}

\begin{methode}[Alcool + sodium : la démarche en 4 étapes]
\begin{enumerate}
\item \textbf{Identifier le groupe hydroxyle} : repérer le groupe \cle{$-$OH} de l'alcool \ce{R-OH}. C'est l'hydrogène de ce groupe (et \emph{uniquement} celui-là) qui est mobile.
\item \textbf{Écrire les produits} : le sodium arrache cet hydrogène ; il se forme un \cle{alcoolate de sodium} (ou alcoolate) \ce{R-O^- Na^+} et un dégagement de \cle{dihydrogène} \ce{H2}.
\item \textbf{Équilibrer} : il faut \textbf{2} moles d'alcool et \textbf{2} moles de sodium pour 1 mole de \ce{H2} :
\[ \ce{2 R-OH + 2 Na -> 2 R-O^- Na^+ + H2 ^} \]
\item \textbf{Nommer l'alcoolate} : remplacer le suffixe « -ol » de l'alcool par « -olate de sodium ». Exemple : éthanol $\to$ éthanolate de sodium.
\end{enumerate}
\textbf{Réflexes} : vérifier la conservation des atomes et des charges ; penser au test d'identification de \ce{H2} (détonation à la flamme, « aboiement ») ; l'alcoolate est une base forte, la solution finale est basique (rougit la phénolphtaléine).
\end{methode}

\begin{exoresolu}[Réaction du sodium avec l'éthanol]
\textbf{Énoncé (type Bac).} Dans un tube à essai sec, on introduit un petit morceau de sodium dans \SI{10}{mL} d'éthanol pur. On observe un dégagement gazeux qui produit une légère détonation à l'approche d'une flamme. Après réaction, on ajoute quelques gouttes de phénolphtaléine : la solution devient rose.
\begin{enumerate}
\item Écrire l'équation-bilan de la réaction et nommer le produit organique formé.
\item Quelle propriété du groupe hydroxyle cette réaction met-elle en évidence ?
\item Sachant que la masse de sodium utilisée est $m = \SI{0,23}{g}$, calculer le volume de dihydrogène dégagé, mesuré dans les conditions où $V_m = \SI{24,0}{L.mol^{-1}}$. On donne $M(\ce{Na}) = \SI{23,0}{g.mol^{-1}}$.
\end{enumerate}
\tcblower
\textbf{Solution détaillée.}
\begin{enumerate}
\item L'éthanol \ce{C2H5OH} réagit avec le sodium. L'hydrogène du groupe $-$OH est arraché :
\[ \ce{2 C2H5OH + 2 Na -> 2 C2H5O^- Na^+ + H2 ^} \]
Le produit organique est l'\textbf{éthanolate de sodium} \ce{C2H5ONa}. Le gaz qui détone à la flamme est le \textbf{dihydrogène} \ce{H2}.
\item Cette réaction met en évidence la \textbf{mobilité de l'hydrogène du groupe hydroxyle} : le caractère légèrement acide de l'alcool (liaison O$-$H polarisée). Le virage au rose de la phénolphtaléine confirme le caractère basique de l'ion éthanolate \ce{C2H5O^-}, base forte.
\item Quantité de sodium :
\[ n(\ce{Na}) = \frac{m}{M} = \frac{0,23}{23,0} = \SI{1,0e-2}{mol} \]
D'après l'équation, $n(\ce{H2}) = \dfrac{n(\ce{Na})}{2} = \SI{5,0e-3}{mol}$.
\[ V(\ce{H2}) = n \times V_m = 5,0\times 10^{-3} \times 24,0 = \SI{0,12}{L} = \SI{120}{mL} \]
\end{enumerate}
\textbf{Conclusion} : le dégagement de \SI{120}{mL} de dihydrogène et la formation d'éthanolate de sodium prouvent la mobilité de l'hydrogène du groupe hydroxyle de l'éthanol.
\end{exoresolu}

\courssub{Fiche méthode 2 : Écrire les équations de déshydratation d'un alcool}

\begin{methode}[Déshydratation intra- ou intermoléculaire : bien choisir]
\begin{enumerate}
\item \textbf{Regarder les conditions opératoires} (c'est elles qui décident !) :
\begin{itemize}
\item \textbf{Intramoléculaire} $\to$ \cle{alcène} : acide sulfurique concentré (ou alumine \ce{Al2O3}), température \textbf{élevée} (environ \SI{180}{\celsius} avec \ce{H2SO4}, \SI{400}{\celsius} sur \ce{Al2O3}).
\item \textbf{Intermoléculaire} $\to$ \cle{éther-oxyde} : acide sulfurique concentré, température \textbf{modérée} (environ \SI{140}{\celsius}), alcool en excès.
\end{itemize}
\item \textbf{Intramoléculaire} : éliminer le $-$OH d'un carbone et un H du carbone \emph{voisin} (carbone en $\alpha$) ; créer la double liaison \ce{C=C} :
\[ \ce{R-CH2-CH2-OH -> R-CH=CH2 + H2O} \]
\item \textbf{Intermoléculaire} : condenser \textbf{deux} molécules d'alcool avec élimination d'\textbf{une} molécule d'eau :
\[ \ce{2 R-OH -> R-O-R + H2O} \]
\item \textbf{Vérifier} : bilan matière (même nombre de C, H, O de part et d'autre) et nommer les produits (alcène : suffixe « -ène » ; éther : « alcoxyalcane » ou nom d'usage, ex. éthoxyéthane).
\end{enumerate}
\textbf{Réflexe} : si l'alcool est dissymétrique, plusieurs alcènes sont possibles ; l'alcène majoritaire est le plus substitué (règle de Zaïtsev).
\end{methode}

\begin{exoresolu}[Déshydratation de l'éthanol : deux expériences]
\textbf{Énoncé (type Bac).} On chauffe de l'éthanol en présence d'acide sulfurique concentré.
\begin{itemize}
\item Expérience A : la température est maintenue vers \SI{180}{\celsius}. Le gaz formé décolore l'eau de brome.
\item Expérience B : la température est maintenue vers \SI{140}{\celsius} avec excès d'éthanol. On recueille un liquide d'odeur caractéristique, de formule \ce{C4H10O}.
\end{itemize}
\begin{enumerate}
\item Pour chaque expérience, écrire l'équation-bilan et nommer le(s) produit(s) organique(s).
\item Préciser le type de déshydratation mis en jeu dans chaque cas.
\end{enumerate}
\tcblower
\textbf{Solution détaillée.}
\begin{enumerate}
\item \textbf{Expérience A} (\SI{180}{\celsius}) : déshydratation \textbf{intramoléculaire}. Une molécule d'éthanol perd $-$OH et un H du carbone voisin :
\[ \ce{CH3-CH2-OH ->[H2SO4,\ \text{\SI{180}{\celsius}}] CH2=CH2 + H2O} \]
Le produit est l'\textbf{éthène} (éthylène) \ce{CH2=CH2}. Le test positif à l'eau de brome (décoloration) confirme la présence de la double liaison \ce{C=C}.

\textbf{Expérience B} (\SI{140}{\celsius}) : déshydratation \textbf{intermoléculaire}. Deux molécules d'éthanol se condensent :
\[ \ce{2 CH3-CH2-OH ->[H2SO4,\ \text{\SI{140}{\celsius}}] CH3-CH2-O-CH2-CH3 + H2O} \]
Le produit est l'\textbf{éthoxyéthane} (éther éthylique), de formule brute \ce{C4H10O}, conforme à l'énoncé.
\item Expérience A : déshydratation intramoléculaire (élimination au sein d'une même molécule). Expérience B : déshydratation intermoléculaire (condensation entre deux molécules).
\end{enumerate}
\textbf{Conclusion} : c'est la \textbf{température} qui oriente la réaction : vers \SI{180}{\celsius} on obtient l'alcène, vers \SI{140}{\celsius} l'éther-oxyde.
\end{exoresolu}

\courssub{Fiche méthode 3 : Réaliser et exploiter l'expérience de la lampe sans flamme}

\begin{methode}[Lampe sans flamme : protocole et exploitation]
\begin{enumerate}
\item \textbf{Dispositif} : un fil de cuivre (ou une spirale de cuivre) chauffé au rouge est plongé dans un erlenmeyer contenant de l'éthanol dont on a chassé l'air ; les vapeurs d'éthanol passent sur le cuivre incandescent.
\item \textbf{Observations à citer} : le fil de cuivre \textbf{reste incandescent} (la réaction est exothermique : elle s'auto-entretient, d'où « lampe sans flamme ») ; odeur piquante (pomme verte) ; le liquide recueilli \textbf{rougit le papier pH} et \textbf{rose le réactif de Schiff} (test des aldéhydes).
\item \textbf{Interpréter} : le cuivre joue le rôle de \cle{catalyseur} ; le dioxygène de l'air oxyde les vapeurs d'éthanol en \textbf{éthanal} :
\[ \ce{2 CH3-CH2-OH + O2 ->[Cu] 2 CH3-CHO + 2 H2O} \]
\item \textbf{Identifier les produits} : éthanal (réactif de Schiff, DNPH) et eau (sulfate de cuivre anhydre qui bleuit).
\end{enumerate}
\textbf{Réflexes} : le cuivre n'apparaît pas dans le bilan (catalyseur, écrit au-dessus de la flèche) ; ne pas confondre cette oxydation catalytique avec la combustion complète.
\end{methode}

\begin{exoresolu}[La lampe sans flamme au laboratoire]
\textbf{Énoncé (type Bac).} Un élève réalise l'expérience de la lampe sans flamme : il plonge une spirale de cuivre portée au rouge dans un flacon contenant des vapeurs d'éthanol. La spirale continue de rougeoyer. Le liquide formé rose le réactif de Schiff et donne un précipité jaune avec la 2,4-DNPH.
\begin{enumerate}
\item Expliquer pourquoi la spirale reste incandescent sans apport extérieur de chaleur.
\item Identifier le produit organique formé et écrire l'équation-bilan de la réaction.
\item Quel est le rôle du cuivre ? Justifier.
\end{enumerate}
\tcblower
\textbf{Solution détaillée.}
\begin{enumerate}
\item La réaction d'oxydation des vapeurs d'éthanol par le dioxygène est \textbf{exothermique} : la chaleur dégagée par la réaction maintient la spirale de cuivre au rouge. La réaction s'auto-entretient : c'est pourquoi on parle de « lampe sans flamme ».
\item Le précipité jaune avec la DNPH indique un composé \textbf{carbonylé} ; le réactif de Schiff rosé précise qu'il s'agit d'un \textbf{aldéhyde}. L'oxydation ménagée de l'éthanol (alcool primaire) donne l'\textbf{éthanal} \ce{CH3-CHO} :
\[ \ce{2 CH3-CH2-OH + O2 ->[Cu] 2 CH3-CHO + 2 H2O} \]
\item Le cuivre est un \textbf{catalyseur} : il accélère la réaction sans figurer dans l'équation-bilan (il est régénéré en fin de réaction). On l'écrit au-dessus de la flèche.
\end{enumerate}
\textbf{Conclusion} : la lampe sans flamme illustre l'oxydation catalytique ménagée de l'éthanol en éthanal, réaction exothermique catalysée par le cuivre.
\end{exoresolu}

\courssub{Fiche méthode 4 : Réaliser l'oxydation ménagée selon la classe de l'alcool}

\begin{methode}[Oxydation ménagée par \ce{KMnO4} : prévoir les produits]
\begin{enumerate}
\item \textbf{Déterminer la classe de l'alcool} : compter le nombre de groupes alkyles portés par le carbone fonctionnel (celui qui porte le $-$OH) : 1 $\to$ primaire, 2 $\to$ secondaire, 3 $\to$ tertiaire.
\item \textbf{Appliquer la règle d'oxydation} (oxydant : \ce{KMnO4} en milieu acide, couple \ce{MnO4^-}/\ce{Mn^2+}) :
\begin{itemize}
\item Alcool \textbf{primaire} $\to$ \cle{aldéhyde} puis, si l'oxydant est en excès, \cle{acide carboxylique} ;
\item Alcool \textbf{secondaire} $\to$ \cle{cétone} (qui ne s'oxyde plus) ;
\item Alcool \textbf{tertiaire} $\to$ \textbf{pas d'oxydation ménagée} (pas d'hydrogène sur le carbone fonctionnel).
\end{itemize}
\item \textbf{Écrire les demi-équations} puis l'équation-bilan. Demi-équation de l'oxydant :
\[ \ce{MnO4^- + 8 H^+ + 5 e^- -> Mn^2+ + 4 H2O} \]
Exemples de demi-équations organiques :
\[ \ce{R-CH2-OH -> R-CHO + 2 H^+ + 2 e^-} \qquad \ce{R-CH2-OH + H2O -> R-COOH + 4 H^+ + 4 e^-} \]
\[ \ce{R1-CH(OH)-R2 -> R1-CO-R2 + 2 H^+ + 2 e^-} \]
\item \textbf{Multiplier pour égaliser les électrons}, additionner, simplifier les \ce{H^+} et \ce{H2O}.
\item \textbf{Observation expérimentale} : décoloration de la solution violette de permanganate. \textbf{Identifier les produits} : DNPH (précipité jaune : carbonylé), réactif de Schiff (aldéhyde), liqueur de Fehling (aldéhyde), papier pH (acide).
\end{enumerate}

\begin{attention}[Piège classique]
L'oxydation ménagée d'un alcool \textbf{ne modifie pas la chaîne carbonée} : le squelette reste identique, seul le groupe fonctionnel change. Et n'oublie jamais : un alcool tertiaire ne réagit pas — écrire « pas de réaction » peut rapporter le point !
\end{attention}
\end{methode}

\begin{exoresolu}[Oxydation ménagée du propan-1-ol et du propan-2-ol]
\textbf{Énoncé (type Bac).} On dispose de deux alcools isomères A et B de formule brute \ce{C3H8O}. A est le propan-1-ol, B le propan-2-ol. On réalise l'oxydation ménagée de chacun par une solution acidifiée de permanganate de potassium.
\begin{enumerate}
\item Donner la classe de chaque alcool.
\item Écrire l'équation-bilan de l'oxydation de A en aldéhyde par \ce{MnO4^-} en milieu acide.
\item Quel produit obtient-on avec B ? Comment l'identifier expérimentalement ?
\item Que se passe-t-il si on poursuit l'oxydation de A avec un excès de permanganate ?
\end{enumerate}
\tcblower
\textbf{Solution détaillée.}
\begin{enumerate}
\item A : propan-1-ol \ce{CH3-CH2-CH2-OH} : le carbone fonctionnel porte un seul groupe alkyle $\to$ alcool \textbf{primaire}. B : propan-2-ol \ce{CH3-CH(OH)-CH3} : deux groupes alkyles $\to$ alcool \textbf{secondaire}.
\item Demi-équations :
\[ \ce{MnO4^- + 8 H^+ + 5 e^- -> Mn^2+ + 4 H2O} \quad (\times 2) \]
\[ \ce{CH3-CH2-CH2-OH -> CH3-CH2-CHO + 2 H^+ + 2 e^-} \quad (\times 5) \]
En additionnant et en simplifiant ($16\,\ce{H^+} - 10\,\ce{H^+} = 6\,\ce{H^+}$) :
\[ \ce{2 MnO4^- + 6 H^+ + 5 CH3-CH2-CH2-OH -> 2 Mn^2+ + 8 H2O + 5 CH3-CH2-CHO} \]
Le produit est le \textbf{propanal}. On observe la décoloration du permanganate.
\item B, alcool secondaire, s'oxyde en \textbf{cétone} : la \textbf{propanone} (acétone) \ce{CH3-CO-CH3} :
\[ \ce{2 MnO4^- + 6 H^+ + 5 CH3-CH(OH)-CH3 -> 2 Mn^2+ + 8 H2O + 5 CH3-CO-CH3} \]
Identification : précipité jaune avec la DNPH (composé carbonylé) mais \textbf{réactif de Schiff négatif} (reste incolore) et liqueur de Fehling négative $\to$ c'est une cétone, pas un aldéhyde.
\item Avec un excès de permanganate, le propanal est oxydé à son tour en \textbf{acide propanoïque} \ce{CH3-CH2-COOH} (la solution devient acide, pH $< 7$).
\end{enumerate}
\textbf{Conclusion} : la classe de l'alcool commande le produit d'oxydation : primaire $\to$ aldéhyde $\to$ acide ; secondaire $\to$ cétone ; tertiaire $\to$ rien.
\end{exoresolu}

\courssub{Fiche méthode 5 : Oxydation ménagée des aldéhydes}

\begin{methode}[Aldéhyde $\to$ acide carboxylique : rédaction propre]
\begin{enumerate}
\item \textbf{Reconnaître l'aldéhyde} : groupe carbonyle \ce{C=O} en \textbf{bout de chaîne} (\ce{R-CHO}), avec un H sur le carbone fonctionnel — c'est ce H qui permet l'oxydation (les cétones n'en ont pas et ne s'oxydent pas).
\item \textbf{Écrire la demi-équation d'oxydation} (milieu acide) :
\[ \ce{R-CHO + H2O -> R-COOH + 2 H^+ + 2 e^-} \]
\item \textbf{Combiner avec l'oxydant} (ex. \ce{MnO4^-}/\ce{Mn^2+} : multiplier l'aldéhyde par 5, le permanganate par 2) et simplifier.
\item \textbf{Tests d'identification à connaître} : réactif de Schiff (rose), liqueur de Fehling (précipité rouge brique \ce{Cu2O}), réactif de Tollens (miroir d'argent) — ces tests distinguent l'aldéhyde de la cétone.
\end{enumerate}
\end{methode}

\begin{exoresolu}[Oxydation de l'éthanal]
\textbf{Énoncé (type Bac).} On verse quelques gouttes d'éthanal dans une solution acidifiée de permanganate de potassium : la solution se décolore. Le produit obtenu rougit le papier pH.
\begin{enumerate}
\item Écrire l'équation-bilan de la réaction.
\item Nommer le produit organique et justifier l'observation au papier pH.
\end{enumerate}
\tcblower
\textbf{Solution détaillée.}
\begin{enumerate}
\item Demi-équations :
\[ \ce{CH3-CHO + H2O -> CH3-COOH + 2 H^+ + 2 e^-} \quad (\times 5) \]
\[ \ce{MnO4^- + 8 H^+ + 5 e^- -> Mn^2+ + 4 H2O} \quad (\times 2) \]
Bilan (après simplification : $16\,\ce{H^+} - 10\,\ce{H^+} = 6\,\ce{H^+}$ ; $8\,\ce{H2O} - 5\,\ce{H2O} = 3\,\ce{H2O}$) :
\[ \ce{5 CH3-CHO + 2 MnO4^- + 6 H^+ -> 5 CH3-COOH + 2 Mn^2+ + 3 H2O} \]
\item Le produit est l'\textbf{acide éthanoïque} (acide acétique) \ce{CH3-COOH}. C'est un acide carboxylique : il libère des ions \ce{H3O+} en solution, d'où le rougissement du papier pH.
\end{enumerate}
\textbf{Conclusion} : l'oxydation ménagée d'un aldéhyde conduit à l'acide carboxylique de même chaîne carbonée ; c'est la suite logique de l'oxydation d'un alcool primaire.
\end{exoresolu}

\courssub{Fiche méthode 6 : Écrire l'équation-bilan d'une estérification}

\begin{methode}[Estérification : acide carboxylique + alcool]
\begin{enumerate}
\item \textbf{Identifier les réactifs} : un \cle{acide carboxylique} $R\text{-}COOH$ et un \cle{alcool} $R'\text{-}OH$.
\item \textbf{Écrire les produits} : un \cle{ester} $R\text{-}COO\text{-}R'$ et de l'\cle{eau}. L'assemblage : le groupe acyle $R\text{-}CO$ de l'acide se lie au groupe $\text{-}O\text{-}R'$ de l'alcool :
\[ R\text{-}COOH + R'\text{-}OH \rightleftharpoons R\text{-}COO\text{-}R' + H_2O \]
\item \textbf{Nommer l'ester} : « alkanoate d'alkyle » : la partie acide donne « -oate », la partie alcool donne l'alkyle. Exemple : \ce{CH3-COO-CH2-CH3} = éthanoate d'éthyle.
\item \textbf{Citer les caractéristiques} : réaction \textbf{lente}, \textbf{limitée} (équilibre, double flèche \ce{<=>}), \textbf{athermique}. Pour l'accélérer : chauffer, catalyseur acide (\ce{H2SO4}). Pour augmenter le rendement : excès d'un réactif ou élimination d'un produit.
\end{enumerate}
\textbf{Réflexe} : vérifier le bilan matière ; la classe de l'alcool influence le rendement à l'équilibre (primaire $\approx$ 66 \%, secondaire $\approx$ 60 \%, tertiaire $\approx$ 5 \% pour un mélange équimolaire).
\end{methode}

\begin{exoresolu}[Synthèse de l'éthanoate d'éthyle]
\textbf{Énoncé (type Bac).} On chauffe en présence de quelques gouttes d'acide sulfurique concentré un mélange de $n_a = \SI{0,20}{mol}$ d'acide éthanoïque et de $n_b = \SI{0,20}{mol}$ d'éthanol. À l'équilibre, il s'est formé \SI{0,13}{mol} d'ester.
\begin{enumerate}
\item Écrire l'équation-bilan de la réaction et nommer l'ester formé.
\item Donner deux caractéristiques de cette réaction et le rôle de l'acide sulfurique.
\item Dresser le tableau d'avancement et calculer le rendement de l'esterification.
\end{enumerate}
\tcblower
\textbf{Solution détaillée.}
\begin{enumerate}
\item Équation-bilan :
\[ \ce{CH3-COOH + CH3-CH2-OH <=> CH3-COO-CH2-CH3 + H2O} \]
L'ester formé est l'\textbf{éthanoate d'éthyle} (odeur fruitée caractéristique).
\item La réaction est \textbf{lente} et \textbf{limitée} (elle aboutit à un équilibre chimique). L'acide sulfurique est un \textbf{catalyseur} : il accélère la réaction sans modifier l'état d'équilibre.
\item Tableau d'avancement (en mol) :
\begin{center}
\begin{tabularx}{\linewidth}{|l|X|X|X|X|}
\hline
\rowcolor{popblueL} État & \ce{CH3COOH} & \ce{C2H5OH} & Ester & \ce{H2O} \\
\hline
Initial ($x=0$) & 0,20 & 0,20 & 0 & 0 \\
\hline
En cours ($x$) & $0,20 - x$ & $0,20 - x$ & $x$ & $x$ \\
\hline
Équilibre ($x_f$) & $0,20 - x_f$ & $0,20 - x_f$ & $x_f = 0,13$ & $0,13$ \\
\hline
\end{tabularx}
\end{center}
Avancement maximal (réactifs en proportions stœchiométriques) : $x_{max} = \SI{0,20}{mol}$.
Avancement final : $x_f = \SI{0,13}{mol}$.
\[ \eta = \frac{x_f}{x_{max}} \times 100 = \frac{0,13}{0,20} \times 100 = 65\ \% \]
\end{enumerate}
\textbf{Conclusion} : le rendement de 65 \% est conforme à la limite théorique ($\approx$ 66 \%) d'une estérification équimolaire avec un alcool primaire : la réaction est bien limitée par l'hydrolyse inverse de l'ester.
\end{exoresolu}

\courssub{Fiche méthode 7 : Préparer un alcool — hydratation d'un alcène et fermentation alcoolique}

\begin{methode}[Deux voies de préparation à maîtriser]
\begin{enumerate}
\item \textbf{Hydratation d'un alcène} : addition d'eau sur la double liaison \ce{C=C}, en présence d'un catalyseur acide :
\[ \ce{R-CH=CH2 + H2O ->[H^+] R-CH2-CH2-OH} \]
Réflexe : pour un alcène dissymétrique, le $-$OH se fixe sur le carbone \textbf{le plus substitué} (règle de Markovnikov) : l'hydratation du propène donne majoritairement le propan-2-ol.
\item \textbf{Fermentation alcoolique} (anaérobie) : transformation du glucose en éthanol et dioxyde de carbone sous l'action d'enzymes (zymase des levures), \textbf{en l'absence d'air} :
\[ \ce{C6H12O6 ->[enzymes] 2 C2H5OH + 2 CO2} \]
Réflexes : vérifier l'équilibrage (6 C, 12 H, 6 O de chaque côté) ; le dégagement de \ce{CO2} trouble l'eau de chaux ; citer les applications locales (vin de palme, bili-bili, matango).
\item \textbf{Choisir la voie} selon l'énoncé : voie industrielle/synthèse $\to$ hydratation ; voie biologique/boissons $\to$ fermentation.
\end{enumerate}
\end{methode}

\begin{exoresolu}[Du palmier à l'éthanol]
\textbf{Énoncé (type Bac).} Le vin de palme, boisson traditionnelle obtenue à partir de la sève du palmier, contient de l'éthanol issu de la fermentation des sucres par des levures naturelles.
\begin{enumerate}
\item Écrire l'équation-bilan de la fermentation alcoolique du glucose. Pourquoi dit-on qu'elle est anaérobie ?
\item Quel volume de dioxyde de carbone (mesuré à $V_m = \SI{24,0}{L.mol^{-1}}$) accompagne la formation de $m = \SI{46}{g}$ d'éthanol ? On donne $M(\ce{C2H5OH}) = \SI{46,0}{g.mol^{-1}}$.
\item Par ailleurs, l'industrie prépare l'éthanol par hydratation de l'éthène. Écrire l'équation de cette réaction.
\end{enumerate}
\tcblower
\textbf{Solution détaillée.}
\begin{enumerate}
\item Équation-bilan de la fermentation alcoolique :
\[ \ce{C6H12O6 ->[zymase] 2 C2H5OH + 2 CO2} \]
Vérification : 6 C, 12 H, 6 O de chaque côté. Elle est dite \textbf{anaérobie} car elle se déroule \textbf{en l'absence de dioxygène} : les levures dégradent le glucose sans utiliser l'air (en présence d'air, l'éthanol serait oxydé en acide éthanoïque — le vin « tourne au vinaigre »).
\item Quantité d'éthanol formé :
\[ n(\ce{C2H5OH}) = \frac{m}{M} = \frac{46}{46,0} = \SI{1,0}{mol} \]
D'après l'équation, $n(\ce{CO2}) = n(\ce{C2H5OH}) = \SI{1,0}{mol}$ (même coefficient stœchiométrique 2).
\[ V(\ce{CO2}) = n \times V_m = 1,0 \times 24,0 = \SI{24}{L} \]
\item Hydratation de l'éthène :
\[ \ce{CH2=CH2 + H2O ->[H^+] CH3-CH2-OH} \]
\end{enumerate}
\textbf{Conclusion} : la fermentation de \SI{180}{g} de glucose produit 2 mol d'éthanol et 2 mol de \ce{CO2} ; ici, \SI{46}{g} d'éthanol s'accompagnent de \SI{24}{L} de dioxyde de carbone. L'hydratation de l'éthène offre une voie de synthèse industrielle complémentaire.
\end{exoresolu}

\begin{attention}[Les 5 erreurs qui coûtent des points au Bac]
\begin{enumerate}
\item \textbf{Écrire \ce{R-OH + Na -> R-ONa + H2} non équilibrée} : il faut des coefficients 2, 2, 2, 1. Le dihydrogène est \ce{H2}, jamais H !
\item \textbf{Confondre les deux déshydratations} : alcène à haute température (\SI{180}{\celsius}), éther-oxyde à température modérée (\SI{140}{\celsius}). Inverser les conditions, c'est inverser les produits — et perdre tous les points de la question.
\item \textbf{Oxyder un alcool tertiaire} : un alcool tertiaire ne subit \textbf{pas} d'oxydation ménagée (pas d'hydrogène sur le carbone fonctionnel). Écris « pas de réaction » et justifie.
\item \textbf{Arrêter l'oxydation d'un alcool primaire à l'aldéhyde quand l'oxydant est en excès} : avec excès de \ce{KMnO4}, on va jusqu'à l'\textbf{acide carboxylique}. Lis bien l'énoncé !
\item \textbf{Oublier les caractéristiques de l'estérification} : réaction \textbf{lente et limitée} — la double flèche \ce{<=>} est obligatoire, et l'acide sulfurique est un catalyseur (il ne déplace pas l'équilibre, il l'atteint plus vite). Bonus : dans la fermentation, n'oublie pas le \ce{CO2} dans le bilan !
\end{enumerate}
\end{attention}

\newpage

\coursec{Exercices}

\courssub{Je vérifie mes connaissances}

\begin{exercice}[QCM : les bons réflexes]
Pour chaque question, choisis la (ou les) bonne(s) réponse(s).
\begin{enumerate}
  \item La réaction d'un alcool avec le sodium dégage :
  \begin{enumerate}
    \item du dioxygène
    \item du dihydrogène
    \item du dioxyde de carbone
    \item de la vapeur d'eau
  \end{enumerate}
  \item La déshydratation intramoléculaire de l'éthanol en présence d'alumine \ce{Al2O3} vers \SI{400}{\celsius} donne :
  \begin{enumerate}
    \item l'éthanal
    \item l'éthylène
    \item l'éther diéthylique
    \item l'acide éthanoïque
  \end{enumerate}
  \item L'oxydation ménagée d'un alcool secondaire conduit à :
  \begin{enumerate}
    \item un aldéhyde
    \item une cétone
    \item un acide carboxylique
    \item un ester
  \end{enumerate}
  \item La réaction d'estérification est :
  \begin{enumerate}
    \item totale et rapide
    \item lente et limitée par une réaction inverse
    \item catalysée par les ions \ce{H3O+}
    \item exothermique
  \end{enumerate}
\end{enumerate}
\end{exercice}

\begin{exercice}[Vrai ou faux — justifie ta réponse]
\begin{enumerate}
  \item Un alcool tertiaire ne subit pas d'oxydation ménagée par le permanganate de potassium.
  \item La fermentation alcoolique du glucose est une réaction aérobie.
  \item Le test à la DNPH permet de distinguer un alcool d'un aldéhyde.
  \item La déshydratation intermoléculaire de deux molécules d'éthanol donne un éther-oxyde.
  \item L'hydratation du propène conduit uniquement au propan-1-ol.
\end{enumerate}
\end{exercice}

\begin{exercice}[Texte à trous : la lampe sans flamme]
Recopie et complète le texte suivant avec les mots : \emph{catalyseur, éthanal, cuivre, incandescent, oxydation, exothermique}.

Dans l'expérience de la lampe sans flamme, un fil de \dots\ldots{} porté au rouge est plongé au-dessus de l'éthanol bouillant. Les vapeurs d'alcool subissent une \dots\ldots{} ménagée au contact du métal qui joue le rôle de \dots\ldots{}. Il se forme de l'\dots\ldots{} à l'odeur piquante. La réaction étant \dots\ldots{}, le fil reste \dots\ldots{} bien que la flamme ait été retirée.
\end{exercice}

\begin{exercice}[Définitions essentielles]
\begin{enumerate}
  \item Définis : oxydation ménagée ; estérification ; fermentation alcoolique.
  \item Qu'appelle-t-on \cle{classe} d'un alcool ? Donne les trois classes et un exemple pour chacune.
  \item Cite les deux types de déshydratation d'un alcool et précise pour chacun la nature du produit obtenu.
\end{enumerate}
\end{exercice}

\courssub{Je m'entraîne}

\begin{exercice}[Réaction avec le sodium]
On fait réagir le propan-1-ol avec le sodium.
\begin{enumerate}
  \item Écris l'équation-bilan de la réaction.
  \item Nomme le produit organique formé.
  \item Quelle propriété du groupe hydroxyle cette réaction met-elle en évidence ?
\end{enumerate}
\end{exercice}

\begin{exercice}[Déshydratations du butan-1-ol]
On chauffe du butan-1-ol en présence d'acide sulfurique concentré.
\begin{enumerate}
  \item Écris l'équation de la déshydratation intramoléculaire et nomme le produit obtenu. Précise les conditions expérimentales.
  \item Écris l'équation de la déshydratation intermoléculaire et nomme le produit obtenu. Précise les conditions expérimentales.
  \item Quel paramètre permet d'orienter la réaction vers l'un ou l'autre produit ?
\end{enumerate}
\end{exercice}

\begin{exercice}[Oxydation ménagée et tests d'identification]
L'oxydation ménagée du propan-1-ol par une solution acidifiée de permanganate de potassium conduit d'abord à un composé A, puis, en excès d'oxydant, à un composé B.
\begin{enumerate}
  \item Donne les formules semi-développées et les noms de A et B.
  \item Propose deux tests permettant d'identifier A (on précisera le réactif et le résultat observé).
  \item Comment mettre en évidence B sans test organique ?
  \item Écris l'équation-bilan ionique de la formation de A. Couples : \ce{MnO4-}/\ce{Mn^2+} et \ce{CH3CH2CHO}/\ce{CH3CH2CH2OH}.
\end{enumerate}
\end{exercice}

\begin{exercice}[Estérification]
On chauffe un mélange équimolaire d'acide éthanoïque et de propan-1-ol en présence de quelques gouttes d'acide sulfurique concentré.
\begin{enumerate}
  \item Écris l'équation-bilan de la réaction.
  \item Nomme l'ester formé.
  \item Cite deux caractéristiques de cette réaction.
  \item Propose deux moyens d'augmenter le rendement en ester.
\end{enumerate}
\end{exercice}

\courssub{Je me prépare au Bac}

\begin{exercice}[Type Bac — Le sodium et l'éthanol]
\emph{Données :} $M(\ce{Na}) = \SI{23,0}{g.mol^{-1}}$ ; $M(\ce{C2H5ONa}) = \SI{68,0}{g.mol^{-1}}$ ; volume molaire dans les conditions de l'expérience : $V_m = \SI{24,0}{L.mol^{-1}}$.

Au laboratoire du lycée, on introduit un morceau de sodium de masse $m = \SI{4,6}{g}$ dans un excès d'éthanol absolu contenu dans un erlenmeyer. On observe une effervescence et le dégagement d'un gaz qui détone à la flamme.
\begin{enumerate}
  \item Identifier le gaz dégagé et le produit organique formé.
  \item Écrire l'équation-bilan de la réaction.
  \item Calculer la quantité de matière de sodium utilisée, puis le volume de gaz dégagé.
  \item Calculer la masse d'éthanolate de sodium formée.
  \item Citer deux précautions à prendre lors de cette manipulation.
\end{enumerate}
\end{exercice}

\begin{exercice}[Type Bac — Étude d'un alcool \ce{C4H10O}]
\emph{Données :} $M(\ce{C}) = \SI{12,0}{g.mol^{-1}}$ ; $M(\ce{H}) = \SI{1,0}{g.mol^{-1}}$ ; $M(\ce{O}) = \SI{16,0}{g.mol^{-1}}$.

Un alcool A de formule brute \ce{C4H10O} subit une déshydratation intramoléculaire en présence d'alumine chauffée ; on obtient majoritairement le but-2-ène.
\begin{enumerate}
  \item Donner les formules semi-développées et les noms des quatre alcools isomères de formule \ce{C4H10O}, en précisant la classe de chacun.
  \item Identifier A en justifiant la réponse.
  \item Écrire l'équation de la déshydratation intramoléculaire de A en précisant les conditions expérimentales.
  \item Écrire l'équation de la déshydratation intermoléculaire de A et nommer le produit formé. Quelles conditions permettent de favoriser cette voie ?
  \item L'oxydation ménagée de A par le permanganate de potassium en milieu acide conduit à un composé D qui réagit avec la DNPH mais pas avec la liqueur de Fehling. Confirmer l'identification de A et donner la formule semi-développée et le nom de D.
  \item Calculer la masse d'alcool A nécessaire pour obtenir théoriquement \SI{5,6}{L} de but-2-ène (volume molaire $V_m = \SI{22,4}{L.mol^{-1}}$), en supposant la déshydratation totale.
\end{enumerate}
\end{exercice}

\begin{exercice}[Type Bac — Synthèse d'un arôme à la brasserie]
\emph{Données :} $M(\ce{C6H12O6}) = \SI{180}{g.mol^{-1}}$ ; $M(\ce{C2H5OH}) = \SI{46,0}{g.mol^{-1}}$ ; densité de l'éthanol $d = 0,79$ ; $V_m = \SI{24,0}{L.mol^{-1}}$.

Une brasserie camerounaise fermente une cuve contenant \SI{500}{kg} de glucose. L'éthanol obtenu servira, en partie, à la préparation d'un arôme : l'éthanoate d'éthyle.
\begin{enumerate}
  \item Écrire l'équation-bilan de la fermentation alcoolique du glucose. Préciser les conditions (enzymes, anaérobiose, température).
  \item Calculer la masse puis le volume d'éthanol théoriquement obtenus.
  \item Calculer le volume de dioxyde de carbone dégagé pendant la fermentation.
  \item Pour préparer l'éthanoate d'éthyle, on oxyde d'abord une partie de l'éthanol en acide éthanoïque par oxydation ménagée (couple \ce{CH3COOH}/\ce{CH3CH2OH}, oxydant \ce{MnO4-}/\ce{Mn^2+} en milieu acide). Écrire les deux demi-équations puis l'équation-bilan ionique de cette oxydation.
  \item Écrire l'équation de la réaction entre l'acide éthanoïque et l'éthanol. Nommer l'ester obtenu et donner deux caractéristiques de cette réaction.
  \item Le rendement de l'estérification à partir d'un mélange équimolaire est de 67 \%. Calculer la quantité de matière d'ester obtenue à partir de \SI{1,00}{mol} d'acide éthanoïque et \SI{1,00}{mol} d'éthanol, puis proposer un moyen d'améliorer ce rendement.
\end{enumerate}
\end{exercice}

\newpage

\coursec{Corrigés des exercices}

\begin{corrige}[Exercice 1 — QCM : les bons réflexes]
\begin{enumerate}
  \item \textbf{Réponse b}. La réaction d'un alcool avec le sodium dégage du \cle{dihydrogène} \ce{H2}, qui détone à la flamme (test de l'aboiement) : $\ce{2 R-OH + 2 Na -> 2 R-O^-Na+ + H2 ^}$.
  \item \textbf{Réponse b}. Vers \SI{400}{\celsius} sur alumine \ce{Al2O3}, la déshydratation est \cle{intramoléculaire} : l'éthanol perd une molécule d'eau et donne l'\cle{éthylène} (éthène) : $\ce{CH3-CH2-OH ->[Al2O3,\ 400\ ^{\circ}C] CH2=CH2 + H2O}$.
  \item \textbf{Réponse b}. L'oxydation ménagée d'un alcool \cle{secondaire} conduit toujours à une \cle{cétone} (ex. : propan-2-ol $\rightarrow$ propanone).
  \item \textbf{Réponses b et c}. L'estérification est \cle{lente} et \cle{limitée} par la réaction inverse (hydrolyse) : c'est un équilibre. Elle est \cle{catalysée} par les ions \ce{H3O+} (apportés par quelques gouttes d'acide sulfurique concentré). Elle n'est ni totale, ni rapide, et elle est pratiquement \emph{athermique} (pas exothermique).
\end{enumerate}
\end{corrige}

\begin{attention}[Points de vigilance — Exercice 1]
Retenir les conditions typiques : \ce{Na} métallique $\rightarrow$ \ce{H2} (gaz) + alcoolate ; \ce{Al2O3}/\SI{400}{\celsius} $\rightarrow$ alcène (intramoléculaire) ; \ce{H2SO4} conc. \SI{140}{\celsius} + excès alcool $\rightarrow$ éther (intermoléculaire) ; \ce{H2SO4} conc. \SI{180}{\celsius} $\rightarrow$ alcène. L'estérification est un équilibre athermique catalysé par \ce{H+}.
\end{attention}

\begin{corrige}[Exercice 2 — Vrai ou faux]
\begin{enumerate}
  \item \textbf{Vrai.} Un alcool tertiaire ne possède \textbf{pas d'atome d'hydrogène} sur le carbone fonctionnel (celui qui porte le \ce{-OH}) ; il ne peut donc pas subir d'oxydation ménagée sans rupture de la chaîne carbonée. La solution de \ce{KMnO4} reste violette.
  \item \textbf{Faux.} La fermentation alcoolique est une réaction \cle{anaérobie} : elle se déroule \emph{en l'absence de dioxygène}, sous l'action des enzymes des levures, vers \SI{30}{\celsius}.
  \item \textbf{Vrai.} La DNPH (2,4-dinitrophénylhydrazine) donne un précipité jaune-orangé avec les composés carbonylés (aldéhydes et cétones) et aucune réaction avec les alcools. Elle permet donc de distinguer un alcool d'un aldéhyde.
  \item \textbf{Vrai.} Deux molécules d'éthanol perdent une molécule d'eau et forment l'éthoxyéthane (éther diéthylique) : $\ce{2 C2H5OH -> C2H5-O-C2H5 + H2O}$.
  \item \textbf{Faux.} Le propène est un alcène \emph{asymétrique} : son hydratation conduit majoritairement au \cle{propan-2-ol} (règle de Markovnikov : l'hydrogène se fixe sur le carbone le plus hydrogéné), le propan-1-ol n'étant que minoritaire.
\end{enumerate}
\end{corrige}

\begin{attention}[Points de vigilance — Exercice 2]
Ne pas confondre anaérobie (sans \ce{O2}) et aérobie. La règle de Markovnikov s'applique à l'addition d'eau sur un alcène asymétrique : l'\ce{H} va sur le C le plus hydrogéné, l'\ce{OH} sur l'autre. Un alcool tertiaire ne s'oxyde pas ménagée (pas de \ce{H} sur le C fonctionnel).
\end{attention}

\begin{corrige}[Exercice 3 — Texte à trous : la lampe sans flamme]
Dans l'expérience de la lampe sans flamme, un fil de \textbf{cuivre} porté au rouge est plongé au-dessus de l'éthanol bouillant. Les vapeurs d'alcool subissent une \textbf{oxydation} ménagée au contact du métal qui joue le rôle de \textbf{catalyseur}. Il se forme de l'\textbf{éthanal} à l'odeur piquante. La réaction étant \textbf{exothermique}, le fil reste \textbf{incandescent} bien que la flamme ait été retirée.
\par\medskip
Équation de la réaction : $\ce{2 CH3-CH2-OH + O2 ->[Cu] 2 CH3-CHO + 2 H2O}$.
\end{corrige}

\begin{attention}[Points de vigilance — Exercice 3]
Le cuivre n'est pas consommé (catalyseur). L'éthanal (\ce{CH3CHO}) a une odeur caractéristique de pomme verte/piquante. L'exothermicité maintient le fil au rouge : c'est une oxydation catalytique hétérogène.
\end{attention}

\begin{corrige}[Exercice 4 — Définitions essentielles]
\begin{enumerate}
  \item \textbf{Oxydation ménagée} : oxydation qui modifie le groupe fonctionnel d'une espèce organique \emph{sans rompre la chaîne carbonée} (ex. : alcool primaire $\rightarrow$ aldéhyde $\rightarrow$ acide carboxylique). \textbf{Estérification} : réaction entre un acide carboxylique et un alcool qui donne un \cle{ester} et de l'eau ; elle est lente, athermique et limitée par l'hydrolyse de l'ester. \textbf{Fermentation alcoolique} : transformation, en l'absence de dioxygène et sous l'action d'enzymes de levures, du glucose en éthanol et dioxyde de carbone : $\ce{C6H12O6 -> 2 C2H5OH + 2 CO2}$.
  \item La \cle{classe} d'un alcool est le nombre de groupes alkyles (de carbones) liés au carbone fonctionnel (celui qui porte le \ce{-OH}). \textbf{Primaire} : un seul alkyle, ex. éthanol \chemfig{CH_3-CH_2-OH} ; \textbf{secondaire} : deux alkyles, ex. propan-2-ol \chemfig{H_3C-CH(-[2]OH)-CH_3} ; \textbf{tertiaire} : trois alkyles, ex. 2-méthylpropan-2-ol \chemfig{H_3C-C(-[2]OH)(-[6]CH_3)-CH_3}.
  \item \textbf{Déshydratation intramoléculaire} : une seule molécule d'alcool perd \ce{H2O} ; on obtient un \cle{alcène} (conditions : \ce{Al2O3} vers \SI{400}{\celsius}, ou \ce{H2SO4} concentré vers \SI{180}{\celsius}). \textbf{Déshydratation intermoléculaire} : deux molécules d'alcool perdent ensemble une molécule d'eau ; on obtient un \cle{éther-oxyde} (conditions : \ce{H2SO4} concentré vers \SI{140}{\celsius}, alcool en excès).
\end{enumerate}
\end{corrige}

\begin{attention}[Points de vigilance — Exercice 4]
Bien distinguer les trois classes d'alcools (nombre de C sur le C-OH). Pour la déshydratation, la température est le paramètre clé : $\approx$\SI{140}{\celsius} $\rightarrow$ éther ; $\approx$\SI{180}{\celsius} $\rightarrow$ alcène. L'estérification est un équilibre : lent, athermique, catalysé \ce{H+}.
\end{attention}

\begin{corrige}[Exercice 5 — Réaction avec le sodium]
\begin{enumerate}
  \item Équation-bilan :
  \[ \ce{2 CH3-CH2-CH2-OH + 2 Na -> 2 CH3-CH2-CH2-O^-Na+ + H2 ^} \]
  \item Le produit organique est le \cle{propan-1-olate de sodium} (ou propanolate de sodium), \ce{CH3CH2CH2ONa}.
  \item Cette réaction met en évidence la \cle{mobilité de l'atome d'hydrogène} du groupe hydroxyle \ce{-OH} : l'alcool se comporte comme un acide faible vis-à-vis du sodium, qui arrache cet hydrogène en le réduisant en \ce{H2}.
\end{enumerate}
\end{corrige}

\begin{attention}[Points de vigilance — Exercice 5]
L'alcoolate formé est ionique (\ce{R-O^- Na+}), soluble dans l'alcool. Le gaz \ce{H2} est inflammable (détonant) : éloigner toute flamme. Manipuler le Na avec des pinces (brûlures par réaction avec l'humidité cutanée).
\end{attention}

\begin{corrige}[Exercice 6 — Déshydratations du butan-1-ol]
\begin{enumerate}
  \item \textbf{Intramoléculaire} (une molécule, élimination de \ce{H2O}) :
  \[ \ce{CH3-CH2-CH2-CH2-OH ->[\text{H2SO4 conc., $\approx$ 180 °C}] CH3-CH2-CH=CH2 + H2O} \]
  Le produit est le \cle{but-1-ène}. Conditions : acide sulfurique concentré et chauffage vers \SI{180}{\celsius} (ou passage des vapeurs sur \ce{Al2O3} vers \SI{400}{\celsius}).
  \item \textbf{Intermoléculaire} (deux molécules) :
  \[ \ce{2 CH3-CH2-CH2-CH2-OH ->[\text{H2SO4 conc., $\approx$ 140 °C}] CH3-CH2-CH2-CH2-O-CH2-CH2-CH2-CH3 + H2O} \]
  Le produit est l'\cle{éther dibutylique} (1-butoxybutane). Conditions : température plus douce, vers \SI{140}{\celsius}, avec excès d'alcool.
  \item Le paramètre déterminant est la \cle{température} (associée à la proportion d'alcool) : température élevée ($\approx$ \SI{180}{\celsius}) $\rightarrow$ alcène ; température modérée ($\approx$ \SI{140}{\celsius}) $\rightarrow$ éther-oxyde.
\end{enumerate}
\end{corrige}

\begin{attention}[Points de vigilance — Exercice 6]
Le butan-1-ol ne donne qu'un seul alcène par déshydratation intramoléculaire (but-1-ène), pas d'isomérie de position possible. L'éther formé est symétrique. Respecter les températures : \SI{140}{\celsius} (éther) vs \SI{180}{\celsius} (alcène).
\end{attention}

\begin{corrige}[Exercice 7 — Oxydation ménagée et tests d'identification]
\begin{enumerate}
  \item Le propan-1-ol est un alcool \textbf{primaire} : son oxydation ménagée donne d'abord l'aldéhyde, puis l'acide carboxylique.
  \begin{itemize}
    \item A : \chemfig{CH_3-CH_2-C(=[1]O)-[7]H}, le \cle{propanal} \ce{CH3CH2CHO} ;
    \item B : \chemfig{CH_3-CH_2-C(=[1]O)-[7]OH}, l'\cle{acide propanoïque} \ce{CH3CH2COOH}.
  \end{itemize}
  \item Deux tests pour identifier A (aldéhyde) :
  \begin{itemize}
    \item \textbf{DNPH} : formation d'un \emph{précipité jaune-orangé} (caractérise le groupe carbonyle) ;
    \item \textbf{Liqueur de Fehling} (bleu, chauffée) : formation d'un \emph{précipité rouge brique} d'oxyde de cuivre(I) \ce{Cu2O} (caractérise l'aldéhyde réducteur). On peut aussi citer le réactif de Tollens (miroir d'argent).
  \end{itemize}
  \item B est un acide carboxylique : on le met en évidence par son \textbf{caractère acide} sans test organique : il fait virer au rouge le papier pH (ou le BBT au jaune), on mesure un pH $< 7$.
  \item Demi-équations (milieu acide) :
  \begin{align*}
    \text{Reduction : } \ce{MnO4- + 8 H+ + 5 e- &-> Mn^2+ + 4 H2O} \quad (\times 2)\\
    \text{Oxydation : } \ce{CH3CH2CH2OH &-> CH3CH2CHO + 2 H+ + 2 e-} \quad (\times 5)
  \end{align*}
  Équation-bilan ionique :
  \[ \ce{2 MnO4- + 5 CH3CH2CH2OH + 6 H+ -> 2 Mn^2+ + 5 CH3CH2CHO + 8 H2O} \]
  \emph{Vérification des charges :} à gauche $2\times(-1)+6\times(+1)=+4$ ; à droite $2\times(+2)=+4$. Les atomes sont également conservés.
\end{enumerate}
\end{corrige}

\begin{attention}[Points de vigilance — Exercice 7]
L'oxydation de l'alcool primaire à l'aldéhyde ne consomme que \textbf{2 électrons} par molécule d'alcool (perte de 2 H), pas 4 (qui correspondrait à l'oxydation jusqu'à l'acide). La demi-équation d'oxydation s'écrit donc $\ce{RCH2OH -> RCHO + 2H+ + 2e-}$ (sans \ce{H2O} à gauche). Le bilan final est correct car les coefficients stœchiométriques (2 et 5) équilibrent les électrons (10 e-).
\end{attention}

\begin{corrige}[Exercice 8 — Estérification]
\begin{enumerate}
  \item Équation-bilan :
  \[ \ce{CH3COOH + CH3CH2CH2OH <=> CH3COOCH2CH2CH3 + H2O} \]
  \item L'ester formé est l'\cle{éthanoate de propyle} (nomenclature IUPAC).
  \item Deux caractéristiques (au choix) : réaction \textbf{lente} ; réaction \textbf{limitée} par la réaction inverse (hydrolyse) — c'est un équilibre ; réaction \textbf{athermique} ; réaction \textbf{catalysée} par les ions \ce{H3O+}.
  \item Pour augmenter le rendement en ester :
  \begin{itemize}
    \item utiliser l'un des réactifs en \textbf{excès} (par exemple l'alcool, moins coûteux) pour déplacer l'équilibre ;
    \item \textbf{éliminer l'ester} (ou l'eau) au fur et à mesure de sa formation par distillation.
  \end{itemize}
  \emph{Remarque :} l'acide sulfurique (catalyseur) accélère la réaction mais ne modifie pas le rendement à l'équilibre.
\end{enumerate}
\end{corrige}

\begin{attention}[Points de vigilance — Exercice 8]
Nomenclature de l'ester : \og éthanoate de \textit{alkyle} \fg{} (acide éthanoïque + alcool). Le catalyseur \ce{H2SO4} concentré apporte \ce{H3O+} et capte l'eau (agent déshydratant), mais ne déplace pas l'équilibre par lui-même. Seuls l'excès de réactif ou l'élimination d'un produit augmentent le rendement.
\end{attention}

\begin{corrige}[Exercice 9 — Type Bac : le sodium et l'éthanol]
\textbf{1.} Le gaz qui détone à la flamme est le \cle{dihydrogène} \ce{H2}. Le produit organique est l'\cle{éthanolate de sodium} \ce{C2H5ONa}. \hfill\textit{(1 pt)}

\textbf{2.} Équation-bilan :
\[ \ce{2 C2H5OH + 2 Na -> 2 C2H5ONa + H2 ^} \] \hfill\textit{(1 pt)}

\textbf{3.} Quantité de sodium : $n(\ce{Na}) = \dfrac{m}{M(\ce{Na})} = \dfrac{4,6}{23,0} = \SI{0,20}{mol}$.

Tableau d'avancement (éthanol en excès, sodium limitant) :
\begin{center}
\begin{tabularx}{\linewidth}{|l|X|X|X|X|}
\hline
\rowcolor{popblueL} État & $2\,\ce{C2H5OH}$ & $+2\,\ce{Na}$ & $\rightarrow 2\,\ce{C2H5ONa}$ & $+\,\ce{H2}$ \\
\hline
Initial & excès & $0,20$ & $0$ & $0$ \\
\hline
Final & excès & $0$ & $0,20$ & $0,10$ \\
\hline
\end{tabularx}
\end{center}
$x_{max} = \dfrac{0,20}{2} = \SI{0,10}{mol}$, donc $n(\ce{H2}) = x_{max} = \SI{0,10}{mol}$ et
\[ V(\ce{H2}) = n(\ce{H2}) \times V_m = 0,10 \times 24,0 = \SI{2,4}{L}. \] \hfill\textit{(2 pts)}

\textbf{4.} $n(\ce{C2H5ONa}) = 2\,x_{max} = \SI{0,20}{mol}$, d'où
\[ m(\ce{C2H5ONa}) = n \times M = 0,20 \times 68,0 = \SI{13,6}{g}. \] \hfill\textit{(1 pt)}

\textbf{5.} Précautions : manipuler le sodium avec des \textbf{pinces} (jamais à mains nues, il réagit avec l'humidité de la peau) ; utiliser un \textbf{petit morceau} de sodium et un excès d'alcool ; éloigner toute \textbf{flamme} du dégagement de \ce{H2} ; porter des lunettes de protection. \hfill\textit{(1 pt)}
\end{corrige}

\begin{attention}[Points de vigilance — Exercice 9]
Ne pas oublier les coefficients stœchiométriques 2 : $n(\ce{H2}) = n(\ce{Na})/2$, et non $n(\ce{Na})$. Erreur fréquente : écrire \ce{C2H5ONa} comme produit \og gazeux \fg{} — seul \ce{H2} est gazeux. Toujours conserver les unités dans les calculs. $V_m = \SI{24,0}{L.mol^{-1}}$ aux CNTP.
\end{attention}

\begin{corrige}[Exercice 10 — Type Bac : étude d'un alcool $C_4H_{10}O$]
\textbf{1.} Les quatre alcools isomères de formule \ce{C4H10O} :
\begin{center}
\begin{tabularx}{\linewidth}{|l|X|l|}
\hline
\rowcolor{popblueL} Formule semi-développée & Nom & Classe \\
\hline
\chemfig{CH_3-CH_2-CH_2-CH_2-OH} & butan-1-ol & primaire \\
\hline
\chemfig{CH_3-CH_2-CH(-[2]OH)-CH_3} & butan-2-ol & secondaire \\
\hline
\chemfig{CH_3-CH(-[2]CH_3)-CH_2-OH} & 2-méthylpropan-1-ol & primaire \\
\hline
\chemfig{H_3C-C(-[2]OH)(-[10]CH_3)-CH_3} & 2-méthylpropan-2-ol & tertiaire \\
\hline
\end{tabularx}
\end{center} \hfill\textit{(2 pts)}

\textbf{2.} La déshydratation intramoléculaire donne majoritairement l'alcène \textbf{le plus substitué} (règle de Zaïtsev). Seul le \textbf{butan-2-ol} peut conduire majoritairement au but-2-ène (le butan-1-ol donnerait le but-1-ène ; les deux autres donneraient le 2-méthylpropène). Donc \textbf{A est le butan-2-ol}, alcool secondaire. \hfill\textit{(1 pt)}

\textbf{3.} Déshydratation intramoléculaire de A :
\[ \ce{CH3-CH2-CH(OH)-CH3 ->[\text{Al2O3, vers 400 °C}] CH3-CH=CH-CH3 + H2O} \]
Conditions : passage des vapeurs d'alcool sur de l'alumine \ce{Al2O3} chauffée vers \SI{400}{\celsius} (ou \ce{H2SO4} concentré vers \SI{180}{\celsius}). \hfill\textit{(1 pt)}

\textbf{4.} Déshydratation intermoléculaire :
\[ \ce{2 CH3-CH2-CH(OH)-CH3 ->[\text{H2SO4 conc., vers 140 °C}] CH3CH2CH(CH3)-O-CH(CH3)CH2CH3 + H2O} \]
Le produit est l'\cle{éther di-sec-butylque} (2-(butan-2-yloxy)butane). On favorise cette voie par une température \textbf{modérée} (vers \SI{140}{\celsius}) et un \textbf{excès d'alcool}. \hfill\textit{(1,5 pt)}

\textbf{5.} D réagit avec la DNPH : c'est un composé \textbf{carbonylé}. Il ne réagit pas avec la liqueur de Fehling : ce n'est \textbf{pas un aldéhyde}, c'est donc une \textbf{cétone}. Or seule l'oxydation ménagée d'un alcool \textbf{secondaire} donne une cétone : cela confirme que A est le butan-2-ol. D est la \cle{butan-2-one} (butanone) : \chemfig{CH_3-CH_2-C(=[1]O)-[7]CH_3}, \ce{CH3-CH2-CO-CH3}. \hfill\textit{(1,5 pt)}

\textbf{6.} D'après l'équation de la question 3, $n(\text{alcool}) = n(\text{but-2-ène})$.
\[ n(\text{but-2-ène}) = \frac{V}{V_m} = \frac{5,6}{24,0} = \SI{0,233}{mol}. \quad (V_m = \SI{24,0}{L.mol^{-1}} \text{ aux CNTP}) \]
Masse molaire de A : $M(\ce{C4H10O}) = 4 \times 12,0 + 10 \times 1,0 + 16,0 = \SI{74,0}{g.mol^{-1}}$.
\[ m(\text{A}) = n \times M = 0,233 \times 74,0 = \SI{17,2}{g}. \] \hfill\textit{(2 pts)}
\end{corrige}

\begin{attention}[Points de vigilance — Exercice 10]
Bien distinguer les deux voies de déshydratation et leurs conditions de température. Le test DNPH est positif pour les aldéhydes \emph{et} les cétones : c'est le test de Fehling (ou Tollens) qui les départage. Vérifier la cohérence : un alcool secondaire ne peut jamais donner d'acide carboxylique par oxydation ménagée. \textbf{Convention :} le volume molaire $V_m = \SI{24,0}{L.mol^{-1}}$ aux conditions normales de température et de pression (CNTP, \SI{20}{\celsius}, \SI{1}{atm}) est la référence du programme ; l'utiliser systématiquement.
\end{attention}

\begin{corrige}[Exercice 11 — Type Bac : synthèse d'un arôme à la brasserie]
\textbf{1.} Équation de la fermentation alcoolique :
\[ \ce{C6H12O6 ->[enzymes\ (zymase)] 2 C2H5OH + 2 CO2 ^} \]
Conditions : présence de \textbf{levures} (enzymes, notamment la zymase), milieu \textbf{anaérobie} (absence de dioxygène), température voisine de \SI{30}{\celsius}. \hfill\textit{(1,5 pt)}

\textbf{2.} $n(\ce{C6H12O6}) = \dfrac{m}{M} = \dfrac{500 \times 10^{3}}{180} = \SI{2,78 \times 10^{3}}{mol}$.

D'après la stœchiométrie : $n(\ce{C2H5OH}) = 2\,n(\ce{C6H12O6}) = \SI{5,56 \times 10^{3}}{mol}$.
\[ m(\ce{C2H5OH}) = n \times M = 5,56 \times 10^{3} \times 46,0 \approx \SI{2,56 \times 10^{5}}{g} \approx \SI{256}{kg}. \]
Volume : $V = \dfrac{m}{\rho}$ avec $\rho = d \times \rho_{eau} = \SI{0,79}{kg.L^{-1}}$ :
\[ V(\ce{C2H5OH}) = \frac{256}{0,79} \approx \SI{324}{L}. \] \hfill\textit{(2,5 pts)}

\textbf{3.} $n(\ce{CO2}) = 2\,n(\ce{C6H12O6}) = \SI{5,56 \times 10^{3}}{mol}$, d'où
\[ V(\ce{CO2}) = n \times V_m = 5,56 \times 10^{3} \times 24,0 \approx \SI{1,33 \times 10^{5}}{L} \approx \SI{133}{m^3}. \] \hfill\textit{(1,5 pt)}

\textbf{4.} Demi-équations (milieu acide) :
\begin{align*}
  \text{Reduction : } \ce{MnO4- + 8 H+ + 5 e- &-> Mn^2+ + 4 H2O} \quad (\times 4)\\
  \text{Oxydation : } \ce{CH3CH2OH + H2O &-> CH3COOH + 4 H+ + 4 e-} \quad (\times 5)
\end{align*}
Équation-bilan ionique :
\[ \ce{4 MnO4- + 5 CH3CH2OH + 12 H+ -> 4 Mn^2+ + 5 CH3COOH + 11 H2O} \]
\emph{Vérification des charges :} gauche $4\times(-1)+12\times(+1)=+8$ ; droite $4\times(+2)=+8$. \hfill\textit{(2 pts)}

\textbf{5.} Équation de l'estérification :
\[ \ce{CH3COOH + CH3CH2OH <=> CH3COOCH2CH3 + H2O} \]
L'ester est l'\cle{éthanoate d'éthyle} (arôme fruité). Caractéristiques : réaction \textbf{lente} et \textbf{limitée} par l'hydrolyse inverse (équilibre) ; réaction \textbf{athermique}, \textbf{catalysée} par \ce{H3O+}. \hfill\textit{(1,5 pt)}

\textbf{6.} Tableau d'avancement (mélange équimolaire) :
\begin{center}
\begin{tabularx}{\linewidth}{|l|X|X|X|X|}
\hline
\rowcolor{popblueL} État & $\ce{CH3COOH}$ & $+\,\ce{C2H5OH}$ & $\rightleftharpoons$ ester & $+\,\ce{H2O}$ \\
\hline
Initial (mol) & $1,00$ & $1,00$ & $0$ & $0$ \\
\hline
Final (mol) & $1,00 - x$ & $1,00 - x$ & $x$ & $x$ \\
\hline
\end{tabularx}
\end{center}
Si la réaction était totale, $x_{max} = \SI{1,00}{mol}$. Avec un rendement de 67 \% :
\[ n(\text{ester}) = 0,67 \times 1,00 = \SI{0,67}{mol}. \]
Pour améliorer le rendement : mettre l'un des réactifs en \textbf{excès} (par exemple l'éthanol, produit sur place à la brasserie) ou \textbf{distiller l'ester} au fur et à mesure de sa formation pour déplacer l'équilibre. \hfill\textit{(2 pts)}
\end{corrige}

\begin{attention}[Points de vigilance — Exercice 11]
Dans la fermentation, ne pas oublier le facteur 2 : 1 mol de glucose donne 2 mol d'éthanol \emph{et} 2 mol de \ce{CO2}. Pour le volume d'éthanol (liquide), utiliser la masse volumique ($\rho = d \times \rho_{eau}$), pas le volume molaire (réservé aux gaz). Le catalyseur \ce{H2SO4} accélère l'estérification mais ne change pas le rendement à l'équilibre. Notation scientifique : privilégier $\times 10^{n}$ à la notation \og e \fg{}.
\end{attention}

\newpage

\coursec{Fiche bilan}

\courssub{Carte mentale de la leçon}

\begin{popfigure}
\begin{tikzpicture}[
  central/.style={circle, draw=popdark, fill=popblue, text=white, font=\bfseries\small, minimum size=2.8cm, align=center},
  branche/.style={rounded corners=4pt, draw=popdark, font=\scriptsize, align=center, inner sep=4pt, text width=3.0cm},
  lien/.style={thick, popdark}
]
\node[central] (C) {Propriétés\\chimiques\\des alcools};
\node[branche, fill=poporangeL] (N1) at (-5.5,2.8) {\textbf{Réaction avec le sodium}\\$\ce{2ROH + 2Na -> 2RO^-Na^+ + H2}$\\H du $-\mathrm{OH}$ mobile $\to$ alcoolate + $\ce{H2}$};
\node[branche, fill=popgreenL] (N2) at (0,3.8) {\textbf{Déshydratation}\\Intra (\SI{180}{\celsius}, $\ce{H2SO4}$) $\to$ alcène\\Inter (\SI{140}{\celsius}, $\ce{H2SO4}$) $\to$ éther-oxyde};
\node[branche, fill=poppinkL] (N3) at (5.5,2.8) {\textbf{Combustion}\\$\ce{C2H5OH + 3O2 -> 2CO2 + 3H2O}$\\Lampe sans flamme (Cu) : $\ce{CH3CHO}$};
\node[branche, fill=poppurpleL] (N4) at (-5.5,-2.8) {\textbf{Oxydation ménagée}\\Alcool I : aldéhyde $\to$ acide\\Alcool II : cétone\\Alcool III : pas de réaction};
\node[branche, fill=popgoldL] (N5) at (0,-3.8) {\textbf{Estérification}\\Acide + alcool $\rightleftharpoons$ ester + eau\\Lente, limitée, athermique};
\node[branche, fill=popblueL] (N6) at (5.5,-2.8) {\textbf{Préparation des alcools}\\Hydratation d'un alcène\\Fermentation alcoolique (anaérobie)};
\draw[lien] (C) -- (N1);
\draw[lien] (C) -- (N2);
\draw[lien] (C) -- (N3);
\draw[lien] (C) -- (N4);
\draw[lien] (C) -- (N5);
\draw[lien] (C) -- (N6);
\end{tikzpicture}
\legende{Carte mentale : les six grandes réactions des alcools au programme de Terminale TI.}
\end{popfigure}

\courssub{L'essentiel à retenir}

\begin{propriete}[Fiche de synthèse — Propriétés chimiques des alcools]
\textbf{Définitions clés}
\begin{itemize}
  \item \cle{Alcool} : composé organique portant un groupe caractéristique hydroxyle $-\mathrm{OH}$ lié à un carbone tétragonal (sp$^3$) ; formule générale $C_nH_{2n+1}OH$.
  \item \cle{Classe d'un alcool} : primaire (I), secondaire (II) ou tertiaire (III) selon que le carbone fonctionnel (porteur du $-\mathrm{OH}$) est lié à 1, 2 ou 3 groupes alkyles.
  \item \cle{Déshydratation} : élimination d'une molécule d'eau ; \emph{intramoléculaire} (un seul alcool) $\to$ alcène, \emph{intermoléculaire} (deux molécules d'alcool) $\to$ éther-oxyde.
  \item \cle{Oxydation ménagée} : oxydation qui conserve la chaîne carbonée (oxydants usuels : $\ce{KMnO4}$/$\ce{H+}$, $\ce{K2Cr2O7}$/$\ce{H+}$, $\ce{O2}$/Cu) ; s'oppose à la combustion (rupture totale de la chaîne en $\ce{CO2}$).
  \item \cle{Estérification} : réaction d'un acide carboxylique avec un alcool donnant un ester et de l'eau ; réaction \textbf{lente, limitée (équilibre), athermique} ($\Delta H \approx 0$).
\end{itemize}

\textbf{Équations-types à connaître par cœur (équilibrées atomes et charges)}
\begin{itemize}
  \item Sodium (mobilité de H du $-\mathrm{OH}$) : $\ce{2C2H5OH + 2Na -> 2C2H5O^-Na^+ + H2}$ (éthylate de sodium + dihydrogène).
  \item Déshydratation intramoléculaire : $\ce{CH3CH2OH ->[{\ce{conc. H2SO4}}, \SI{180}{\celsius}] CH2=CH2 + H2O}$ (éthène).
  \item Déshydratation intermoléculaire : $\ce{2CH3CH2OH ->[{\ce{conc. H2SO4}}, \SI{140}{\celsius}] CH3CH2-O-CH2CH3 + H2O}$ (oxyde de diéthyle).
  \item Combustion complète : $\ce{C2H5OH + 3O2 -> 2CO2 + 3H2O}$ (très exothermique).
  \item Oxydation ménagée alcool primaire (éthanol) : $\ce{CH3CH2OH ->[oxydation] CH3CHO ->[oxydation] CH3COOH}$ (éthanal puis acide éthanoïque).
  \item Oxydation ménagée alcool secondaire (propan-2-ol) : $\ce{CH3-CHOH-CH3 ->[oxydation] CH3-CO-CH3}$ (propanone).
  \item Alcool tertiaire : pas d'oxydation ménagée (absence d'H sur le carbone fonctionnel).
  \item Oxydation ménagée des aldéhydes : $\ce{R-CHO + [O] -> R-COOH}$ (ex. $\ce{CH3CHO ->[KMnO4/H+] CH3COOH}$).
  \item Estérification (ex. acide éthanoïque + éthanol) : $\ce{CH3COOH + C2H5OH <=> CH3COOC2H5 + H2O}$ (éthanoate d'éthyle).
  \item Hydratation d'un alcène (industriel) : $\ce{CH2=CH2 + H2O ->[{\ce{H3PO4}/\text{support}}, \SI{300}{\celsius}, \text{haute pression}] CH3CH2OH}$.
  \item Fermentation alcoolique (anaérobie) : $\ce{C6H12O6 ->[\text{levures/enzymes}] 2C2H5OH + 2CO2}$ (glucose $\to$ éthanol).
\end{itemize}

\textbf{Tableau récapitulatif de l'oxydation ménagée}
\begin{center}
\begin{tabularx}{\linewidth}{|l|X|X|X|}
\hline
\rowcolor{popblueL} \textbf{Classe} & \textbf{1\textsuperscript{er} produit} & \textbf{Si excès d'oxydant} & \textbf{Tests d'identification du produit} \\
\hline
Primaire (I) & Aldéhyde (R-CHO) & Acide carboxylique (R-COOH) & DNPH (précipité jaune) \textbf{puis} Schiff (rose) ou Fehling (rouge) \\
\hline
Secondaire (II) & Cétone (R-CO-R') & Rien (stabilité) & DNPH (précipité jaune) ; Schiff et Fehling \textbf{négatifs} \\
\hline
Tertiaire (III) & \multicolumn{2}{X|}{Pas d'oxydation ménagée (pas de H sur le carbone fonctionnel)} & --- \\
\hline
\end{tabularx}
\end{center}

\textbf{Méthodes express}
\begin{itemize}
  \item \cle{Identifier la classe} : compter les groupes alkyles (chaînes carbonées) directement liés au carbone porteur du $-\mathrm{OH}$.
  \item \cle{Prévoir le produit d'oxydation ménagée} : I $\to$ aldéhyde (contrôlé) $\to$ acide (excès) ; II $\to$ cétone ; III $\to$ aucune réaction.
  \item \cle{Distinguer aldéhyde et cétone} : tous deux donnent un précipité jaune-orangé avec la \cle{DNPH} (2,4-dinitrophénylhydrazine) ; \textbf{seul l'aldéhyde} rosit le réactif de \cle{Schiff} et réduit la liqueur de \cle{Fehling} (précipité rouge $\ce{Cu2O}$).
  \item \cle{Lampe sans flamme} : spirale de cuivre chauffée au rouge, plongée dans les vapeurs d'éthanol ; le cuivre reste incandescent (catalyseur) et l'odeur de pomme (éthanal) se dégage : $\ce{CH3CH2OH + 1/2 O2 ->[Cu] CH3CHO + H2O}$.
  \item \cle{Équilibrer une combustion} : 1) C $\to$ $\ce{CO2}$ ; 2) H $\to$ $\ce{H2O}$ ; 3) Ajuster $\ce{O2}$ (attention : l'alcool apporte déjà de l'O).
\end{itemize}
\end{propriete}

\begin{aretenir}[Checklist avant le Bac]
\begin{itemize}
  \item $\square$ Je sais écrire la formule générale d'un alcool $C_nH_{2n+1}OH$ et déterminer sa classe (I, II, III) à partir de sa formule semi-développée.
  \item $\square$ J'équilibre la réaction d'un alcool avec le sodium : $\ce{2ROH + 2Na -> 2RONa + H2}$ et je nomme l'alcoolate formé (ex. éthylate de sodium).
  \item $\square$ Je distingue clairement déshydratation intramoléculaire (\SI{180}{\celsius}, alcène) et intermoléculaire (\SI{140}{\celsius}, éther-oxyde) et j'en écris les équations.
  \item $\square$ J'équilibre la combustion complète de n'importe quel alcool $C_nH_{2n+1}OH$.
  \item $\square$ Je décris l'expérience de la lampe sans flamme (cuivre, vapeurs d'éthanol, odeur de pomme/éthanal) et j'écris l'équation catalytique.
  \item $\square$ Je prévois le(s) produit(s) d'oxydation ménagée selon la classe : I $\to$ aldéhyde/acide, II $\to$ cétone, III $\to$ pas de réaction ; je justifie l'inertie de l'alcool tertiaire.
  \item $\square$ Je sais oxyder un aldéhyde en acide carboxylique ($\ce{KMnO4}$/$\ce{H+}$ ou $\ce{K2Cr2O7}$/$\ce{H+}$) et j'écris l'équation.
  \item $\square$ Je maîtrise les tests de reconnaissance : DNPH (carbonyle), Schiff et Fehling (aldéhyde spécifique).
  \item $\square$ J'écris l'équation d'estérification avec une double flèche $\rightleftharpoons$ et je cite ses trois caractéristiques : \textbf{lente, limitée (équilibre), athermique}.
  \item $\square$ J'écris l'équation d'hydratation d'un alcène (catalyseur $\ce{H3PO4}$, conditions industrielles) et celle de la fermentation alcoolique du glucose (enzymes, anaérobie, $\ce{CO2}$).
  \item $\square$ Je relie la fermentation au contexte camerounais : \emph{vin de palme}, \emph{bili-bili}, \emph{jus de bissap} fermenté ; production artisanale d'éthanol, dégagement visible de $\ce{CO2}$ (bullissement), milieu anaérobie strict.
\end{itemize}
\end{aretenir}

\findecours

\end{document}
